\section{Contact calculations}\label{app:contacts}
Throughout this appendix, suppose that two nondegenerate independent laws
attain the maximum ratio $C=C_2\ge49/120$. The companion upper bound gives
$C<227/500$. Each law has at most three support points, and at least one
is a genuine triple: three distinct points with positive masses.
Initially normalize the total variance to one. Write $F$ for the selected
closed event probability, $D=F-\Phi(z)$, and $R$ for the sum of the two
third absolute moments. Strict Lyapunov convexity gives $R>1/\sqrt2$.
A representation of the threshold is a pair of support points whose sum
is $z$. The row-count notation for the event is defined in
Proposition~\ref{prop:finite-contact-exclusion}.

That proposition disposes of the zero-event, largest-point, and
negative-mixed-derivative cases. The remaining arguments are indexed
below. Transposition interchanges the two marginal laws. The certificate
identifiers refer to domains in
\path{certificates/finite/current-results.json}; the partitions and
formulas are checked as described in Section~\ref{sec:verification}.

{\small
\begin{longtable}{>{\raggedright\arraybackslash}p{.19\linewidth}>{\raggedright\arraybackslash}p{.25\linewidth}>{\raggedright\arraybackslash}p{.45\linewidth}}
\toprule
Marginals and event & Representations & Argument or certificate \\
\midrule
\endhead
Triple/pair $211$ & Unique &
Section~\ref{contact:one-triple-bernoulli-contacts} \\
Triple/pair $220$ & Unique middle/high &
Section~\ref{contact:one-triple-bernoulli-rectangle} \\
Triple/pair $221$ & Unique middle/high &
Section~\ref{contact:one-triple-bernoulli-upper-shape} \\
Triple/pair $210$ & Unique middle/low &
Section~\ref{contact:one-triple-bernoulli-middle-210} \\
Triple/pair $210$ & Unique low/high &
Section~\ref{contact:low-210-algebraic-reduction}; \texttt{low210} \\
Triple/pair $110$ & Unique &
Section~\ref{contact:lower-110-contact-reduction}; \texttt{lower110} \\
Triple/pair $211$ & Two &
Section~\ref{contact:one-triple-bernoulli-collision-211} \\
Triple/pair $221$ & Two, upper-gap collision &
Section~\ref{contact:collision-221-contact-reduction}; \texttt{collision221} \\
Triple/pair $210$ & Two, lower-gap collision &
Sections~\ref{contact:collision-210-contact-reduction}--\ref{contact:collision-210-raw-law-replay}; \texttt{collision210} \\
Two triples $220$ & Unique &
Sections~\ref{contact:two-triple-rectangle-exclusion}--\ref{contact:two-triple-rectangle-exclusion-continued} \\
Two triples $221,320$ & Unique &
Section~\ref{contact:two-triple-asymmetric-exclusion} \\
Two triples $321$ & Unique middle/middle &
Section~\ref{contact:two-triple-triangle-exclusion} \\
Two triples $322,331$ & Two &
Section~\ref{contact:two-triple-collision-322-exclusion} \\
Two triples $321$ & Three &
Sections~\ref{contact:three-owner-321-contact-reduction}--\ref{contact:three-owner-321-scalar-closure}; \texttt{three\_owner321} \\
Two triples $321$ & Two adjacent &
Sections~\ref{contact:two-owner-321-contact-reduction}--\ref{contact:adjacent-owner-321-scalar-reduction}; \texttt{adjacent\_owner321} and its completion \\
Two triples $321$ & Two endpoints &
Sections~\ref{contact:two-owner-321-contact-reduction}, \ref{contact:endpoint-owner-321-scalar-reduction}; \texttt{endpoint\_owner321} \\
Two triples $221,320$ & Two &
Sections~\ref{contact:two-triple-collision-221-reduction}--\ref{contact:positive-parameter}; \texttt{tt221\_negative}, \texttt{tt221\_positive} \\
\bottomrule
\end{longtable}
}

\paragraph{Changes of scale}
The discrepancy $D$ and ratio $C$ are unchanged by a common positive
rescaling. In any raw scale, let $V$ be the total variance, $R_0$ the
sum of the centered third absolute moments, and $o$ the centered
threshold. Then
\[
 z=\frac{o}{\sqrt V},\qquad R_{\rm normalized}=\frac{R_0}{V^{3/2}},
 \qquad C=\frac{DV^{3/2}}{R_0}.
\]
Two choices of raw scale are used:
\begin{center}\small
\begin{tabular}{>{\raggedright\arraybackslash}p{.20\linewidth}>{\raggedright\arraybackslash}p{.69\linewidth}}
\toprule
Scale & Definition and conversion\\
\midrule
Contact scale &
Divide the variance-one supports by
$K=R_{\rm normalized}+z\phi(z)/(3C)$.
Then $V=K^{-2}$ and $\tau=CK^3=C/V^{3/2}$.
After division by $\tau$, the cubic majorant has quadratic coefficient
$-3/2$.\\[4pt]
Unit-span scale &
Translate the supports to start at zero and divide by the chosen
outer support span. Center these raw laws when computing $V,R_0,o$.
Setting a support span to one does not set the contact parameter to one.\\
\bottomrule
\end{tabular}
\end{center}
Each family states its scale. In particular the three-representation
calculations use the unit-span scale, while the collided $221$ systems
use the contact scale.

\subsection{Variance derivatives at a triple}\label{contact:three-point-variance-curvature}
The total variance is one in this calculation.

Let $x_1<x_2<x_3$ be the support of a genuine triple. Its second divided
difference, normalized by $[x_1,x_2,x_3]x^2=1$, satisfies
\begin{equation}\label{eq:contact-curvature}
 [x_1,x_2,x_3]|x|^3\le\frac{15D+\phi'''(z)}{12C}.
\end{equation}
Indeed, take a nonzero signed measure $\eta$ on these points with zero
mass and first moment. Put $v=\int x^2\,d\eta$ and
$r=\int|x|^3\,d\eta$. The Vandermonde determinant gives $v\ne0$, and
$r/v=[x_1,x_2,x_3]|x|^3$. Both signs of $t$ are allowed in $P+t\eta$.
The total variance becomes $V(t)=1+tv$, the sum of third absolute moments becomes $R+tr$,
and the event probability becomes $g(t)=F+tg'$. Maximality gives
\[
 J(t)=V(t)^{3/2}\{g(t)-\Phi(z/\sqrt{V(t)})\}-C(R+tr)\le0,
 \qquad J(0)=0.
\]
With $A=(3D+z\phi(z))/2$, differentiation gives
\[
 g'+Av-Cr=0,\qquad
 J''(0)=v^2\left\{3C\frac rv-\frac{15D+\phi'''(z)}4\right\}\le0.
\]
This proves \eqref{eq:contact-curvature}.

For any three real knots of diameter $h$,
\begin{equation}\label{eq:contact-dd-diameter}
 [x_1,x_2,x_3]|x|^3\ge h/2.
\end{equation}
On one side of zero the divided difference is the sum of the absolute
values. Otherwise scale $h$ to one and reflect to write the knots as
$-a,c,1-a$, where $0<a<1$ and $0\le c\le1-a$. The divided difference
$q$ obeys
\[
 q=1+c-2a+\frac{2a^3}{a+c},\qquad
 (a+c)(q-1/2)=2a(a-1/2)^2+c(c+1/2-a).
\]
The right side is nonnegative for $a\le1/2$. For $a>1/2$, put
$d=a-1/2$ and rewrite it as $(c-d/2)^2+(2a-1/4)d^2\ge0$.
Equality requires the symmetric knots $(-h/2,0,h/2)$.
\subsection{A joint curvature and diameter bound}\label{contact:joint-curvature-diameter}
For every centered variance-one law,
\begin{equation}\label{eq:joint-normal-bound}
 15\{F(z)-\Phi(z)\}+\phi'''(z)<79/10.
\end{equation}
For $z\ge0$ use $F(z)\le1$. On $[0,\sqrt3]$ the derivative of
$15(1-\Phi(z))+\phi'''(z)$ is
$(z^4-6z^2-12)\phi(z)<0$; its value at zero is $15/2$. Beyond $\sqrt3$
the third derivative is nonpositive, giving the same upper bound.

For $z=-x<0$, Cantelli gives
$D\le d(x)=(1+x^2)^{-1}-(1-\Phi(x))$.
The ratio $2x/\{(1+x^2)^2\phi(x)\}$ increases, since its logarithmic
derivative is $(x^2-1)^2/\{x(1+x^2)\}$. At $x=1/2$ the ratio exceeds
one; moreover $1-\Phi(1/2)>3/10$. Hence $d$ decreases on $[1/2,\infty)$
and $d(1/2)<1/2$. On $[1/2,\sqrt3]$, $\phi'''(-x)\le0$, so the
quantity in \eqref{eq:joint-normal-bound} is less than $15/2$.
For $x\ge\sqrt3$, $d(x)<1/4$ and Fourier inversion gives
$\|\phi'''\|_\infty\le2/\pi<2/3$, yielding $53/12$.

It remains to consider $0\le x\le1/2$. The bounds
$1-\Phi(x)\ge1/2-\phi(0)x$,
$\phi(x)\ge\phi(0)(1-x^2/2)$, and $\phi(0)<2/5$ give the upper bound
\[
 U(x)=\frac{15}{1+x^2}-\frac{15}2+\frac{24}5x+x^3-\frac{x^5}5.
\]
Here
\[
 10(1+x^2)\{79/10-U(x)\}
 =4-48x+154x^2-58x^3-8x^5+2x^7=:p(x).
\]
On each interval $[i/64,(i+1)/64]$, $0\le i<32$, expand
$p((i+t)/64)$ in the degree-seven Bernstein basis. Direct rational
conversion gives a positive coefficient minimum
$11964025693/3848290697216$. For completeness this assertion means the
following finite rational identity and comparison: if
$p(x)=\sum_{j=0}^7a_jx^j$, compute
\[
 c_{i,k}=\sum_{j=k}^7 a_j\binom jk\frac{i^{j-k}}{64^j},\qquad
 b_{i,l}=\sum_{k=0}^l c_{i,k}\frac{\binom lk}{\binom7k}
 \quad(0\le i<32,\ 0\le l\le7).
\]
The minimum of these 256 fractions is the stated number.
The nonnegative Bernstein basis functions sum to one, proving positivity.

Combining \eqref{eq:contact-curvature}--\eqref{eq:joint-normal-bound}
with $C\ge49/120$ shows that every normalized triple diameter is less
than $158/49$. In raw coordinates of total variance $V$ and triple
diameter $A_i$ this is
\begin{equation}\label{eq:raw-triple-diameter}
 24964V>2401A_i^2.
\end{equation}
If the two raw spans equal one, it follows that
\[
 V>2401/24964,
 \qquad R\ge V^{3/2}/\sqrt2>\frac{49^3}{158^3\sqrt2}>1/48.
\]
The last strict comparison follows on squaring from
$49^6\,48^2>2\,158^6$.
\subsection{Threshold and density bounds}\label{contact:fixed-c2-threshold-collar}


For every centered variance-one law with positive closed-CDF discrepancy

\[
 D=F(z)-\Phi(z)\ge d_0:=49/(120\sqrt2),
\]

one necessarily has

\[
 \boxed{-5/4<z<3/5,\qquad \phi(z)>9/50.}
\tag{C.3.1}\label{contact:fixed-c2-threshold-collar-1}
\]

This scalar lemma does not assume stationarity or a finite support. Its
application to a violating fixed-$C_2$ maximum uses only the already proved
$R_{\mathrm{normalized}}\ge 1/\sqrt{2}$ and $C\ge49/120$. It does not supply a uniform
threshold bound for arbitrary numbers of summands, whose normalized
additive third absolute moment sum can be arbitrarily small.

\paragraph{The negative threshold}

For $z=-s<0$, Cantelli's inequality, including the closed event, gives

\[
 D\le g(s):={1\over1+s^2}-\Phi(-s).
\]

The ratio

\[
 H(s)={2s\over(1+s^2)^2\phi(s)}
\]

satisfies

\[
 {d\over ds}\log H(s)={(s^2-1)^2\over s(1+s^2)}\ge0.
\]

Also $2s/(1+s^2)^2=640/1681>1/4$ at $s=5/4$, while
$\phi(5/4)<\phi(1)<1/4$. The last inequality follows from $2 \pi e>16$, using
$\pi>3$ and $e>8/3$. Thus $H(s)>1$ and $g'(s)<0$ for every $s\ge 5/4$.

Taylor's formula with the signed remainder gives, for $u\ge 0$,
$\exp(-u)\le 1-u+u^2/2-u^3/6+u^4/24$. At $s=5/4$, let

\[
 I=s-s^3/6+s^5/40-s^7/336+s^9/3456.
\]

Integrating the polynomial and using $\phi(0)<2/5$ gives

\[
 \Phi(-s)>1/2-(2/5)I,
\]
\[
 g(5/4)<16/41-1/2+(2/5)I
 ={37147949759\over130006646784}=:B_-.
\]

Both $B_-$ and $d_0$ are positive, and the exact comparison is

\[
 d_0^2-B_-^2
 ={727351123347544172903\over422543205200493438566400}>0.
\]

Consequently $D\ge d_0$ excludes every $z\le -5/4$.

\paragraph{The positive threshold}

The classical $\pi<22/7$ implies $\phi(0)>79/200$, because

\[
 7/44-(79/200)^2=1349/440000>0.
\]

For $x\ge 0$, $\exp(-x^2/2)\ge 1-x^2/2$. Hence at $s=3/5$,

\[
 1-\Phi(s)<1/2-(79/200)(s-s^3/6)=13861/50000=:B_+.
\]

The exact difference $d_0^2-B_+^2=146635361/22500000000$ is positive.
Since $F(z)\le 1$, all $z\ge 3/5$ have $D\le 1-\Phi(3/5)<d_0$.

\paragraph{A density bound and the raw contact slope}

By \eqref{contact:fixed-c2-threshold-collar-1}, $|z|<5/4$, so $\phi(z)>\phi(5/4)$. Put $x=25/32$. The degree-five
Taylor polynomial is a lower bound for $\exp(-x)$ for all $x\ge 0$, giving

\[
 \phi(5/4)>{79\over200}\sum_{j=0}^{5}{(-25/32)^j\over j!}
 ={29108943797\over161061273600}>9/50.
\]

The final exact slack is $117914549/161061273600>0$.

The raw contact normalizations used in the contact systems have
$G>1/2$, $\sqrt{V}<2$, and $\sigma \sqrt{V}=G \phi(z)$. Therefore

\[
 \boxed{\sigma>9/200.}
\tag{C.3.2}\label{contact:fixed-c2-threshold-collar-2}
\]
\subsection{The raw contact scale}\label{contact:fixed-c2-contact-scale-collar}
Here supports are divided by $K$, and $V,R,o$ denote their raw total variance, sum of third absolute moments, and centered threshold.



The collided triple/pair $210$ and two-triple $221$ contact charts satisfy

\[
 C=V^{3/2}/G,\quad D=R/G,\quad z=o/\sqrt V,
 \quad\sigma\sqrt V=G\phi(z),
 \quad3(V-R)=\sigma o.
\]

Consequently

\[
 \boxed{C/\sqrt V=D+z\phi(z)/3.}
\tag{C.4.1}\label{contact:fixed-c2-contact-scale-collar-1}
\]

The threshold result
Section~\ref{contact:fixed-c2-threshold-collar} gives
$-5/4<z<3/5$. On $0\le z<3/5$, $z \phi(z)$ increases, and
$\phi(3/5)<1/3$ as verified below. For negative $z$, the product is negative.
Thus in all these contact charts $z \phi(z)<1/5$. Combining \eqref{contact:fixed-c2-contact-scale-collar-1} with
$D<11/20$ and $C\ge49/120$ proves

\[
 {1\over\sqrt V}<{{11\over20}+{1\over15}\over49/120}
 ={74\over49},\qquad
 \boxed{V>2401/5476.}
\tag{C.4.2}\label{contact:fixed-c2-contact-scale-collar-2}
\]

If the actual representation satisfies $o\le 0$, then $z\le 0$ and the stronger bound is

\[
 {1\over\sqrt V}<{11/20\over49/120}={66\over49},\qquad
 \boxed{V>2401/4356.}
\tag{C.4.3}\label{contact:fixed-c2-contact-scale-collar-3}
\]

The entire negative $210$ chart has $o=c-ph<0$, so \eqref{contact:fixed-c2-contact-scale-collar-3} holds on that whole
candidate class. In another chart it can be used when the enclosing interval proves $o\le 0$ for every surviving candidate. It cannot
be applied to a box that still permits both signs without an additional
valid case split.

These scale formulas are specific to the chart normalization fixing the
quadratic coefficient of the cubic contact potential. They are not asserted
for the unit-span three-representation $321$ chart, which uses a different raw scale.

For the elementary density endpoint, the classical Archimedean lower bound
$\pi>223/71$ gives $\phi(0)<399/1000$, because

\[
 (399/1000)^2-71/446=1823/223000000>0.
\]

The degree-four Taylor upper bound for $\exp(-x)$ at $x=9/50$ now gives

\[
 \phi(3/5)<{399\over1000}\sum_{j=0}^4{(-9/50)^j\over j!}
 ={16663671213\over50000000000}<1/3.
\]
\subsection{One triple: contact identities and threshold shapes}\label{contact:one-triple-bernoulli-contacts}
In this subsection the support coordinates have total variance one; factors of $K$ are retained explicitly.



\paragraph{Derivative identities for the Bernoulli law}

Write the Bernoulli support as $-u,v$, with $u,v>0$, and set

\[
 h_B=u+v,\qquad K_B=\frac{u^2+v^2}{u+v}.
\]

Its variance is $uv$, its third absolute moment is $uv K_B$, and
its signed second moment is $M_B=uv(v-u)/(u+v)$. Its dual derivative is

\[
 f_B'(x)=3C\{x|x|-Kx-M_B\}-\phi(z).
\tag{C.5.1}\label{contact:one-triple-bernoulli-contacts-1}
\]

The nonselected point is a smooth minimum. Applying \eqref{contact:one-triple-bernoulli-contacts-1} there gives

\[
 \begin{array}{c|c|c}
 \text{Bernoulli representation}&\text{exact identity}&\text{consequence}\\ \hline
 -u&\phi(z)=3Cv(K_B-K)&K_B>K,\quad 2v>K\\
 v&\phi(z)=3Cu(K-K_B)&K_B<K,\quad 2u>K.
 \end{array}
\tag{C.5.2}\label{contact:one-triple-bernoulli-contacts-2}
\]

The second derivative gives the last inequalities. Equality would be
a stationary inflection rather than a local minimum. Also

\[
 K_B\geq\frac{h_B}{2},\qquad K_B<h_B.
\tag{C.5.3}\label{contact:one-triple-bernoulli-contacts-3}
\]

\paragraph{Exhaustive grid and three elementary eliminations}

List the number of included Bernoulli support points in each successive
triple row. The ten possible downsets are

\[
 000,100,110,111,200,210,211,220,221,222.
\]

A unique selected representation must be a maximal included cell. A genuine
triple cannot have its largest support point as the selected representation, by the two smooth-minimum
argument in the proof of Proposition~\ref{prop:finite-contact-exclusion}. This excludes $111$ and $222$; $000$ has
negative signed discrepancy.

Shapes $100$ and $200$ have a selected lowest triple point, while both
nonselected triple points have deletion-CDF value zero. They are thus
two equal zero minima. The cubic identity gives them as $-K,K$, with

\[
 f_T(x)=C|x|^3-\frac32 CKx^2+\frac12 CK^3.
\]

The selected triple point is $-a$, with $a>K$, and let its mass be $l$.
All event probability comes from that triple point, so

\[
 F=l f_T(-a)=Cl(a-K)^2(a+K/2)<Cl a^3\leq CR.
\]

This contradicts $F-\Phi(z)=CR>0$. Shape $200$ need not be the full
sum's lower endpoint, so this slightly extends the endpoint argument.

For shape $211$, the only possible representation is triple-low/Bernoulli-high.
The two nonselected triple points have the same deletion level, hence
are the symmetric smooth minima $-K,K$. Write the selected triple point
as $-a$, with $a>K$. The included triple-high/Bernoulli-low cell is
strictly below the representation, giving

\[
 K-u<-a+v,\qquad h_B>a+K>2K.
\]

But \eqref{contact:one-triple-bernoulli-contacts-2}--\eqref{contact:one-triple-bernoulli-contacts-3}, for a high Bernoulli representation, require
$h_B/2\le K_B<K$. This is a contradiction. Thus $211$ is also excluded.

Exactly five representation-specific configurations remain: $110$ with the
middle/low representation, the two representations of $210$, and $220$ or $221$ with
the middle/high representation.

\paragraph{Exact data for the five remaining configurations}

When the selected representation uses the middle triple point, write the physical triple support as $-a,c,b$.
Its two nonselected minima satisfy

\[
 \frac{K}{2}<a<K<b,\qquad a^2+b^2=K(a+b).
\tag{C.5.4}\label{contact:one-triple-bernoulli-contacts-4}
\]

Let its masses be $l,m,r$, put $q_B=v/(u+v)$, and let
$u_T=a+c$, $v_T=b-c$. The actual deletion levels and strict geometry are:

\begin{center}\small
\begin{tabular}{llll}\toprule
Shape and representation & Triple levels & Bernoulli levels & Gap condition\\\midrule
$110$, middle/low & $(q_B,q_B,0)$ & $(l+m,0)$ & $h_B>u_T$\\
$210$, middle/low & $(1,q_B,0)$ & $(l+m,l)$ & $h_B<u_T$\\
$220$, middle/high & $(1,1,0)$ & $(l+m,l+m)$ & $h_B<v_T$\\
$221$, middle/high & $(1,1,q_B)$ & $(1,l+m)$ & $h_B>v_T$\\
\bottomrule\end{tabular}\end{center}

In the middle/low $210$ case, \eqref{contact:one-triple-bernoulli-contacts-2}--\eqref{contact:one-triple-bernoulli-contacts-3} give

\[
 a+c>h_B>K_B>K,\qquad c>K-a>0.
\tag{C.5.5}\label{contact:one-triple-bernoulli-contacts-5}
\]

For the low/high representation of $210$, instead write the triple support as
$-t,-a,b$, where $t>a$ and its two nonselected minima $-a,b$ satisfy
\eqref{contact:one-triple-bernoulli-contacts-4}. Its deletion levels are $(1,q_B,0)$, the Bernoulli deletion levels
are $(l+m,l)$, and its exact strict geometry is

\[
 t-a<h_B<t+b,\qquad K_B<K.
\tag{C.5.6}\label{contact:one-triple-bernoulli-contacts-6}
\]

Every table entry follows directly from the indicated threshold cell;
none replaces the actual convolution probabilities by free parameters.

\subsection{Two triples: symmetry of a unique-representation rectangle}\label{contact:two-triple-rectangle-exclusion}


\paragraph{Common shape coordinates}

Use the full-law cubic dual from
Lemma~\ref{lem:full-law-contacts}, and normalize the total variance to one. Put

\[
 K=\frac{2A}{3C}=R+\frac{z\phi(z)}{3C}>0.
\]

For coordinate $i$, write its support as $K(-a_i,c_i,b_i)$, where
$-a_i,b_i$ are the two nonselected smooth minima and $c_i$ is the selected
middle point. These are dimensionless support coordinates. The minimum
equations and the strictly decreasing contact levels give

\[
 \frac12<a_i<1<b_i,
 \qquad a_i^2+b_i^2=a_i+b_i,
 \qquad b_i=\frac{1+\sqrt{1+4a_i-4a_i^2}}2.
\tag{C.6.1}\label{contact:two-triple-rectangle-exclusion-1}
\]

Set

\[
 h_i=a_i+b_i,\quad u_i=a_i+c_i,\quad v_i=b_i-c_i,
\]

\[
 D(a)=\frac12(b-a)(a^2+4ab+b^2),\qquad
 H(a,c)=3\{a(1-a)+c-c|c|\}.
\tag{C.6.2}\label{contact:two-triple-rectangle-exclusion-2}
\]

Thus $C K^3 D(a_i)$ is the difference of the two minimum contact
levels, and $C K^2 H(a_i,c_i)$ is the negative slope at the representation.
Let $m_i$ be the middle mass and $q_i$ the combined mass of the first
two support points. The unique threshold representation and centering imply

\[
 \sigma=\frac{\phi(z)}{C K^2}>0,\qquad
 m_i=\frac\sigma{H(a_i,c_i)},\qquad
 q_i=\frac{b_i}{h_i}+\frac{\sigma u_i}{h_iH(a_i,c_i)}.
\tag{C.6.3}\label{contact:two-triple-rectangle-exclusion-3}
\]

For shape $220$, the dual equals $q_j$ at the first and middle support
points of coordinate $i$, and equals zero at its last support point.
Consequently

\[
 f_i(c_iK)=f_i(-a_iK),\qquad
 q_j=C K^3D(a_i).
\tag{C.6.4}\label{contact:two-triple-rectangle-exclusion-4}
\]

In particular,

\[
 q_1D(a_1)=q_2D(a_2).
\tag{C.6.5}\label{contact:two-triple-rectangle-exclusion-5}
\]

\paragraph{The contact equations force identical laws}

For each $a$ in $(1/2,1)$, equation $f(c)=f(-a)$ selects one middle
point on the descending part of the cubic between its local maximum and
its positive minimum. Its normalized location is a strictly increasing
function of $a$.

For $1/2<a\le 3/4$, factorization on the negative half-line gives

\[
 c=2a-3/2,\qquad H=(a+c)^2.
\tag{C.6.6}\label{contact:two-triple-rectangle-exclusion-6}
\]

For $3/4\le a<1$, the middle point is nonnegative and satisfies $c<a$.
Indeed, the linear coefficient of the dual is negative, so
$f(Ka)-f(-Ka)=-6CK^3a^2(1-a)<0$, whereas $f$ decreases
strictly along this positive part of the descending branch. Set
$k=c/a$, so $0\le k<1$. The equal-height equation is exactly

\[
 a=\frac{3(1+k)^2}{2(k^3+3k+2)},\qquad c=ka.
\tag{C.6.7}\label{contact:two-triple-rectangle-exclusion-7}
\]

Its derivative is

\[
 a'(k)=\frac{3(1+k)(1-k)(k^2+4k+1)}{2(k^3+3k+2)^2}>0
 \quad(0\leq k<1).
\tag{C.6.8}\label{contact:two-triple-rectangle-exclusion-8}
\]

We claim that, for every fixed $\sigma>0$, the expression in \eqref{contact:two-triple-rectangle-exclusion-3} obeys

\[
 a\longmapsto q(a)D(a)
 =D(a)\frac b{a+b}
  +\sigma\frac{D(a)(a+c)}{(a+b)H(a,c)}
 \quad\text{strictly decreasing on }(1/2,1).
\tag{C.6.9}\label{contact:two-triple-rectangle-exclusion-9}
\]

First, direct differentiation using \eqref{contact:two-triple-rectangle-exclusion-1} gives
$D'(a)=-3(2a-1)(a+b)<0$. Also $b$ decreases with $a$, so
$b/(a+b)$ strictly decreases. Hence the first term in \eqref{contact:two-triple-rectangle-exclusion-9} strictly
decreases. For the second term, put

\[
 T(a)=\frac{D(a)(a+c)}{(a+b)H(a,c)}.
\]

On the nonpositive branch, \eqref{contact:two-triple-rectangle-exclusion-6} gives $T=D/[(a+b)(a+c)]$.
Here $D$ decreases, while $a+b$ and $a+c$ increase, so $T$ decreases.

On the positive branch, the cubic factors at its positive minimum:

\[
 D(a)=(b-c)^2(c+2b-3/2),\qquad
 H(a,c)=3(b-c)(b+c-1).
\]

Thus the exact function to check is

\[
 T(k)=\frac{(b-c)(a+c)(c+2b-3/2)}
                  {3(a+b)(b+c-1)},
\tag{C.6.10}\label{contact:two-triple-rectangle-exclusion-10}
\]

with \eqref{contact:two-triple-rectangle-exclusion-1}, \eqref{contact:two-triple-rectangle-exclusion-7}. The denominator stays positive through $k=1$, so \eqref{contact:two-triple-rectangle-exclusion-10}
and its derivative extend continuously to the closed interval $[0,1]$.
The fixed certificate \texttt{rectangle}, checked by the rules in
Section~\ref{sec:verification}, proves

\[
 \boxed{\quad T'(k)<-1/4\quad\text{for every }0\leq k\leq1.\quad}
\tag{C.6.11}\label{contact:two-triple-rectangle-exclusion-11}
\]

The two formulas for $T$ agree at $a=3/4$, so \eqref{contact:two-triple-rectangle-exclusion-9} follows. Equation
\eqref{contact:two-triple-rectangle-exclusion-5} now gives $a_1=a_2$. Equations \eqref{contact:two-triple-rectangle-exclusion-1}, \eqref{contact:two-triple-rectangle-exclusion-6}--\eqref{contact:two-triple-rectangle-exclusion-8}, and \eqref{contact:two-triple-rectangle-exclusion-3} then give
identical supports, middle masses, and outer masses for the two
coordinates. In particular $c_1=c_2=c$.
\subsection{Two triples: nonpositive rectangle representations}\label{contact:two-triple-nonpositive-rectangle}
Here $-a,c,b$ are the support points in the variance-one scale, with $K$ retained explicitly.



\paragraph{The equal-height contacts}

Write the discrepancy as $D=C R$, the selected threshold as $z$, and set

\[
 K=\frac{2A}{3C}=R+\frac{z\phi(z)}{3C}.
\]

For each coordinate, write its two nonselected smooth minima as $-a,b$
and its selected middle atom as $c$, of mass $m$. The cubic dual has
derivative

\[
 f'(u)=3C u|u|-3CKu-L.
\]

The two minimum equations imply

\[
 a^2+b^2=K(a+b),\qquad L=3Ca(K-a).
\tag{C.7.1}\label{contact:two-triple-nonpositive-rectangle-1}
\]

The decreasing deletion CDF has strictly larger value at $-a$ than at
$b$ in shape $220$. The minimum-height difference is

\[
 f(-a)-f(b)=\frac{C}{2}(b-a)(a^2+4ab+b^2)>0.
\]

Thus $b>a$. The strict local-minimum conditions and \eqref{contact:two-triple-nonpositive-rectangle-1} consequently give

\[
 K/2<a<K<b,\qquad
 b=\frac{K+\sqrt{K^2+4Ka-4a^2}}2.
\tag{C.7.2}\label{contact:two-triple-nonpositive-rectangle-2}
\]

In this rectangle, the deletion CDF has the same value at $-a$ and $c$,
so $f(c)=f(-a)$. Suppose $c\le 0$. On the negative half-line, direct
factorization gives

\[
 f(c)-f(-a)=C(c+a)^2(2a-3K/2-c).
\]

Since the support points are distinct, this yields

\[
 c=2a-3K/2,\qquad K/2<a\leq3K/4,
 \qquad -f'(c)=C(a+c)^2.
\tag{C.7.3}\label{contact:two-triple-nonpositive-rectangle-3}
\]

The unique-representation support derivative is $m f'(c)=-\phi(z)$. Therefore

\[
 m=\frac{\phi(z)}{C(a+c)^2}.
\tag{C.7.4}\label{contact:two-triple-nonpositive-rectangle-4}
\]

The use of \eqref{contact:two-triple-nonpositive-rectangle-4} is exactly where uniqueness of the selected representation
is needed. The proof does not extend \eqref{contact:two-triple-nonpositive-rectangle-4} to a collision.

\paragraph{A variance loss forced by the selected mass}

For a centered law on $-a,c,b$ with middle mass $m$, its variance is

\[
 v=ab-m(a+c)(b-c)
   =ab-\frac{\phi(z)}C\frac{b-c}{a+c}.
\tag{C.7.5}\label{contact:two-triple-nonpositive-rectangle-5}
\]

This follows by solving the two remaining masses, or by averaging
the identity $x^2+(a-b)x-ab=0$ at the two outer support points.

On $K/2<a\le 3K/4$, the function $ab$ in \eqref{contact:two-triple-nonpositive-rectangle-2} increases with $a$. Its
maximum is

\[
 \frac{3(2+\sqrt7)}{16}K^2<\frac78K^2.
\]

Indeed, after division by $K$, the derivative of $x b(x)$ is positive
for $1/2\le x\le 3/4$, while $\sqrt{7}<8/3$ proves the displayed bound.
Also $b>9K/8$, since the minimum of $b$ in this range is
$(2+\sqrt{7})K/4$ and $\sqrt{7}>5/2$. Equation \eqref{contact:two-triple-nonpositive-rectangle-3} gives
$a+c\le 3K/4$, so

\[
 \frac{b-c}{a+c}>\frac32.
\]

Consequently every such coordinate satisfies

\[
 \boxed{\quad v<\frac78K^2-\frac32\frac{\phi(z)}C.\quad}
\tag{C.7.6}\label{contact:two-triple-nonpositive-rectangle-6}
\]

If both selected middle points are nonpositive, their sum is $z=-t$
with $0\le t<K$, by \eqref{contact:two-triple-nonpositive-rectangle-3}. Adding \eqref{contact:two-triple-nonpositive-rectangle-6} over the two coordinates gives

\[
 1<\frac74K^2-3\frac{\phi(t)}C.
\tag{C.7.7}\label{contact:two-triple-nonpositive-rectangle-7}
\]

\paragraph{The right side is always less than one}

The proved global upper bound gives $C<1/2$, and $C\ge49/120>2/5$.
The public variance-only discrepancy bound, with
$\kappa=271/500$, gives $CR<\kappa$. Hence

\[
 K<\frac{\kappa-t\phi(t)/3}{C}.
\tag{C.7.8}\label{contact:two-triple-nonpositive-rectangle-8}
\]

First, \eqref{contact:two-triple-nonpositive-rectangle-8} and $t<K$ imply $t<\kappa/(2/5)<7/5$ and

\[
 t\{2/5+\phi(t)/3\}<\kappa.
\]

The left side increases on $[0,7/5]$: its derivative is
$2/5+(1-t^2)\phi(t)/3>2/5-(24/25)(2/5)/3>0$.
At $t=6/5$ it exceeds

\[
 \frac65\left(\frac25+\frac{19}{300}\right)
 =\frac{139}{250}>\frac{271}{500},
\]

using $\phi(6/5)>19/100$. Thus $t<6/5$.

Put $B(t)=\kappa-t \phi(t)/3$. Since
$\sup_{t\ge 0}t \phi(t)=\phi(1)<1/4$, we have $B(t)>9/20$.
For $2/5\le C\le 1/2$, the derivative in $C$ of

\[
 \frac{7B(t)^2}{4C^2}-\frac{3\phi(t)}C
\]

has the sign of $3C \phi(t)-(7/2)B(t)^2$, which is negative because
$\phi(t)<2/5$, $C\le 1/2$, and $B(t)>9/20$. Equations \eqref{contact:two-triple-nonpositive-rectangle-7}--\eqref{contact:two-triple-nonpositive-rectangle-8}
therefore imply $1<U(t)$, where

\[
 U(t)=\frac{7B(t)^2}{4(2/5)^2}-\frac{3\phi(t)}{2/5},
 \qquad 0\leq t\leq6/5.
\]

This scalar function has no interior maximum. In fact, with $c_0=2/5$,

\[
 \frac{U'(t)}{\phi(t)}
 =\frac{7B(t)(t^2-1)}{6c_0^2}+\frac{3t}{c_0},
\]

whose derivative is

\[
 \frac7{6c_0^2}
 \left\{\frac{(t^2-1)^2\phi(t)}3+2tB(t)\right\}
 +\frac3{c_0}>0.
\]

Thus $U'$ changes sign at most once, from negative to positive. At the
two endpoints, the elementary bounds $\phi(0)>39/100$ and
$\phi(6/5)>19/100$ give

\[
 U(0)<\frac{46087}{160000},\qquad
 U(6/5)<\frac{152023}{160000}<1.
\]

For the second comparison, the expression defining $U$ is decreasing in
$\phi(t)$ at fixed $t$, so replacing $\phi(6/5)$ by $19/100$ gives an
upper bound. This contradicts $1<U(t)$ and completes the exclusion.

The numerical-looking Gaussian comparisons above require no floating
point arithmetic. For example $\pi<22/7$ gives $\phi(0)>199/500$.
At $x=18/25$, the Taylor bound

\[
 e^x\leq1+x+x^2/2+x^3/6+
              \frac{x^4}{24(1-x/5)}<25/12
\]

gives $\phi(6/5)>(199/500)(12/25)>19/100$.
Also $\pi e>8$ gives $\phi(1)<1/4$, and $\pi>25/8$ gives
$\phi(0)<2/5$.
\subsection{Two triples: positive rectangle representations}\label{contact:two-triple-rectangle-exclusion-continued}
Return to the contact-scale coordinates $(-a,c,b)$ of Section~\ref{contact:two-triple-rectangle-exclusion}.



\paragraph{Positive representations}

Suppose $c>0$ in the identical normalized laws obtained above. Omit the
coordinate index. The rectangle condition is

\[
 b-a>2c,\qquad\text{equivalently }v=b-c>u=a+c.
\tag{C.8.12}\label{contact:two-triple-rectangle-exclusion-12}
\]

The inequality is strict because a collision between a corner sum and
the selected middle-middle sum would give a second representation of the threshold.

Each coordinate has variance $1/2$. Writing its middle mass as $m$,
centering and the exact variance formula give

\[
 \frac12=K^2(ab-muv),\qquad q=\frac{b+mu}{h}.
\]

Eliminating $m$ yields

\[
 q=1+\frac{ac}{hv}-\frac1{2hvK^2}.
\tag{C.8.13}\label{contact:two-triple-rectangle-exclusion-13}
\]

By \eqref{contact:two-triple-rectangle-exclusion-1}, $h\le 2$, $ab\le 1$, and $b<(1+\sqrt{2})/2<5/4$.
Also $v<b$, while \eqref{contact:two-triple-rectangle-exclusion-12} gives $hv>2u^2$. Therefore

\[
 hv<5/2,\qquad
 \frac{ac}{hv}<\frac{ac}{2(a+c)^2}\leq\frac18,
\]

and \eqref{contact:two-triple-rectangle-exclusion-13} implies

\[
 q<Q(K):=\frac98-\frac1{5K^2},\qquad K\geq1/\sqrt2.
\tag{C.8.14}\label{contact:two-triple-rectangle-exclusion-14}
\]

The selected threshold is $z=2Kc>0$, and the event probability is $q^2$.
Since $C R=C K-z \phi(z)/3$, its ratio identity would require

\[
 0=q^2-\Phi(z)-CK+z\phi(z)/3
   <q^2-\frac12-\frac25K.
\tag{C.8.15}\label{contact:two-triple-rectangle-exclusion-15}
\]

Here $C\ge49/120>2/5$, and $\Phi(z)\ge 1/2+z \phi(z)$ for $z\ge 0$.
For $K\ge 5/4$, the right side of \eqref{contact:two-triple-rectangle-exclusion-15} is at most
$1/2-(2/5)K\le 0$, since $q\le 1$, giving a contradiction.

For $1/\sqrt{2}\le K\le 5/4$, define

\[
 G(K)=Q(K)^2-1/2-(2/5)K.
\]

Its second derivative is

\[
 G''(K)=\frac4{5K^4}\left(\frac1{K^2}-\frac{27}8\right)<0.
\]

Also $G'(5/4)=654/78125>0$. Thus $G$ increases throughout this
interval and

\[
 G(K)\leq G(5/4)=\left(\frac{997}{1000}\right)^2-1
 =-\frac{5991}{10^6}<0.
\]

Equations \eqref{contact:two-triple-rectangle-exclusion-14}--\eqref{contact:two-triple-rectangle-exclusion-15} again give a contradiction. This excludes positive
representations of the $220$ event; the only numerical input is \eqref{contact:two-triple-rectangle-exclusion-11}.
\subsection{One triple: the rectangle}\label{contact:one-triple-bernoulli-rectangle}


\paragraph{Normalized contacts and exact variance}

Normalize total variance to one, and use the dimensionless supports

\[
 T:\quad K(-a,c,b),\qquad B:\quad K(-d,e),
 \qquad K=R+\frac{z\phi(z)}{3C}.
\]

The representation is the middle triple point and the high Bernoulli point, so
$z=K(c+e)$. The nonselected triple points are smooth minima with
deletion levels $1,0$; the representation has level $1$. Consequently

\[
 \frac12<a<1<b,\quad a^2+b^2=a+b,\quad
 f_T(Kc)=f_T(-Ka),\quad -\frac12<c<1.
\tag{C.9.1}\label{contact:one-triple-bernoulli-rectangle-1}
\]

Both Bernoulli deletion levels are the same, namely $q=\Pp(T\le Kc)$.
Its low point is a smooth minimum, and its high point has negative
derivative. The selected Bernoulli point identity gives

\[
 d>\frac12,\quad \frac{d^2+e^2}{d+e}<1.
\]

Set $k=e/d>0$. Factoring the equal-height equation first, without yet
restricting $k$, gives

\[
 d=d(k)=\frac{3(1+k)^2}{2(k^3+3k+2)},\qquad e=kd(k),
 \qquad k>0.
\tag{C.9.2}\label{contact:one-triple-bernoulli-rectangle-2}
\]

Substitution into the strict Bernoulli representation inequality then gives

\[
 1-\frac{d^2+e^2}{d+e}
 =\frac{(1-k)(k^2+4k+1)}{2(k^3+3k+2)}>0.
\]

Hence $k<1$. Differentiating \eqref{contact:one-triple-bernoulli-rectangle-2} on $(0,1)$ shows $d'(k)>0$, while
$d(0)=3/4$ and $d(1)=1$. Consequently $3/4<d<1$ and $0<e<d$.

This is the same directly factored equation as in
Section~\ref{contact:two-triple-rectangle-exclusion},
with the same factorization.

Let $m,r$ be the triple's middle and high masses. Centering gives

\[
 r=\frac{a-m(a+c)}{a+b}>0,
 \qquad \operatorname{Var}(T)
   =K^2\{ac+r(a+b)(b-c)\}.
\]

Writing $h=a+b$, exact total variance is therefore

\[
 \boxed{\quad 1=K^2\{ac+de+r h(b-c)\}.\quad}
\tag{C.9.3}\label{contact:one-triple-bernoulli-rectangle-3}
\]

The full event is simply $T\le Kc$, so $F=1-r$.

\paragraph{The selected standardized threshold is positive}

If $c>0$ this is immediate. Suppose $c\le 0$. The triple equal-height
factorization yields

\[
 \frac12<a\leq\frac34,\qquad c=2a-\frac32,
 \qquad H_T=-f_T'(Kc)/(CK^2)=(a+c)^2.
\]

Put $\sigma=\phi(z)/(CK^2)$. The unique-representation derivative identity is
$m H_T=\sigma$. Since the high mass in \eqref{contact:one-triple-bernoulli-rectangle-3} is positive,

\[
 \sigma<a(a+c)\leq\frac9{16}.
\tag{C.9.4}\label{contact:one-triple-bernoulli-rectangle-4}
\]

For the Bernoulli law, its smooth low-point derivative instead gives

\[
 \sigma=3d\left(1-\frac{d^2+e^2}{d+e}\right)
 =\frac{9(1+k)^2(1-k)(k^2+4k+1)}
        {4(k^3+3k+2)^2}.
\tag{C.9.5}\label{contact:one-triple-bernoulli-rectangle-5}
\]

We claim that \eqref{contact:one-triple-bernoulli-rectangle-4}--\eqref{contact:one-triple-bernoulli-rectangle-5} imply $e>1/2$. The functions $d(k)$ and
$e(k)=kd(k)$ strictly increase on $(0,1)$, by differentiating \eqref{contact:one-triple-bernoulli-rectangle-2}, and

\[
 e(11/20)=\frac{10571}{20354}>\frac12.
\]

If $e\le 1/2$, then $0<k<11/20$. The difference between the numerator
in \eqref{contact:one-triple-bernoulli-rectangle-5} and the value $9/16$ has the sign of

\[
 k B(k),\qquad B(k)=8+7k-20k^2-26k^3-4k^4-k^5.
\]

The polynomial $B$ is strictly concave on $[0,11/20]$, because
$B''=-40-156k-48k^2-20k^3<0$, and its two endpoint values are

\[
 B(0)=8,\qquad B(11/20)=\frac{3385269}{3200000}>0.
\]

Thus $\sigma>9/16$, contradicting \eqref{contact:one-triple-bernoulli-rectangle-4}. Hence $e>1/2$. Since
$c>-1/2$, we have proved in both sign cases that

\[
 \boxed{\quad z=K(c+e)>0.\quad}
\tag{C.9.6}\label{contact:one-triple-bernoulli-rectangle-6}
\]

\paragraph{Elementary uniform bounds}

For every centered coordinate, its third absolute moment is at least
its variance to the power $3/2$. With two coordinates and total variance
one, this gives $R\ge 1/\sqrt{2}>7/10$. Also $C\ge49/120$, so
$C^2>1/6$ and $CR>343/1200$. Since $F\le 1$,

\[
 \Phi(z)<\frac{857}{1200}.
\]

The elementary bounds $\phi(0)>39/100$ and $\exp(-x)\ge 1-x$ give

\[
 \Phi(3/5)>
 \frac12+\frac{39}{100}\left\{\frac35-
                         \frac{(3/5)^3}{6}\right\}
 =\frac{17999}{25000}>\frac{857}{1200}.
\]

Consequently $0<z<3/5$. We will also use

\[
 \phi(z)>\phi(3/5)>\frac{33}{100}.
\tag{C.9.7}\label{contact:one-triple-bernoulli-rectangle-7}
\]

For a rational verification of the last bound, $\pi<22/7$ implies
$\phi(0)>199/500$; the cubic alternating exponential lower bound gives

\[
 e^{-9/50}>1-\frac9{50}+\frac{(9/50)^2}{2}
                    -\frac{(9/50)^3}{6}
 =\frac{208807}{250000}.
\]

Their product is $41552593/125000000>33/100$.

Set $N(z)=\Phi(z)-z \phi(z)/3$. The ratio identity is

\[
 1-r=N(z)+CK,
\]

and the decreasing positive Gaussian density gives
$N(z)\ge 1/2+(2/3)z \phi(z)$. Thus

\[
 K<\frac1{2C}<\frac54,
 \qquad r<\frac12-CK-\frac23z\phi(z).
\tag{C.9.8}\label{contact:one-triple-bernoulli-rectangle-8}
\]

\paragraph{Variance contradiction}

We need only the following support estimates:

\[
 ac+de<c+e+\frac14,\qquad h(b-c)<3.
\tag{C.9.9}\label{contact:one-triple-bernoulli-rectangle-9}
\]

For the first, $d<1$ and $e>0$ give $de<e$. If $c>0$, also
$ac<c$; if $c\le 0$, then
$ac=c+(1-a)|c|<c+1/4$, by \eqref{contact:one-triple-bernoulli-rectangle-1}.

For the second, when $c>0$ we have $a>3/4$, and

\[
 h(b-c)<hb\leq\frac{17+7\sqrt7}{16}<\frac94<3.
\]

Here $(hb)'=h(3-4a)/(2b-1)\le 0$, so the bound is its value at $a=3/4$.
When $c\le 0$, substitute $c=2a-3/2$. Direct differentiation gives

\[
 \{h(b-c)\}'=
 \frac{2(a^2-b^2-6ab+3b)}{2b-1}<0
 \qquad(1/2<a\le3/4),
\]

because $a^2<b^2$ and $6ab\ge 3b$. Its upper endpoint at $a=1/2$ is
$(1+\sqrt{2}/2)^2=3/2+\sqrt{2}<3$. This proves \eqref{contact:one-triple-bernoulli-rectangle-9}.

Insert \eqref{contact:one-triple-bernoulli-rectangle-8}--\eqref{contact:one-triple-bernoulli-rectangle-9} into the exact variance \eqref{contact:one-triple-bernoulli-rectangle-3}:

\[
 \begin{aligned}
 1
 &<Kz+\frac{K^2}{4}+3K^2r\\
 &<\frac{K^2}{4}+3K^2(1/2-CK)
       +zK\{1-2K\phi(z)\}.
 \end{aligned}
\tag{C.9.10}\label{contact:one-triple-bernoulli-rectangle-10}
\]

Each term has an elementary upper bound. Equations \eqref{contact:one-triple-bernoulli-rectangle-7}--\eqref{contact:one-triple-bernoulli-rectangle-8} give

\[
 \frac{K^2}{4}<\frac{25}{64},\qquad
 3K^2(1/2-CK)\leq\frac1{18C^2}<\frac13.
\]

For the last term in \eqref{contact:one-triple-bernoulli-rectangle-10}, if its brace is nonpositive the term is
already nonpositive. Otherwise, using $z<3/5$ and $\phi(z)>33/100$,

\[
 zK\{1-2K\phi(z)\}
 <\frac35\max_{x\geq0}\left(x-\frac{33}{50}x^2\right)
 =\frac5{22}.
\]

Therefore \eqref{contact:one-triple-bernoulli-rectangle-10} yields the impossible strict inequality

\[
 \boxed{\quad 1<\frac{25}{64}+\frac13+\frac5{22}
           =\frac{2009}{2112}<1.\quad}
\]

This completes the exclusion of all unique-representation $220$ configurations.
\subsection{One triple: the upper shape}\label{contact:one-triple-bernoulli-upper-shape}

We exclude a maximizing pair with a unique middle/high representation
and event $221$. Use the contact scale:


\[
 T:\quad K(-a,c,b),\qquad B:\quad K(-d,e),\qquad
 K=R+\frac{z\phi(z)}{3C}.
\]

The selected representation is middle/high, with $z=K(c+e)$. Let $r$ be the
triple's high mass, let $p=d/(d+e)$ be the Bernoulli high mass, and put
$\varepsilon=rp$. The actual event probability is

\[
 F=1-rp=1-\varepsilon.
\tag{C.10.1}\label{contact:one-triple-bernoulli-upper-shape-1}
\]

\paragraph{Contact signs force a positive threshold}

The triple deletion levels are $(1,1,1-p)$, so the two nonselected
smooth minima and the selected middle point satisfy

\[
 \frac12<a<1<b,\quad a^2+b^2=a+b,\quad
 f_T(Kc)=f_T(-Ka),\quad -\frac12<c<1.
\tag{C.10.2}\label{contact:one-triple-bernoulli-upper-shape-2}
\]

The Bernoulli low point is a smooth minimum, and its high representation has
negative derivative. The representation identity gives

\[
 d>\frac12,\qquad \frac{d^2+e^2}{d+e}<1.
\tag{C.10.3}\label{contact:one-triple-bernoulli-upper-shape-3}
\]

We also need $d<1$, which does not follow from \eqref{contact:one-triple-bernoulli-upper-shape-3} alone. The actual
Bernoulli deletion levels are $1,1-r$. Set $t=e/d>0$ and
$P(t)=t^3+3t+2$. Subtracting the two contact heights, using the smooth
derivative equation at $-Kd$, gives exactly

\[
 -r=f_B(Ke)-f_B(-Kd)
   =CK^3 d^2 P(t)
      \left\{d-\frac{3(1+t)^2}{2P(t)}\right\}<0.
\tag{C.10.4}\label{contact:one-triple-bernoulli-upper-shape-4}
\]

Since

\[
 2P(t)-3(1+t)^2=(t-1)^2(2t+1)\geq0,
\]

equation \eqref{contact:one-triple-bernoulli-upper-shape-4} gives $d<1$. The strict grid geometry for $221$ is
$d+e>b-c$. Consequently

\[
 c+e>b-d>0,\qquad z>0.
\tag{C.10.5}\label{contact:one-triple-bernoulli-upper-shape-5}
\]

Moreover \eqref{contact:one-triple-bernoulli-upper-shape-3} and $d>1/2$ imply

\[
 \frac{1+t^2}{1+t}<2,
 \qquad 0<t<1+\sqrt2<\frac{17}{7}.
\tag{C.10.6}\label{contact:one-triple-bernoulli-upper-shape-6}
\]

\paragraph{Universal scalar bounds and exact variance}

The elementary Gaussian comparisons in
Section~\ref{contact:one-triple-bernoulli-rectangle},
apply here. They only use $F\le 1$, $R\ge 1/\sqrt{2}$,
$C\ge49/120$, and $z>0$, all valid here. Thus

\[
 0<z<\frac35,
 \qquad \phi(z)>\frac{33}{100},
 \qquad C^2>\frac16.
\tag{C.10.7}\label{contact:one-triple-bernoulli-upper-shape-7}
\]

Writing $N(z)=\Phi(z)-z \phi(z)/3$, the ratio identity \eqref{contact:one-triple-bernoulli-upper-shape-1} is

\[
 1-\varepsilon=N(z)+CK.
\]

Since $N(z)\ge 1/2+(2/3)z \phi(z)$,

\[
 \varepsilon<\frac12-CK-\frac23z\phi(z).
\tag{C.10.8}\label{contact:one-triple-bernoulli-upper-shape-8}
\]

Centering gives the exact variance formula

\[
 1=K^2\{ac+de+r(a+b)(b-c)\}.
\tag{C.10.9}\label{contact:one-triple-bernoulli-upper-shape-9}
\]

The equal-height triple contact in \eqref{contact:one-triple-bernoulli-upper-shape-2} gives the support estimate
$(a+b)(b-c)<3$, proved separately on the two sign branches in
Section~\ref{contact:one-triple-bernoulli-rectangle}. Also

\[
 ac+de<c+e+\delta,
 \qquad \delta=\frac14-(1-d)e.
\tag{C.10.10}\label{contact:one-triple-bernoulli-upper-shape-10}
\]

Indeed, $(a-1)c<1/4$: it is negative for positive $c$, and is less
than $1/4$ when $-1/2<c\le 0$ and $1/2<a<1$. Let

\[
 L=3(1+t),\qquad p=\frac1{1+t}.
\]

Then \eqref{contact:one-triple-bernoulli-upper-shape-1}, \eqref{contact:one-triple-bernoulli-upper-shape-9}--\eqref{contact:one-triple-bernoulli-upper-shape-10} imply

\[
 1<Kz+\delta K^2+L K^2\varepsilon.
\tag{C.10.11}\label{contact:one-triple-bernoulli-upper-shape-11}
\]

\paragraph{A sharper intercept when the high Bernoulli mass is small}

For $0<t\le 1$, simply use $\delta<1/4$. For $1\le t<17/7$, the slope
condition \eqref{contact:one-triple-bernoulli-upper-shape-3} gives

\[
 \frac12<d<d_*:=\frac{1+t}{1+t^2}\leq1.
\]

The function $d(1-d)$ decreases on $(1/2,1)$. Hence

\[
 (1-d)e=t d(1-d)
 >t d_*(1-d_*)
 =\frac{t^2(t^2-1)}{(1+t^2)^2}
 \geq\frac5{12}(t-1).
\tag{C.10.12}\label{contact:one-triple-bernoulli-upper-shape-12}
\]

The last inequality is a direct polynomial comparison. After dividing
by $t-1$ when $t>1$, its numerator is

\[
 P_0(t)=-5t^4+12t^3+2t^2-5.
\]

Its derivative $4t(-5t^2+9t+1)$ changes sign at most once on the
positive axis, from positive to negative. Thus its minimum on $[1,17/7]$
is at an endpoint, and

\[
 P_0(1)=4,\qquad P_0(17/7)=\frac{11404}{2401}>0.
\]

The assertion also holds at $t=1$ directly. We have proved

\[
 \delta<\bar\delta(t):=
 \begin{cases}
  1/4,&0<t\leq1,\\
  2/3-5t/12,&1\leq t<17/7.
 \end{cases}
\tag{C.10.13}\label{contact:one-triple-bernoulli-upper-shape-13}
\]

\paragraph{Maximization of a cubic variance bound}

Insert \eqref{contact:one-triple-bernoulli-upper-shape-8}, \eqref{contact:one-triple-bernoulli-upper-shape-13} into \eqref{contact:one-triple-bernoulli-upper-shape-11}, and put
$A=\bar\delta(t)+L/2>0$. We obtain

\[
 1<A K^2-LCK^3
      +zK\left\{1-\frac23 L K\phi(z)\right\}.
\tag{C.10.14}\label{contact:one-triple-bernoulli-upper-shape-14}
\]

The first two terms have the exact maximum

\[
 \max_{K\geq0}(A K^2-LCK^3)
  =\frac{4A^3}{27L^2C^2}
  <\frac{8A^3}{9L^2},
\tag{C.10.15}\label{contact:one-triple-bernoulli-upper-shape-15}
\]

using \eqref{contact:one-triple-bernoulli-upper-shape-7}. For the remaining term, a nonpositive brace gives a
nonpositive contribution. When the brace is positive, \eqref{contact:one-triple-bernoulli-upper-shape-7} gives

\[
 zK\left\{1-\frac23L K\phi(z)\right\}
 <\frac35\max_{x\geq0}
        \left(x-\frac{11}{50}Lx^2\right)
 =\frac{15}{22L}.
\tag{C.10.16}\label{contact:one-triple-bernoulli-upper-shape-16}
\]

It remains to check the two explicit upper functions in
\eqref{contact:one-triple-bernoulli-upper-shape-15}--\eqref{contact:one-triple-bernoulli-upper-shape-16}. Write $x=1+t$.

For $0<t\le 1$, we have $1<x\le 2$, $L=3x$, and $A=(6x+1)/4$.
The bound equals

\[
 H_0(x)=\frac{(6x+1)^3}{648x^2}+\frac5{22x}
 =\frac{x}{3}+\frac16+\frac{101}{396x}+\frac1{648x^2}.
\]

This function is strictly convex. Its maximum on $[1,2]$ is at an
endpoint, with the larger value at $x=2$:

\[
 H_0(x)\leq H_0(2)=\frac{27407}{28512}<1.
\tag{C.10.17}\label{contact:one-triple-bernoulli-upper-shape-17}
\]

For $1\le t<17/7$, we have $2\le x<24/7$ and
$A=(13/12)(x+1)$. The bound is

\[
 H_1(x)=\frac{2197}{17496}\frac{(x+1)^3}{x^2}
                  +\frac5{22x}.
\]

Put $c_0=2197/17496$. Then

\[
 H_1(x)=c_0 x+3c_0+\frac{3c_0+5/22}{x}+\frac{c_0}{x^2},
\]

again strictly convex. Its two endpoint values are

\[
 H_1(2)=\frac{27407}{28512}<1,
 \qquad H_1(24/7)=\frac{771397337}{775982592}<1.
\tag{C.10.18}\label{contact:one-triple-bernoulli-upper-shape-18}
\]

The second endpoint is larger, and its positive deficit is
$4585255/775982592$. Convexity bounds the whole interval by those
endpoints. Equations \eqref{contact:one-triple-bernoulli-upper-shape-14}--\eqref{contact:one-triple-bernoulli-upper-shape-18} force $1<1$, a contradiction. This
excludes every unique-representation $221$ configuration .
\subsection{One triple: the middle representation of the triangular event}\label{contact:one-triple-bernoulli-middle-210}


Consider shape $210$ with the middle-triple/low-Bernoulli representation.

Normalize total variance to one and write the dimensionless supports as

\[
 T:\quad K(-a,c,b),\qquad B:\quad K(-d,e),\qquad
 K=R+\frac{z\phi(z)}{3C}>0.
\]

Let the triple masses be $l,m,r$, all positive. The selected threshold
is $z=K(c-d)$. The actual deletion levels are

\[
 T:\quad (1,q_B,0),\qquad
 B:\quad (l+m,l),\qquad q_B=\frac{e}{d+e}.
\tag{C.11.1}\label{contact:one-triple-bernoulli-middle-210-1}
\]

The strict representation geometry is

\[
 d+e<a+c.
\tag{C.11.2}\label{contact:one-triple-bernoulli-middle-210-2}
\]

\paragraph{The positive Bernoulli minimum must be at least K}

The nonselected Bernoulli point $Ke$ is a smooth minimum, so $e>1/2$.
Suppose $e<1$. Its cubic dual then also has a smooth negative minimum
at $-K \alpha$, where

\[
 \alpha=\frac{1+\sqrt{1+4e-4e^2}}2>1,
 \qquad \alpha^2+e^2=\alpha+e.
\]

The difference of the minimum heights is exactly

\[
 f_B(-K\alpha)-f_B(Ke)
 =\frac{CK^3}{2}(e-\alpha)
                     (\alpha^2+4\alpha e+e^2)<0.
\tag{C.11.3}\label{contact:one-triple-bernoulli-middle-210-3}
\]

But $f_B(Ke)=l$ by \eqref{contact:one-triple-bernoulli-middle-210-1}. The deletion threshold at the new negative
minimum is

\[
 z+K\alpha=K(c-d+\alpha)>K(-a+e+\alpha)>-Ka,
\]

using \eqref{contact:one-triple-bernoulli-middle-210-2}. Its deletion CDF is therefore at least $l$. Equation \eqref{contact:one-triple-bernoulli-middle-210-3}
contradicts the global cubic majorant at this point. Hence

\[
 \boxed{\quad e\geq1.\quad}
\tag{C.11.4}\label{contact:one-triple-bernoulli-middle-210-4}
\]

The nonselected triple minima have levels $1,0$, so

\[
 \frac12<a<1<b,\qquad a^2+b^2=a+b.
\tag{C.11.5}\label{contact:one-triple-bernoulli-middle-210-5}
\]

Equations \eqref{contact:one-triple-bernoulli-middle-210-2}, \eqref{contact:one-triple-bernoulli-middle-210-4}--\eqref{contact:one-triple-bernoulli-middle-210-5} give

\[
 c-d>e-a>0,\qquad z>0,
\]

and in particular $c>0$.

\paragraph{Contact height bounds the negative triple minimum}

Define the dimensionless minimum-height gap

\[
 D(a)=\frac12(b-a)(a^2+4ab+b^2),\qquad
 b=\frac{1+\sqrt{1+4a-4a^2}}2.
\]

The levels $1,0$ in \eqref{contact:one-triple-bernoulli-middle-210-1} give the exact identity

\[
 CK^3 D(a)=1.
\tag{C.11.6}\label{contact:one-triple-bernoulli-middle-210-6}
\]

The scalar comparisons in Section~\ref{contact:one-triple-bernoulli-rectangle} give

\[
 0<z<\frac35,\qquad \phi(z)>\frac{33}{100},\qquad
 C^2>\frac16,
\tag{C.11.7}\label{contact:one-triple-bernoulli-middle-210-7}
\]

because $R\ge 1/\sqrt{2}$, $F\le 1$, and $C\ge49/120$. Also
$F=N(z)+CK$, where $N(z)=\Phi(z)-z \phi(z)/3>1/2$; hence $CK<1/2$.
Together with \eqref{contact:one-triple-bernoulli-middle-210-6},

\[
 1=CK^3D(a)<\frac{D(a)}{8C^2},\qquad D(a)>8C^2>\frac43.
\tag{C.11.8}\label{contact:one-triple-bernoulli-middle-210-8}
\]

Direct differentiation gives

\[
 D'(a)=-3(2a-1)(a+b)<0.
\]

At $a=2/3$, its value is

\[
 D(2/3)=\frac{71+17\sqrt{17}}{108}
 <\frac{191}{144}<\frac43,
\]

using $\sqrt{17}<17/4$. Therefore \eqref{contact:one-triple-bernoulli-middle-210-8} forces

\[
 \boxed{\quad a<\frac23.\quad}
\tag{C.11.9}\label{contact:one-triple-bernoulli-middle-210-9}
\]

\paragraph{The selected triple gap is bounded away from one}

\[
 D(a)<D(1/2)=\frac34+\frac{\sqrt2}{2}<\frac32.
\]

Equation \eqref{contact:one-triple-bernoulli-middle-210-6} consequently gives

\[
 K^3>\frac{1000}{681}>\left(\frac{25}{22}\right)^3.
\tag{C.11.10}\label{contact:one-triple-bernoulli-middle-210-10}
\]

The second strict difference is the positive rational
$7375/7251288$. Put $\sigma=\phi(z)/(CK^2)$. From \eqref{contact:one-triple-bernoulli-middle-210-6}--\eqref{contact:one-triple-bernoulli-middle-210-8}, \eqref{contact:one-triple-bernoulli-middle-210-10},

\[
 \sigma=\phi(z)D(a)K
 >\frac{33}{100}\frac43\frac{25}{22}=\frac12.
\tag{C.11.11}\label{contact:one-triple-bernoulli-middle-210-11}
\]

The Bernoulli low-representation identity is

\[
 \sigma=3e\left\{\frac{d^2+e^2}{d+e}-1\right\}>0.
\]

Let $u=a+c$. Since the Bernoulli third-moment slope is strictly smaller
than its support diameter, and \eqref{contact:one-triple-bernoulli-middle-210-2} gives $d+e<u$,

\[
 \sigma<3u(u-1).
\]

Combining this with \eqref{contact:one-triple-bernoulli-middle-210-11}, and noting $3(8/7)(1/7)=24/49<1/2$,
we obtain

\[
 \boxed{\quad u>\frac87.\quad}
\tag{C.11.12}\label{contact:one-triple-bernoulli-middle-210-12}
\]

\paragraph{The actual variance Hessian gives the opposite bound}

Let $\gamma=[-Ka,Kc,Kb]|x|^3$ be the third-moment slope in the centered
triple's variance direction. Because $c>0$, its normalized value is

\[
 g(a,u):=\frac\gamma K
   =b+u-2a+\frac{2a^3}{(a+b)u}.
\tag{C.11.13}\label{contact:one-triple-bernoulli-middle-210-13}
\]

On $1/2<a\le 2/3$, $u\ge 8/7$, this function increases with $u$ and
decreases with $a$. For the first derivative,

\[
 g_u=1-\frac{2a^3}{(a+b)u^2}>0,
\]

since $a+b>2a$ and $u>a$. For the second, write $h=a+b$.
Here $b'<0$, $h'>0$, and

\[
 g_a=b'-2+\frac2u\left(\frac{3a^2}{h}
                          -\frac{a^3h'}{h^2}\right)
 <-2+\frac{6a^2}{uh}
 <-2+\frac{14}{9}<0,
\]

using $a\le 2/3$, $h>3/2$, $u\ge 8/7$. Equations \eqref{contact:one-triple-bernoulli-middle-210-9}, \eqref{contact:one-triple-bernoulli-middle-210-12} therefore imply

\[
 \frac\gamma K>g(2/3,8/7)
   =\frac{499+35\sqrt{17}}{504}
   >\frac{71}{56}>\frac54.
\tag{C.11.14}\label{contact:one-triple-bernoulli-middle-210-14}
\]

The strict lower bound uses only $\sqrt{17}>4$.

On the other hand, both signs of the centered triple variance direction
are admissible, because all three masses are positive. The actual
second-order condition derived in
Section~\ref{contact:three-point-variance-curvature}
is

\[
 3C\gamma-\frac{15CR+\phi'''(z)}4\leq0.
\]

Use $R=K-z \phi(z)/(3C)$ and
$\phi'''(z)=(3z-z^3)\phi(z)$ to rewrite it as

\[
 12C\gamma\leq15CK-z(z^2+2)\phi(z)<15CK,
\]

because $z>0$. Thus $\gamma/K<5/4$, contradicting \eqref{contact:one-triple-bernoulli-middle-210-14}.
This excludes the middle/low $210$ representation .
\subsection{One triple: the low representation of the triangular event}\label{contact:low-210-algebraic-reduction}


For the low-triple/high-Bernoulli representation of $210$, divide support coordinates by
$K=R+z \phi(z)/(3C)>0$.

\paragraph{Exact contact reduction}

Write

\[
 T:\quad -(a+w),-a,b,\qquad B:\quad -d,e,
 \qquad w>0,
\]

with triple masses $\ell,m,r$. The triple's two nonselected minima have
levels $q,0$, while its selected lowest point has level $1$.
The Bernoulli masses at $-d,e$ are $q,p$, respectively, and its deletion
levels are $\ell+m,\ell$. The representation geometry is

\[
 w<d+e<a+b+w.
\tag{C.12.1}\label{contact:low-210-algebraic-reduction-1}
\]

The smooth triple minima give

\[
 \frac12<a<1<b,\quad a^2+b^2=a+b,
 \quad b=\frac{1+\sqrt{1+4a-4a^2}}2.
\]

Put

\[
 h=a+b,\quad
 D=\frac12(b-a)(a^2+4ab+b^2),\quad
 P=w^2(3a-3/2+w),\quad \tau=CK^3.
\]

The minimum-height gap is $q=\tau D$. Factoring the negative cubic at
its smooth minimum gives the other level gap as
$p=\tau P$. Therefore

\[
 \boxed{\quad \tau=\frac1{D+P},\qquad
 p=\frac P{D+P},\qquad q=\frac D{D+P}.\quad}
\tag{C.12.2}\label{contact:low-210-algebraic-reduction-2}
\]

The Bernoulli negative smooth minimum has $d>1/2$. Its high-representation
derivative gives $k_B=(d^2+e^2)/(d+e)<1$. The strict level difference
$f_B(e)-f_B(-d)=-m<0$ also gives $d<1$, by the exact factorization
in Section~\ref{contact:one-triple-bernoulli-upper-shape}.

Thus, after choosing $a,w,d$, every quantity is determined. Write

\[
 x=d/p=d+e,\quad e=qx,\quad r_B=p^2+q^2,\quad k_B=r_B x,
\]

\[
 \ell=\frac{d(1-k_B)}{w(2a-1+w)},\quad
 m=\frac{b-\ell(h+w)}h,\quad r=1-\ell-m.
\tag{C.12.3}\label{contact:low-210-algebraic-reduction-3}
\]

The first formula follows by equating the two unique-representation slope
identities:

\[
 \sigma=3\ell w(2a-1+w)=3d(1-k_B),
 \qquad \sigma=\frac{\phi(z)}{CK^2}.
\]

The other two formulas in \eqref{contact:low-210-algebraic-reduction-3} are centering and total mass.

It remains to impose the exact Bernoulli height difference. Put

\[
 M=q^3+3p^2q+2p^3.
\]

Using its smooth derivative at $-d$, the difference is

\[
 \boxed{\quad m=\tau x^2(3/2-Mx).\quad}
\tag{C.12.4}\label{contact:low-210-algebraic-reduction-4}
\]

For fixed $a,w$, this is a cubic equation in $x$. In particular the
CDF levels and coordinate probabilities have not been replaced by
independent optimization parameters.

\paragraph{Variance and two necessary ratio inequalities}

The dimensionless variances are

\[
 V_T=ab+\ell w(h+w),\qquad V_B=pq x^2,\qquad V_0=V_T+V_B.
\]

Total variance one gives $K=V_0^{-1/2}$, hence

\[
 C=\tau V_0^{3/2},\qquad
 \boxed{\quad \tau^2V_0^3<(227/500)^2.\quad}
\tag{C.12.5}\label{contact:low-210-algebraic-reduction-5}
\]

The dimensionless Bernoulli third absolute moment is
$R_B=pq x^3(p^2+q^2)$. The triple dual has constant term
$\tau B_0$, where $B_0=b^2(2b-3/2)$, since it vanishes at its positive
smooth minimum $b$. Averaging that dual over the centered triple and
using $F-CR=\Phi(z)>0$ gives the further necessary inequality

\[
 \boxed{\quad B_0-\frac32 V_T-R_B>0.\quad}
\tag{C.12.6}\label{contact:low-210-algebraic-reduction-6}
\]

\paragraph{Proof of the compact closed domain}

Since $V_T>ab$ and $R_B>0$, equation \eqref{contact:low-210-algebraic-reduction-6} implies

\[
 0<B_0-\frac32ab=\frac{b}{2}(b^2-3a^2).
\]

If $a\ge 3/4$, then $b<5/4$ gives
$b^2<25/16<27/16\le 3a^2$, a contradiction. Thus $a<3/4$.
The strictly decreasing height function satisfies

\[
 1<D(3/4)<D(a)<D(1/2)<3/2.
\]

Here $D(3/4)=(17+7\sqrt{7})/32>1$, and
$D(1/2)=3/4+\sqrt{2}/2<3/2$.

The inequalities $d>1/2$ and $k_B<1$ give
$e/d<1+\sqrt{2}<3$, hence $q<3/4$. Equation \eqref{contact:low-210-algebraic-reduction-2} implies
$P>D/3>1/3$. If $w\le 1/2$, then
$P=w^2(3a-3/2+w)\le (1/4)(3/4+1/2)=5/16$, a contradiction.
Thus $w>1/2$.

Conversely, \eqref{contact:low-210-algebraic-reduction-1} and $d=px<1$ give $pw<1$, or $(w-1)P<D$.
If $w\ge 3/2$, then $P>w^3$ and
$(w-1)P\ge (1/2)(27/8)=27/16>3/2>D$, again a contradiction.
Thus $w<3/2$.

Every candidate therefore lies in the interior of the single closed box

\[
 \boxed{\quad (a,w,d)\in[1/2,3/4]\times[1/2,3/2]\times[1/2,1].\quad}
\tag{C.12.7}\label{contact:low-210-algebraic-reduction-7}
\]

Keeping all its faces makes the computational cover larger than needed.
The fixed certificate \texttt{low210} proves that no point of
\eqref{contact:low-210-algebraic-reduction-7} satisfies all the necessary
contact, probability, variance, and ratio conditions above. Thus this
representation cannot occur at a maximizing pair. The enclosure and
partition rules are given in Section~\ref{sec:verification}.
\subsection{One triple: the lower event}\label{contact:lower-110-contact-reduction}


For shape $110$, the unique selected representation is the triple's middle point and the Bernoulli's low point.

\paragraph{Two equal-height charts}

Divide all support points by
$K=R+z \phi(z)/(3C)>0$, and write

\[
 T:\quad -a,c,b,\qquad B:\quad -d,e.
\]

The triple's deletion levels are $(q,q,0)$, where $q=e/(d+e)$.
Thus its nonselected minima and selected point satisfy

\[
 \frac12<a<1<b,\quad a^2+b^2=a+b,\quad f_T(c)=f_T(-a).
\]

Both possible representation signs have exact closed parameter charts $0\le k\le 1$:

\[
 \begin{array}{c|c|c}
 \text{chart}&a&c\\ \hline
 \text{negative}&1/2+k/4&2a-3/2\\
 \text{positive}&3(1+k)^2/[2(k^3+3k+2)]&ka.
 \end{array}
\tag{C.13.1}\label{contact:lower-110-contact-reduction-1}
\]

The two charts share the $c=0$ boundary and cover the entire genuine
equal-height representation range. Their other faces include limiting contacts
that need not be genuine triples. Keep all faces in the certificate.

Put

\[
 b=\frac{1+\sqrt{1+4a-4a^2}}2,\quad h=a+b,\quad
 u=a+c,\quad v=b-c,
\]

\[
 D=\frac12(b-a)(a^2+4ab+b^2),\quad
 H=3\{a(1-a)+c-c|c|\}.
\]

On the negative chart the stable exact formula is $H=u^2$. On the
positive chart it is $H=3[a(1-a)+c(1-c)]$. In either case $0<H\le 3/2$
at a genuine contact.

\paragraph{The middle mass replaces the singular derivative quotient}

Let the triple middle mass be the independent parameter $m$, with
$0<m<1$. Its other masses and its representation derivative parameter are

\[
 \ell=\frac{b-mv}{h},\qquad r=\frac{a-mu}{h},
 \qquad \sigma=mH=\frac{\phi(z)}{CK^2}>0.
\tag{C.13.2}\label{contact:lower-110-contact-reduction-2}
\]

The Bernoulli's nonselected positive point is a smooth minimum at height
zero. Its entire cubic dual majorizes a nonnegative CDF. If $1/2<e<1$,
the second smooth minimum is at the negative point

\[
 -\alpha=-\frac{1+\sqrt{1+4e-4e^2}}2,\qquad \alpha>1,
\]

and its height is strictly below the positive minimum's height zero.
This is impossible for the global majorant. Hence $e\ge 1$.

The low Bernoulli representation identity is

\[
 \sigma=3e\left\{\frac{d^2+e^2}{d+e}-1\right\}.
\tag{C.13.3}\label{contact:lower-110-contact-reduction-3}
\]

Since $\sigma<3/2$ and $e\ge 1$, its third-moment slope is less than $3/2$,
so its diameter is less than three. Consequently $d<2$ and

\[
 0<t:=d/e<2.
\]

Set $f=(t^2+1)/(t+1)$. Solving the positive quadratic root in \eqref{contact:lower-110-contact-reduction-3}
determines the Bernoulli support and probabilities without dividing by $H$:

\[
 \boxed{\quad e=\frac{1+\sqrt{1+(4/3)f\sigma}}{2f},\quad
 d=te,\quad q=\frac1{1+t},\quad k_B=ef.\quad}
\tag{C.13.4}\label{contact:lower-110-contact-reduction-4}
\]

Every actual candidate therefore lies in one of the two closed cubes

\[
 \boxed{\quad (k,m,t)\in[0,1]\times[0,1]\times[0,2].\quad}
\tag{C.13.5}\label{contact:lower-110-contact-reduction-5}
\]

The use of $m$ as a parameter avoids artificial singularities where the
representation derivative $H$ vanishes on a chart face.

\paragraph{Exact height, ratio, and normal-density equations}

The strict $110$ representation geometry is $d+e>u$. The minimum-level difference
gives

\[
 \tau:=CK^3=q/D.
\]

For the Bernoulli dual, whose positive smooth minimum has height zero,
the low-representation contact value is the dimensionless cubic

\[
 L=e^3(t^3+3t+2)-\frac32e^2(1+t)^2.
\]

Its deletion-CDF value is $\ell+m=1-r$. Thus the exact remaining height
equation, cleared of its positive denominator, is

\[
 \boxed{\quad (1-r)D-qL=0.\quad}
\tag{C.13.6}\label{contact:lower-110-contact-reduction-6}
\]

The dimensionless variances and Bernoulli third absolute moment are

\[
 V_T=ab-muv,\quad V_B=de,\quad V_0=V_T+V_B,\quad R_B=de\,k_B.
\]

Total variance one gives $K=V_0^{-1/2}$, $z=(c-d)/\sqrt{V_0}$, and

\[
 C=\frac{qV_0^{3/2}}D.
\]

The established upper and lower ratio bounds require

\[
 \boxed{\quad (49/120)^2D^2\le q^2V_0^3<(227/500)^2D^2.\quad}
\tag{C.13.7}\label{contact:lower-110-contact-reduction-7}
\]

As in the low/high $210$ reduction, averaging the triple cubic and using
$F-CR=\Phi(z)>0$ supplies the strict necessary inequality

\[
 \boxed{\quad b^2(2b-3/2)-\frac32V_T-R_B>0.\quad}
\tag{C.13.8}\label{contact:lower-110-contact-reduction-8}
\]

Finally the definition of $\sigma$ in \eqref{contact:lower-110-contact-reduction-2} is essential: pure algebraic
contacts alone can have ratios below $227/500$. Its exact squared form is

\[
 \boxed{\quad
 \sigma^2 q^2 V_0
 -\frac{D^2}{2\pi}\exp\!\left(-\frac{(c-d)^2}{V_0}\right)=0.
 \quad}
\tag{C.13.9}\label{contact:lower-110-contact-reduction-9}
\]

Squaring is harmless as a necessary condition. It only enlarges the
candidate set, while all actual quantities are positive. Equation \eqref{contact:lower-110-contact-reduction-9}
is evaluated only on boxes with a strictly positive variance enclosure.
Other necessary constraints remain available on wider boxes.
The fixed certificate \texttt{lower110} proves that neither closed
chart in \eqref{contact:lower-110-contact-reduction-5} contains a point
satisfying all these necessary conditions, including the ratio floor
$C=49/120$. Section~\ref{sec:verification} gives the enclosure and
partition rules. This excludes the unique $110$ representation.
\subsection{One triple: collision across the full span}\label{contact:one-triple-bernoulli-collision-211}


The collided $211$ event has its two maximal cells at the selected threshold.

Normalize total variance to one. Write $R$ for the sum of third absolute moments,
$z$ for the threshold, $\phi_0=\phi(0)$, and

\[
 K=R+\frac{z\phi(z)}{3C}>0.
\tag{C.14.1}\label{contact:one-triple-bernoulli-collision-211-1}
\]

The global cubic dual is the one in
Lemma~\ref{lem:full-law-contacts}.
It majorizes the decreasing deletion CDF, touches it at every support
point, and has nonpositive derivative at every contact. At a nonselected
support point its contact is a smooth local minimum. The two threshold representations use the actual marginal masses.

\paragraph{Collision geometry and the sole smooth triple contact}

The selected fiber of the collided $211$ event is exactly
$(triple-low, Bernoulli-high)$ and $(triple-high, Bernoulli-low)$.
The Bernoulli diameter therefore equals the triple's total diameter.
In units of $K$, write the supports as

\[
 T/K:\quad -t,-a,b,\qquad B/K:\quad -d,e,
 \qquad t>a,\quad b,d,e>0.
\]

Set

\[
 h=t+b=d+e,\qquad u=t-a>0,\qquad
 p=\frac d h,\quad q=\frac e h,\qquad \tau=CK^3.
\tag{C.14.2}\label{contact:one-triple-bernoulli-collision-211-2}
\]

The Bernoulli probabilities at its low and high atoms are respectively
$q,p$. The triple deletion levels are exactly $(1,q,q)$.

The middle triple atom is nonselected and is therefore a smooth minimum.
It must be negative. If it were a positive minimum, the cubic dual would
have positive derivative at the still larger triple contact; zero cannot
be a local minimum because its quadratic coefficient is $-3CK/2<0$.
The negative minimum can thus be written $-Ka$, with $a>1/2$.

After dividing its dual values by $\tau$ and subtracting a constant, the
triple dual is

\[
 g(x)=|x|^3-\frac32x^2-3a(1-a)x.
\tag{C.14.3}\label{contact:one-triple-bernoulli-collision-211-3}
\]

The equality of the two contact levels at $-a,b$ gives $g(b)=g(-a)$.
Put $k=b/a>0$ and $P_0=k^3+3k+2$. Direct factorization gives

\[
 a=\frac{3(1+k)^2}{2P_0}.
\tag{C.14.4}\label{contact:one-triple-bernoulli-collision-211-4}
\]

Moreover,

\[
 \frac{g'(b)}3
 =a\{a(1+k^2)-(1+k)\}
 =\frac{a(1+k)(k-1)(k^2+4k+1)}{2P_0}.
\]

The contact derivative is nonpositive, so $0<k\le 1$. The identities

\[
 2P_0-3(1+k)^2=(k-1)^2(2k+1),
\]

\[
 a-\frac34=\frac{3k(1+2k-k^2)}{4P_0}>0
\]

therefore prove

\[
 \boxed{\qquad \frac34<a\leq1,\qquad 0<b\leq1.\qquad}
\tag{C.14.5}\label{contact:one-triple-bernoulli-collision-211-5}
\]

The boundary $a=b=1$ is retained. A selected contact may have zero
derivative at a collision, so it cannot be discarded as a smooth-representation
contradiction.

The remaining triple contact level difference is $1-q=p$. Factoring the
negative cubic at its minimum gives

\[
 \boxed{\quad p=\tau u^2(t+2a-3/2).\quad}
\tag{C.14.6}\label{contact:one-triple-bernoulli-collision-211-6}
\]

Every factor is positive, including $t+2a-3/2>t$ by \eqref{contact:one-triple-bernoulli-collision-211-5}.

\paragraph{Combine the first and second variance derivatives}

Let $\eta$ be the centered signed three-point direction whose variance
derivative is one, as in
Proposition~\ref{prop:finite-contact-exclusion},
The support and threshold stay fixed when the triple law is
changed to $P_T+s \eta$; both signs of sufficiently small $s$ are admissible.
Since the support stays fixed, the collided closed event presents no
differentiability issue.

Let $\gamma=\int |x|^3 \eta(dx)$ and let $L$ be the corresponding CDF
derivative. The levels $(1,q,q)$ and physical gaps $Ku,K(a+b)$ give

\[
 \boxed{\quad L=\frac{p}{K^2u h}>0.\quad}
\tag{C.14.7}\label{contact:one-triple-bernoulli-collision-211-7}
\]

Along this direction the signed excess at price $C$ is exactly

\[
 J(s)=F+sL-\Phi\!\left(\frac z{\sqrt{1+s}}\right)
       -C(R+s\gamma)(1+s)^{-3/2}.
\]

At the ratio maximum, $J'(0)=0$ and $J''(0)\le 0$. Differentiation gives

\[
 C\gamma=L+\frac32CR+\frac12z\phi(z),
\]

\[
 3C\gamma-\frac{15CR+\phi'''(z)}4\leq0.
\]

Using $\phi'''(z)=(3z-z^3)\phi(z)$ and eliminating $\gamma$ yields the exact
necessary condition

\[
 \boxed{\quad 12L+3CR+(z^3+3z)\phi(z)\leq0.\quad}
\tag{C.14.8}\label{contact:one-triple-bernoulli-collision-211-8}
\]

In particular $z=-s<0$, since $L,C,R$ are positive. Equation \eqref{contact:one-triple-bernoulli-collision-211-1} now
reads $CR=CK+s \phi(s)/3$, so \eqref{contact:one-triple-bernoulli-collision-211-8} becomes

\[
 12L+3CK\leq s(s^2+2)\phi(s)\leq2s\phi_0.
\tag{C.14.9}\label{contact:one-triple-bernoulli-collision-211-9}
\]

The last inequality is elementary: the derivative of
$(s^2+2)\phi(s)$ is $-s^3 \phi(s)$, so that function decreases on
$[0,\infty)$ from $2\phi_0$.

\paragraph{Geometry makes the positive CDF derivative too large}

The threshold is the high-triple/low-Bernoulli sum, hence

\[
 s=K(d-b)<Kt.
\tag{C.14.10}\label{contact:one-triple-bernoulli-collision-211-10}
\]

Here $d<h=t+b$ because the Bernoulli law is nondegenerate. Since $L>0$,
\eqref{contact:one-triple-bernoulli-collision-211-9} gives $3CK<2s \phi_0$. Combining this with \eqref{contact:one-triple-bernoulli-collision-211-10} and
$C\ge49/120>\phi_0$ gives

\[
 \boxed{\quad t>\frac sK>\frac{3C}{2\phi_0}>\frac32.\quad}
\tag{C.14.11}\label{contact:one-triple-bernoulli-collision-211-11}
\]

Substitute the exact probability identity \eqref{contact:one-triple-bernoulli-collision-211-6} into \eqref{contact:one-triple-bernoulli-collision-211-7}. Equations
\eqref{contact:one-triple-bernoulli-collision-211-5}, \eqref{contact:one-triple-bernoulli-collision-211-10}, and \eqref{contact:one-triple-bernoulli-collision-211-11} imply

\[
 \begin{aligned}
 L
 &=CK\,\frac{(t-a)(t+2a-3/2)}{t+b}\\
 &>CK\,\frac{t(t-1)}{t+1}\\
 &>Cs\,\frac{t-1}{t+1}\\
 &>\frac{Cs}{5}.
 \end{aligned}
\tag{C.14.12}\label{contact:one-triple-bernoulli-collision-211-12}
\]

The first strict comparison uses $a\le 1$, $b\le 1$, and $2a-3/2>0$;
the final one uses $t>3/2$. Inserting \eqref{contact:one-triple-bernoulli-collision-211-12} into \eqref{contact:one-triple-bernoulli-collision-211-9} yields

\[
 \frac{12Cs}{5}<2s\phi_0,
 \qquad\text{hence}\qquad C<\frac56\phi_0.
\]

The identity \eqref{contact:one-triple-bernoulli-collision-211-8} is also a general necessary condition for any genuine
triple at a full-law ratio maximum, with its actual unit-variance CDF
derivative $L$. In particular $L\ge 0$ forces a negative selected threshold.
Its successful use here depends on the exact collided $211$ geometry and
the actual Bernoulli probability in \eqref{contact:one-triple-bernoulli-collision-211-6}--\eqref{contact:one-triple-bernoulli-collision-211-7}.
\subsection{One triple: collision across the upper gap}\label{contact:collision-221-contact-reduction}


Normalize variance to one and write

\[
 K=R+z\phi(z)/(3C)>0,\qquad \tau=CK^3,\qquad
 \sigma=\phi(z)/(CK^2)>0.
\]

All following support variables are divided by $K$.

\paragraph{The exact compact support and probability charts}

The selected fiber is $(triple-middle, Bernoulli-high)$ and
$(triple-high, Bernoulli-low)$. Write

\[
 T:\ -a,c,b,\qquad B:\ -d,e,
 \qquad a,b,d,e>0,\quad -a<c<b.
\]

The collision forces

\[
 h=b-c=d+e,\qquad p=d/h,\quad q=e/h,
\tag{C.15.1}\label{contact:collision-221-contact-reduction-1}
\]

where $q,p$ are the low and high Bernoulli probabilities. The triple
deletion levels are $(1,1,q)$; its low point is the sole nonselected
contact and is a smooth negative minimum. Thus $a>1/2$ and, up to a
constant, the divided triple dual is

\[
 g(x)=|x|^3-\frac32x^2-3a(1-a)x.
\tag{C.15.2}\label{contact:collision-221-contact-reduction-2}
\]

The middle contact satisfies $g(c)=g(-a)$ and $g'(c)\le0$.
The two equal-height charts in Section~\ref{contact:lower-110-contact-reduction} apply:

\[
 \begin{array}{c|c|c}
 &a&c\\ \hline
 \text{negative},\ 0\leq k\leq1 &1/2+k/4 &(k-1)/2\\
 \text{positive},\ 0\leq k\leq1
   &3(1+k)^2/[2(k^3+3k+2)] &ka.
 \end{array}
\tag{C.15.3}\label{contact:collision-221-contact-reduction-3}
\]

At a genuine configuration $a<1$: the common endpoint $a=c=1$ would
put the still larger contact $b$ to the right of the positive minimum.
The highest triple contact has $g'(b)\le 0$, so

\[
 0<b\leq B(a):=\frac{1+\sqrt{1+4a(1-a)}}2.
\]

Use $b=w B(a)$, $0\le w\le 1$, retaining the unnecessary $b\le c$ region
and rejecting it by the strict gap condition $h>0$.

Let the triple middle mass $m$ be an independent parameter. Its low and
high masses, and the useful contact derivatives, are

\[
 A=a+b,\quad u=a+c,\qquad
 \ell=\frac{b-mh}{A},\qquad r=\frac{a-mu}{A},
\tag{C.15.4}\label{contact:collision-221-contact-reduction-4}
\]

\[
 H_c=-g'(c)=3\{a(1-a)+c-c|c|\},\qquad
 H_b=-g'(b)=3\{a(1-a)+b-b^2\}\geq0.
\]

On the negative chart use the exact identity $H_c=u^2$. Averaging the
dual derivative gives

\[
 \boxed{\quad \sigma=mH_c+rH_b.\quad}
\tag{C.15.5}\label{contact:collision-221-contact-reduction-5}
\]

Thus every candidate lies in one of the two full closed four-cubes

\[
 \boxed{\quad (k,w,p,m)\in[0,1]^4.\quad}
\tag{C.15.6}\label{contact:collision-221-contact-reduction-6}
\]

The denominator $A$ is bounded below by $1/2$ throughout these cubes,
including extraneous boundary points. No small-gap or probability
division is needed in this parameterization.

\paragraph{Actual contact and collision-preserving support equations}

Put

\[
 U=g(-a)-g(b)
   =a^2(3/2-2a)-b^3+\frac32b^2+3a(1-a)b.
\]

The triple contact levels give

\[
 \boxed{\quad p=\tau U,\qquad U>0.\quad}
\tag{C.15.7}\label{contact:collision-221-contact-reduction-7}
\]

For the Bernoulli law, set $k_B=h(p^2+q^2)$. Its divided derivative at
the high support point is

\[
 f'_B(Ke)/(CK^2)=3e(k_B-1)-\sigma.
\]

\[
 m f'_T(Kc)=p f'_B(Ke).
\]

Substitute \eqref{contact:collision-221-contact-reduction-5} and the displayed derivative. The exact necessary
collision equation is

\[
 \boxed{\quad q mH_c-p rH_b-3pq h(1-k_B)=0.\quad}
\tag{C.15.8}\label{contact:collision-221-contact-reduction-8}
\]

This support motion preserves the collision and uses the actual marginal masses.

The Bernoulli deletion levels are $(1,\ell+m)=(1,1-r)$. Its signed second
moment is $M_B=de(e-d)/h$, and its divided derivative everywhere is
$3[x|x|-x-M_B]-\sigma$. Integrating between $-d,e$ gives its exact value
difference

\[
 f_B(Ke)-f_B(-Kd)
 =\tau\{J-\sigma h\}=-r,
\]

\[
 J=(e-d)\{(e-d)^2-3h/2\}.
\]

Using \eqref{contact:collision-221-contact-reduction-7} clears the denominator:

\[
 \boxed{\quad p(J-\sigma h)+rU=0.\quad}
\tag{C.15.9}\label{contact:collision-221-contact-reduction-9}
\]

\paragraph{Variance, normal comparisons, and exact bounds}

Use the following dimensionless moments:

\[
 V_T=ab-muh,\qquad V_B=de,\qquad V_0=V_T+V_B,
\]

\[
 R_T=\ell a^3+m|c|^3+r b^3,\qquad R_B=de\,k_B,
 \qquad R_0=R_T+R_B.
\]

Normalization and \eqref{contact:collision-221-contact-reduction-7} give

\[
 K=V_0^{-1/2},\qquad z=\frac{o}{\sqrt{V_0}},\quad o=c+e,
 \qquad C=\frac{pV_0^{3/2}}U.
\tag{C.15.10}\label{contact:collision-221-contact-reduction-10}
\]

The definition of $K$ and \eqref{contact:collision-221-contact-reduction-5} imply the further necessary polynomial
identity

\[
 \boxed{\quad 3(V_0-R_0)-\sigma o=0.\quad}
\tag{C.15.11}\label{contact:collision-221-contact-reduction-11}
\]

The bounds $1/2<K<8/5$ in \eqref{eq:contact-basic-collars}
also give the strict necessary bound

\[
 \boxed{\quad 25/64<V_0<4.\quad}
\tag{C.15.11a}\label{contact:collision-221-contact-reduction-11a}
\]

This holds before any contact division and removes zero-variance limits.

The strict ratio interval gives

\[
 (49/120)^2U^2\le p^2V_0^3<(227/500)^2U^2.
\tag{C.15.12}\label{contact:collision-221-contact-reduction-12}
\]

These denominator-cleared tests are available even when a box's variance
enclosure crosses zero. On boxes with a strictly positive variance
enclosure, also impose the unsquared density and actual CDF equations

\[
 \boxed{\quad \sigma p\sqrt{V_0}-U\phi(z)=0,\quad}
\tag{C.15.13}\label{contact:collision-221-contact-reduction-13}
\]

\[
 \boxed{\quad (1-pr)U-pR_0-U\Phi(z)=0.\quad}
\tag{C.15.14}\label{contact:collision-221-contact-reduction-14}
\]

The true event probability is $F=1-pr$. A simple necessary CDF positivity
test without transcendental evaluation follows by averaging the triple
dual, which equals one at $-a$:

\[
 \boxed{\quad U-p\{a^2(3/2-2a)+(3/2)V_T+R_B\}>0.\quad}
\tag{C.15.15}\label{contact:collision-221-contact-reduction-15}
\]

At the maximizing pair, each marginal variance lies strictly between
$13/200$ and $187/200$ of the total variance, by
\eqref{eq:contact-variance-collar}. Hence

\[
 \boxed{\quad V_T>\frac{13}{200}V_0,
                 \qquad V_B>\frac{13}{200}V_0.\quad}
\tag{C.15.16}\label{contact:collision-221-contact-reduction-16}
\]

Finally, the triple's actual unit-variance CDF derivative is
$L=-p/(K^2 hA)$. The first/second variance identity \eqref{contact:one-triple-bernoulli-collision-211-8} of
Section~\ref{contact:one-triple-bernoulli-collision-211}
applies without changing the closed event. Substitute \eqref{contact:collision-221-contact-reduction-10}, \eqref{contact:collision-221-contact-reduction-13}, and
multiply by the positive quantity $U h A V_0/p$. The necessary
second-order condition becomes the polynomial inequality

\[
 \boxed{\quad
 -12V_0^2U+3R_0hAV_0+\sigma hAo(o^2+3V_0)\leq0.
 \quad}
\tag{C.15.17}\label{contact:collision-221-contact-reduction-17}
\]
The fixed certificate \texttt{collision221} proves that the two closed
support charts above contain no point satisfying their contact equations,
probability conditions, and necessary inequalities. Its strict-floor
checks include $C=49/120$, as described in Section~\ref{sec:verification}.
Thus a collision across the upper gap cannot occur at a maximizing pair.

\subsection{One triple: collision across the lower gap}\label{contact:collision-210-contact-reduction}

Use the contact scale: the support coordinates below are divided by $K$.


\[
 K=R+z\phi(z)/(3C),\quad \tau=CK^3,
 \quad \sigma=\phi(z)/(CK^2).
\]

The bounds in \eqref{eq:contact-basic-collars} give

\[
 \boxed{\quad 1/2<K<8/5,\qquad 1/20<\tau<2.\quad}
\tag{C.16.1}\label{contact:collision-210-contact-reduction-1}
\]

The lower bound for $\tau$ uses $C>2/5$; the upper bound uses
$C<227/500$ and $(227/500)(8/5)^3<2$.

\paragraph{Compact support charts from the zero minimum}

Divide all support points by $K$ and write

\[
 T:\ -t,c,b,\qquad B:\ -d,e,
 \quad t,b,d,e>0,\quad -t<c<b.
\]

The selected fiber is $(triple-low, Bernoulli-high)$ and
$(triple-middle, Bernoulli-low)$, so

\[
 h=t+c=d+e,\qquad p=d/h,\quad q=e/h.
\tag{C.16.2}\label{contact:collision-210-contact-reduction-2}
\]

The low/high Bernoulli probabilities are $q,p$; the triple deletion
levels are $(1,q,0)$. Its high contact is nonselected and is a smooth
positive minimum of height zero. The whole dual is nonnegative, so
$b\ge 1$. Indeed a positive minimum with $1/2<b<1$ has a strictly lower
negative local minimum, contradicting the nonnegative global majorant.

The triple dual divided by $\tau$ is exactly

\[
 g_b(x)=|x|^3-\frac32x^2-3b(b-1)x+B_0,
 \qquad B_0=b^2(2b-3/2).
\tag{C.16.3}\label{contact:collision-210-contact-reduction-3}
\]

It vanishes at $b$. Put

\[
 G=g_b(-t)=t^3-\frac32t^2+3b(b-1)t+B_0,
 \qquad G_c=g_b(c).
\]

The contact values give

\[
 \boxed{\quad \tau=1/G,\quad q=G_c/G,\quad p=1-q,
                  \quad 1/2<G<20.\quad}
\tag{C.16.4}\label{contact:collision-210-contact-reduction-4}
\]

The height bounds in \eqref{contact:collision-210-contact-reduction-4} follow from \eqref{contact:collision-210-contact-reduction-1}. They also bound the supports
without using any infinite-profile radius theorem. If $b\ge 5/2$, then
$g_b(-t)$ increases for $t\ge 0$, since its derivative is
$3[(t-1/2)^2+b(b-1)-1/4]>0$. Hence
$G\ge B_0\ge 175/8>20$, impossible. Since $b\ge 1$,

\[
 G\geq t^3-\frac32t^2+\frac12.
\]

For $t\ge 7/2$, the right side is increasing and is at least 25,
also impossible. Thus

\[
 1\leq b<5/2,\qquad 0<t<7/2.
\tag{C.16.5}\label{contact:collision-210-contact-reduction-5}
\]

Use independent parameters $(s_b,s_t,w,m) \in [0,1]^4$ and maps

\[
 b=1+\frac32s_b,\qquad t=\frac72s_t,
\]

\[
 c=-t(1-w)\quad\text{on the negative chart},\qquad
 c=wb\quad\text{on the positive chart}.
\tag{C.16.6}\label{contact:collision-210-contact-reduction-6}
\]

Here $m$ is the triple middle mass. Both closed cubes are retained,
including all extra boundary points and the shared $c=0$ face.

Let $A=t+b$, $v=b-c$. Centering determines

\[
 \ell=\frac{b-mv}{A},\qquad r=\frac{t-mh}{A}.
\tag{C.16.7}\label{contact:collision-210-contact-reduction-7}
\]

The denominator $A\ge 1$ throughout the whole parameter domain.

\paragraph{Exact collision and Bernoulli contact equations}

Define the nonnegative selected-contact slopes

\[
 H_\ell=3\{t^2-t+b(b-1)\},\qquad
 H_m=3\{b(b-1)+c-c|c|\}.
\]

Averaging the triple dual derivative gives

\[
 \boxed{\quad \sigma=\ell H_\ell+mH_m>0.\quad}
\tag{C.16.8}\label{contact:collision-210-contact-reduction-8}
\]

Set $d=ph$, $e=qh$, and $k_B=h(p^2+q^2)$. Move the raw low triple
point and the raw high Bernoulli point by opposite amounts, holding
all other raw support points and all probabilities fixed. Both selected
raw sums and the closed event stay fixed. Recenter the same two
marginals and move the threshold by their total mean shift. The
two-sided ratio derivative gives

\[
 \ell f'_T(-Kt)=p f'_B(Ke).
\]

Using the Bernoulli derivative
$f'_B(Ke)/(CK^2)=3e(k_B-1)-\sigma$ and \eqref{contact:collision-210-contact-reduction-8}, this is

\[
 \boxed{\quad q\ell H_\ell-pmH_m-3pq h(1-k_B)=0.\quad}
\tag{C.16.9}\label{contact:collision-210-contact-reduction-9}
\]

The Bernoulli deletion levels are $(\ell+m,\ell)$. Integrating its divided
derivative, as in the collided $221$ reduction, gives

\[
 \tau\{J-\sigma h\}=-m,\qquad
 J=(e-d)\{(e-d)^2-3h/2\}.
\]

Since $\tau=1/G$, its denominator-cleared contact equation is

\[
 \boxed{\quad J-\sigma h+mG=0.\quad}
\tag{C.16.10}\label{contact:collision-210-contact-reduction-10}
\]

\paragraph{Normalization, actual CDF, and second-order constraints}

Write the dimensionless variance and sum of third absolute moments as

\[
 V_T=tb-mhv,\quad V_B=de,\quad V_0=V_T+V_B,
\]

\[
 R_T=\ell t^3+m|c|^3+r b^3,\quad R_B=de\,k_B,
 \quad R_0=R_T+R_B.
\]

Then

\[
 K=V_0^{-1/2},\qquad C=V_0^{3/2}/G,\qquad
 o=c-d,\qquad z=o/\sqrt{V_0}.
\]

The definition of $K$ gives the necessary polynomial identity

\[
 \boxed{\quad 3(V_0-R_0)-\sigma o=0.\quad}
\tag{C.16.11}\label{contact:collision-210-contact-reduction-11}
\]

The compact bounds and ratio interval imply

\[
 25/64<V_0<4,\qquad
 (49/120)^2G^2\le V_0^3<(227/500)^2G^2.
\tag{C.16.12}\label{contact:collision-210-contact-reduction-12}
\]

At the maximizing pair, each marginal variance lies strictly between
$13/200$ and $187/200$ of the total variance, by
\eqref{eq:contact-variance-collar}. Hence

\[
 V_T>\frac{13}{200}V_0,\qquad
 V_B>\frac{13}{200}V_0.
\tag{C.16.13}\label{contact:collision-210-contact-reduction-13}
\]

The actual event probability is $F=\ell+mq$. The density and CDF equations
are

\[
 \boxed{\quad \sigma\sqrt{V_0}-G\phi(z)=0,\quad}
\tag{C.16.14}\label{contact:collision-210-contact-reduction-14}
\]

\[
 \boxed{\quad G\{\ell+mq-\Phi(z)\}-R_0=0.\quad}
\tag{C.16.15}\label{contact:collision-210-contact-reduction-15}
\]

Averaging the zero-minimum triple dual \eqref{contact:collision-210-contact-reduction-3} also supplies the necessary
positivity test

\[
 \boxed{\quad B_0-\frac32V_T-R_B>0.\quad}
\tag{C.16.16}\label{contact:collision-210-contact-reduction-16}
\]

For the triple's centered unit-variance direction, the actual CDF
derivative at the fixed closed event is

\[
 L=\frac{(pA-h)V_0}{Ahv}.
\]

Combine its first and second variance derivatives and use \eqref{contact:collision-210-contact-reduction-14}. After
multiplication by the positive factor $G A h v V_0$, the necessary
second-order condition is

\[
 \boxed{\quad
 12(pA-h)GV_0^2+3R_0AhvV_0
       +\sigma Ahv\,o(o^2+3V_0)\leq0.
 \quad}
\tag{C.16.17}\label{contact:collision-210-contact-reduction-17}
\]

All probabilities in these formulas are the actual marginal probabilities;
they have not been relaxed to independent deletion-CDF parameters.
\subsection{Lower-gap collision: elimination of the middle mass}\label{contact:collision-210-mass-elimination}
Continue with the supports divided by $K$ in Section~\ref{contact:collision-210-contact-reduction}.



\paragraph{A positive affine coefficient}

Retain the previous notation

\[
 A=t+b,\quad h=t+c,\quad v=b-c,\quad
 \ell=(b-mv)/A,\quad r=(t-mh)/A,
\]

with $H_t,H_m\ge 0$, $\sigma=\ell H_t+m H_m>0$, and Bernoulli probabilities
$q,p>0$. The collision support equation is

\[
 q\ell H_t-pmH_m-3pq h\{1-h(p^2+q^2)\}=0.
\]

Substitution of $\ell$ makes it affine in $m$, and gives exactly

\[
 \boxed{
 m=\frac{q b H_t-3A p q h\{1-h(p^2+q^2)\}}
           {qvH_t+A pH_m}.}
\tag{C.17.1}\label{contact:collision-210-mass-elimination-1}
\]

The denominator $L=qvH_t+A pH_m$ is strictly positive: every coefficient
is positive, both slopes are nonnegative, and $\sigma>0$ excludes their
simultaneous vanishing. Thus there is no additional zero-denominator
collision case to retain.

Let $N=G_c$, $P=G-G_c$, $W=P^2+N^2$, so $q=N/G$, $p=P/G$. Define

\[
 \Delta=G^3(NvH_t+A PH_m)=G^4 L,
\]
\[
 M=N bH_tG^3-3A PN hG^2+3A PN h^2W.
\tag{C.17.2}\label{contact:collision-210-mass-elimination-2}
\]

Then $m=M/\Delta$. The four natural mass conditions become polynomial:

\[
 M>0,\quad\Delta-M>0,\quad
 b\Delta-vM>0,\quad t\Delta-hM>0.
\tag{C.17.3}\label{contact:collision-210-mass-elimination-3}
\]

\paragraph{A uniform positive denominator bound}

The canonical mixture-width theorem applies to this triple and its
Bernoulli partner without an optimality assumption. Consequently every
violating candidate in this system must satisfy

\[
 \min\{bh,tv\}>V/3>25/192.
\tag{C.17.4}\label{contact:collision-210-mass-elimination-4}
\]

The Bernoulli variance bound gives

\[
 pq h^2>\frac{13}{200}V>\frac{13}{512}.
\]

As $h<A=t+b<6$, this yields
$p,q>pq>13/18432$. Moreover $A\ge 1$.

There is also a uniform selected-slope bound. With normalized variance
one, $D=C R_{\mathrm{normalized}}>49/(120 \sqrt{2})>1/5$. Cantelli's inequality gives
$|z|<2$: for $z\le -2$, the actual CDF is at most $1/5$; for $z\ge 2$, the
normal upper tail is at most $1/5$. Either would contradict $D>1/5$.
The elementary estimates $\pi<22/7$ and $e<11/4$ imply

\[
 2\pi e^4<\frac{44}{7}\left(\frac{11}{4}\right)^4<400,
 \qquad\phi(2)>1/20.
\]

Therefore the actual density equation, $\sigma \sqrt{V}=G \phi(z)$, and
$G>1/2$, $\sqrt{V}<2$ give $\sigma>1/80$.
Since $\ell,m<1$, one has $H_t+H_m\ge \sigma$, and hence

\[
 L>\frac{13}{18432}\frac{25}{672}\frac1{80}.
\]

Multiplication by $G^4>1/16$ proves the explicit bound

\[
 \boxed{\Delta>\frac{65}{3170893824}>\frac1{50000000}.}
\tag{C.17.5}\label{contact:collision-210-mass-elimination-5}
\]

The final comparison has exact numerator slack
$65\,50000000-3170893824=79106176$. Thus the eliminated mass is a smooth
rational function on an open neighborhood of every candidate in the
compact support charts. In interval arithmetic, \eqref{contact:collision-210-mass-elimination-5} is only a necessary
test; division is permitted only if the direct enclosure of $\Delta$ is
positive on the whole current box.

\paragraph{The two three-parameter support charts}

Delete the independent mass variable from the two charts. For
$(s_b,s_t,w) \in [0,1]^3$, set

\[
 b=1+3s_b/2,\qquad t=7s_t/2,
\]

and take $c=-t(1-w)$ on the negative chart or $c=wb$ on the positive chart.
Compute $m$ by \eqref{contact:collision-210-mass-elimination-1} or \eqref{contact:collision-210-mass-elimination-2} and retain all genuine-support, probability,
height, slope, ratio, variance, width and global-dual constraints. Every
original candidate has a preimage in one of these two closed three-cubes;
the shared $c=0$ face remains covered. Conversely, any candidate satisfying
the retained tests and four remaining equations satisfies the eliminated
support equation identically. This statement only identifies necessary
systems; it does not assert that every solution is a maximizing law.

\paragraph{Four explicit equations without an independent mass variable}

The following formulas permit denominator-free evaluation before any
normal function. Define

\[
 L_*=b\Delta-vM,\quad H_*=t\Delta-hM,
 \quad S_*=L_*H_t+A MH_m,
 \quad O=cN-tP,
\]
\[
 V_*=(tb\Delta-Mhv)G^2+\Delta PN h^2,
\]
\[
 R_*=(L_*t^3+A M|c|^3+H_*b^3)G^4
                 +A\Delta PN h^3W,
\quad F_*=GL_*+A MN.
\]

Then the actual quantities are

\[
 V=\frac{V_*}{\Delta G^2},\quad
 R=\frac{R_*}{A\Delta G^4},\quad
 \sigma=\frac{S_*}{A\Delta},\quad
 F=\frac{F_*}{A\Delta G},\quad
 z=O\sqrt{\Delta/V_*}.
\tag{C.17.6}\label{contact:collision-210-mass-elimination-6}
\]

The retained Bernoulli height, scale, squared density and CDF equations
are, respectively,

\[
 A\Delta h(N-P)\{h^2(N-P)^2-3hG^2/2\}
                  -S_*hG^3+A MG^4=0,
\]
\[
 3(AG^2V_*-R_*)-S_*OG^3=0,
\]
\[
 S_*^2V_*-A^2\Delta^3G^4\phi(z)^2=0,
\]
\[
 G^4\{F_*-A\Delta G\Phi(z)\}-R_*=0.
\tag{C.17.7}\label{contact:collision-210-mass-elimination-7}
\]

Each clearing factor is positive at every candidate. In particular the
two equations preceding the normal functions are exact polynomial
equations in the three support parameters on either sign chart. No
resultant or Groebner calculation is needed for this dimension reduction.

The ratio test is
$(49/120)^2 \Delta^3 G^8\le V_*^3<(227/500)^2 \Delta^3 G^8$.
The added width tests are
$3bh \Delta G^2>V_*$ and $3tv \Delta G^2>V_*$.
\subsection{Lower-gap collision: direct necessary equations}\label{contact:collision-210-raw-law-replay}
Use the contact scale of Section~\ref{contact:collision-210-contact-reduction}.



\paragraph{Independent law coordinates}

Write the centered triple supports $(-t,c,b)$, set $A=t+b$, $h=t+c$,
$v=b-c$, and use its variance $W=V_T$ as the mass coordinate. Exact
centering gives

\[
 \ell={W+bc\over Ah},\qquad m={tb-W\over hv},
 \qquad r={W-tc\over Av}. \tag{C.18.1}\label{contact:collision-210-raw-law-replay-1}
\]

These formulas are intersected with normalization, the direct second
moment, and the usual middle-mass expression $W=tb-mhv$.
For the Bernoulli, low/high probabilities are $q,p$, supports are
$(-d,e)$, and $d=ph$, $e=qh$. Its third absolute moment is evaluated as

\[
 V_B=de,\qquad R_B=V_B\left(h-{2V_B\over h}\right),
 \qquad h/2\le h-2V_B/h\le h. \tag{C.18.2}\label{contact:collision-210-raw-law-replay-2}
\]

\[
 R_*={tb(t^2+b^2)\over A},\qquad M_*={tb(b-t)\over A},
\]

the exact moments are

\[
 R_T=R_*+\gamma(W-tb),\qquad
 \E[T|T|]=M_*+\eta(W-tb). \tag{C.18.3}\label{contact:collision-210-raw-law-replay-3}
\]

The direct three-term sum of third absolute moments is also intersected. The bounds
$A/2\le \gamma\le 2A$ and $-1\le \eta\le 1$ are the established divided-
difference bounds. All adjacent and outer gap divisions are performed
separately with their positive candidate bounds; a product enclosure that
straddles zero is never inverted.

\paragraph{Dual contacts and support pairings}

The triple majorant is evaluated using the factorization

\[
 g_b(x)=(x-b)^2(x+2b-3/2)+2(-x)_+^3. \tag{C.18.4}\label{contact:collision-210-raw-law-replay-4}
\]

Thus $G=g_b(-t)$, $q=g_b(c)/G$, and $p=1-q$. A differently arranged
expanded value of $G$ is also intersected. The high minimum gives

\[
 \sigma=3\{b(b-1)-\E[T|T|]\}=\ell H_t+mH_m,
\]

where $H_t=3(t^2-t+b(b-1))$ and
$H_m=3(b(b-1)+c-c|c|)$ are nonnegative. Put

\[
 Q=3pq h\{1-h(p^2+q^2)\}.
\]

The two collision-preserving derivative pairings are evaluated as

\[
 \ell H_t=p\sigma+Q,\qquad mH_m=q\sigma-Q. \tag{C.18.5}\label{contact:collision-210-raw-law-replay-5}
\]

The Bernoulli height equation is

\[
 mG=\sigma h-(e-d)\{(e-d)^2-3h/2\}. \tag{C.18.6}\label{contact:collision-210-raw-law-replay-6}
\]

\paragraph{Normalization and inverse identities}

The independently stored values satisfy the contact equations

\[
 C=V^{3/2}/G,\quad D=R/G=F-\Phi(z),\quad
 \sigma\sqrt V=G\phi(z),\quad 3(V-R)=\sigma o,
 \quad o=c-d=e-t=z\sqrt V. \tag{C.18.7}\label{contact:collision-210-raw-law-replay-7}
\]

They are evaluated both forward and backward, including
$V=(CG)^{2/3}$, $\sqrt{V}=G \phi/\sigma$, and the exact inverse density

\[
 z^2=-\log(2\pi\phi(z)^2). \tag{C.18.8}\label{contact:collision-210-raw-law-replay-8}
\]

Equation \eqref{contact:collision-210-raw-law-replay-8} has a known negative square-root branch on the negative chart. On the positive chart a branch is used only when the
current candidate interval has proved its sign; otherwise its symmetric
hull is retained. A nonnegative interval for $z^2$ is intersected before
taking a square root. The CDF is evaluated through the error function.

The following bounds hold at every candidate:
$49/120\le C\le 227/500$, $49/(120 \sqrt{2})\le D\le 11/20$,
$-5/4\le z\le 3/5$, $\phi\ge 9/50$,
$V\ge 2401/5476$ (or $2401/4356$ on the negative chart), and $V\le 4$.
It also uses the established coordinate variance fraction bounds
$13V/200\le V_i\le 187V/200$ and H\"older third-moment bounds. References are
Section~\ref{contact:fixed-c2-threshold-collar},
Section~\ref{contact:fixed-c2-contact-scale-collar}, and
Section~\ref{contact:collision-210-mass-elimination}.
Strict bounds are relaxed to closed bounds for enclosure.

The canonical mixture-width inequalities $3bh>V$, $3tv>V$ are also used
backward to obtain necessary bounds on each positive factor. The elementary
Bernoulli bound $h\ge 2 \sqrt{V_B}$ is included. Initial gap and probability
floors are exactly the proved floors: $h\ge \sqrt{13/128}$,
$v\ge 25/672$, $p,q\ge 13/18432$, and $t\ge 13/1280$.
The fixed certificate \texttt{collision210} proves that no point of
either closed support chart satisfies the necessary conditions in
Sections~\ref{contact:collision-210-contact-reduction}--\ref{contact:collision-210-raw-law-replay}.
The enclosure and partition rules are given in Section~\ref{sec:verification}.
This excludes the lower-gap collision.

\subsection{Two triples: asymmetric events with a unique representation}\label{contact:two-triple-asymmetric-exclusion}



Normalize total variance to one and use the dimensionless coordinates of
Section~\ref{contact:two-triple-rectangle-exclusion}. Thus the two coordinate supports are $K(-a_i,c_i,b_i)$, with

\[
 \frac12<a_i<1<b_i,\qquad a_i^2+b_i^2=a_i+b_i,
 \qquad h_i=a_i+b_i,
\]

and the two nonselected minima have height difference

\[
 C K^3D(a_i),\qquad
 D(a)=\frac12(b-a)(a^2+4ab+b^2).
\]

The decreasing deletion CDF gives these same strict inequalities for
$a_i,b_i$ in $221$ and $320$, because the two nonselected contact levels
are distinct. Along the curve $a^2+b^2=a+b$,

\[
 D'(a)=-3(2a-1)(a+b)<0,
 \qquad
 h'(a)=\frac{2(b-a)}{2b-1}>0.
\tag{C.19.1}\label{contact:two-triple-asymmetric-exclusion-1}
\]

For shape $221$, the first coordinate's three deletion-CDF values are
$q_2,q_2,l_2$, where $l_2$ is the lowest mass, $m_2$ the middle mass,
and $q_2=l_2+m_2$. The second coordinate's three values are
$1,q_1,0$. Therefore

\[
 C K^3D(a_1)=m_2<1=C K^3D(a_2).
\]

Equation \eqref{contact:two-triple-asymmetric-exclusion-1} forces $a_1>a_2$ and hence $h_1>h_2$.

On the other hand, the selected representation is middle-middle. The high-low
corner is included, and the low-high corner is excluded. The unique threshold representation
makes both relevant comparisons strict:

\[
 b_1-a_2<c_1+c_2<-a_1+b_2.
\]

With $u_i=a_i+c_i$, these are exactly

\[
 h_1<u_1+u_2<h_2,
\]

contradicting $h_1>h_2$. Interchanging the two coordinates transposes
$221$ to $320$, so that shape is excluded as well.
\subsection{Two triples: triangular event with a unique representation}\label{contact:two-triple-triangle-exclusion}


\paragraph{Contact geometry and notation}

Normalize total variance to one. Write the discrepancy as $D=C R$, the
selected threshold as $z$, and put

\[
 K=\frac{2A}{3C}=R+\frac{z\phi(z)}{3C}.
\]

The full-law cubic contacts from
Proposition~\ref{prop:finite-contact-exclusion}
give supports $K(-a_i,c_i,b_i)$, with the selected point in the middle
of each coordinate. The two nonselected points are smooth minima.
In shape $321$ their deletion-CDF values are $1$ and the other
coordinate's lowest mass, so the minimum-height difference is positive.
Consequently

\[
 \frac12<a_i<1<b_i,\qquad
 a_i^2+b_i^2=a_i+b_i=:h_i\leq2.
\tag{C.20.1}\label{contact:two-triple-triangle-exclusion-1}
\]

Define $u_i=a_i+c_i$ and $v_i=b_i-c_i$. The two cross corners are
strictly below the selected middle-middle sum. Thus

\[
 v_1<u_2,\qquad v_2<u_1,
 \qquad
 c_1+c_2>\max\{b_1-a_2,b_2-a_1\}>0.
\tag{C.20.2}\label{contact:two-triple-triangle-exclusion-2}
\]

In particular $z=K(c_1+c_2)>0$ and $K>0$. Let $l_i,m_i,r_i$
denote the lowest, middle, and highest masses, and put $p_i=m_i+r_i$.
The actual selected event probability and its complement are

\[
 F=1-p_1p_2+m_1m_2,\qquad
 P_{\rm out}=r_1p_2+r_2p_1-r_1r_2.
\tag{C.20.3}\label{contact:two-triple-triangle-exclusion-3}
\]

\paragraph{A divided-difference inequality on the actual triangle}

For a pair $a,b$ satisfying \eqref{contact:two-triple-triangle-exclusion-1}, let

\[
 d_a(c)=[-a,c,b]|x|^3.
\]

The divided-difference convention is $[x_0,x_1,x_2]x^2=1$.
Direct interpolation gives

\[
 d_a(c)=
 \begin{cases}
 b+c-a+\dfrac{2a^3}{(a+b)(a+c)},&c\geq0,\\[4pt]
 a-b-c+\dfrac{2b^3}{(a+b)(b-c)},&c\leq0.
 \end{cases}
\tag{C.20.4}\label{contact:two-triple-triangle-exclusion-4}
\]

This is a strictly convex function of $c$ on $(-a,b)$, with

\[
 d_a(0)=1,\qquad d_a'(0)=\frac{b-a}{a+b}>0.
\tag{C.20.5}\label{contact:two-triple-triangle-exclusion-5}
\]

We prove that \eqref{contact:two-triple-triangle-exclusion-2} implies

\[
 \boxed{\quad d_{a_1}(c_1)+d_{a_2}(c_2)>2.\quad}
\tag{C.20.6}\label{contact:two-triple-triangle-exclusion-6}
\]

Relabel so $a_1\le a_2$. Then $b_1\ge b_2$, $h_1\le h_2$, and the
slopes $k_i=(b_i-a_i)/h_i$ satisfy $k_1\ge k_2$. Put
$z_0=b_2-a_1>0$. We need one derivative comparison:

\[
 d_{a_2}'(z_0)\geq k_1.
\tag{C.20.7}\label{contact:two-triple-triangle-exclusion-7}
\]

By \eqref{contact:two-triple-triangle-exclusion-4}, this is equivalent to

\[
 a_2^3h_1\leq a_1h_2(h_2-a_1)^2.
\]

Use $h_1\le 2$ and $h_2\ge a_2+1$. With $a=a_1$, $b=a_2$, it is
sufficient that

\[
 a(b+1)(b+1-a)^2\geq2b^3,
 \qquad 1/2\leq a\leq b\leq1.
\tag{C.20.8}\label{contact:two-triple-triangle-exclusion-8}
\]

For fixed $b$, the function $a(b+1-a)^2$ has at most one interior
critical point, a maximum. Its minimum on $[1/2,b]$ is at an endpoint.
At $a=b$, \eqref{contact:two-triple-triangle-exclusion-8} follows from $b+1\ge 2b^2$. At $a=1/2$, it follows from

\[
 (b+1)(b+1/2)^2-4b^3
 =\frac12+(1-b)(3b^2+b-1/4)>0.
\]

This proves \eqref{contact:two-triple-triangle-exclusion-7}.

If both middle points are nonnegative, \eqref{contact:two-triple-triangle-exclusion-6} follows immediately from
\eqref{contact:two-triple-triangle-exclusion-5}, since their sum is positive. If $c_1<0$, then \eqref{contact:two-triple-triangle-exclusion-2} gives
$c_2>z_0-c_1$. Along $s$ from $c_1$ to zero, strict convexity and
\eqref{contact:two-triple-triangle-exclusion-7} give

\[
 \frac{d}{ds}\{d_{a_1}(s)+d_{a_2}(z_0-s)\}
 \leq k_1-d_{a_2}'(z_0)\leq0.
\]

The sum is therefore at least $1+d_{a_2}(z_0)>2$. The replacement
argument stays within the middle-knot intervals because
$z_0-c_1<c_2<b_2$. If $c_2<0$, the same argument uses
$d_{a_1}'(z_0)\ge k_1\ge k_2$. This proves \eqref{contact:two-triple-triangle-exclusion-6}, including mixed signs.

\paragraph{The actual probability Hessian forces K>1}

Use the centered unit-variance signed probability directions on the two
three-point supports displayed in Proposition~\ref{prop:finite-contact-exclusion}. Their third-moment
derivatives are

\[
 \gamma_i=K d_{a_i}(c_i).
\]

For their two perturbation parameters $s,t$, the total variance is
$1+s+t$, the sum of third absolute moments is $R+\gamma_1 s+\gamma_2 t$, and the selected CDF
is $F+f_1s+f_2t+g st$. Both signs of each sufficiently small parameter
are admissible. At the maximizing ratio, differentiate

\[
 J(s,t)=(1+s+t)^{3/2}
 \{F(s,t)-\Phi(z/\sqrt{1+s+t})\}
 -C(R+\gamma_1s+\gamma_2t).
\]

It has a local maximum zero. First stationarity and direct second
differentiation give its Hessian entries

\[
 J_{ii}=3C\gamma_i-B_*=-b_i^*,\qquad
 J_{12}=g-\frac{b_1^*+b_2^*}{2},
\]

\[
 B_*=(15CR+\phi'''(z))/4,\qquad b_i^*=B_*-3C\gamma_i\geq0.
\tag{C.20.9}\label{contact:two-triple-triangle-exclusion-9}
\]

This is the two-coordinate extension of
Section~\ref{contact:three-point-variance-curvature}.
Negative semidefiniteness implies

\[
 g\leq\frac{b_1^*+b_2^*}{2}+\sqrt{b_1^*b_2^*}
 \leq b_1^*+b_2^*.
\tag{C.20.10}\label{contact:two-triple-triangle-exclusion-10}
\]

The actual $321$ grid gives, exactly,

\[
 g=\frac{1-v_1v_2/(h_1h_2)}{K^4u_1v_1u_2v_2}
   >\frac3{4K^4}.
\tag{C.20.11}\label{contact:two-triple-triangle-exclusion-11}
\]

Indeed \eqref{contact:two-triple-triangle-exclusion-2} gives $h_1h_2>(v_1+v_2)^2\ge 4v_1v_2$, and
$u_iv_i\le h_i^2/4\le 1$.

Since $CR=CK-z \phi(z)/3$ and $\phi'''(z)=(3z-z^3)\phi(z)$, equations
\eqref{contact:two-triple-triangle-exclusion-6}, \eqref{contact:two-triple-triangle-exclusion-9}--\eqref{contact:two-triple-triangle-exclusion-11} yield

\[
 \frac3{4K^4}
 <\frac32CK-\frac12z(z^2+2)\phi(z).
\tag{C.20.12}\label{contact:two-triple-triangle-exclusion-12}
\]

Put $N(z)=\Phi(z)-z \phi(z)/3$. It is strictly increasing and exceeds
$1/2$ for positive $z$. The ratio identity is $CK=F-N(z)$, so $CK<1/2$.
Equation \eqref{contact:two-triple-triangle-exclusion-12} now implies

\[
 \boxed{\quad K>1,\qquad K<1/(2C).\quad}
\tag{C.20.13}\label{contact:two-triple-triangle-exclusion-13}
\]

\paragraph{The Hessian also forces z<1/4}

Let

\[
 L(z)=1-N(z)=Q(z)+z\phi(z)/3,
 \qquad
 A_*(z)=\frac32Q(z)-\frac12z(z^2+1)\phi(z).
\]

We have $CK<L(z)$. Equation \eqref{contact:two-triple-triangle-exclusion-12}, and $C\ge49/120$, imply

\[
 \frac34<K^4A_*(z)
       <\left(\frac{L(z)}{49/120}\right)^4A_*(z).
\tag{C.20.14}\label{contact:two-triple-triangle-exclusion-14}
\]

The middle expression must be positive, so \eqref{contact:two-triple-triangle-exclusion-14} is only used when
$A_*(z)>0$. First, $z<1$: the moment inequality $R\ge 1/\sqrt{2}$ gives
$CR>1/4$, whereas $Q(1)<\phi(1)<1/4$. A threshold $z\ge 1$ would thus
contradict $CR=F-\Phi(z)\le Q(z)$.

On $[0,1]$, both $L$ and $A_*$ strictly decrease, since

\[
 L'(z)=-(2+z^2)\phi(z)/3,
 \qquad
 A_*'(z)=-(2+z^2-z^4/2)\phi(z).
\]

At $z=1/4$, use $Q(z)\le 1/2-z \phi(z)$ and $\phi(1/4)>3/8$ to get

\[
 L(1/4)<7/16,
 \qquad A_*(1/4)<573/1024.
\]

If $z\ge 1/4$, \eqref{contact:two-triple-triangle-exclusion-14} would consequently give

\[
 \frac34<\left(\frac{15}{14}\right)^4\frac{573}{1024}<\frac34,
\]

a contradiction. The exact slack in the last comparison is
$495363/39337984$. We have proved

\[
 \boxed{\quad0<z<1/4.\quad}
\tag{C.20.15}\label{contact:two-triple-triangle-exclusion-15}
\]

The Gaussian comparisons are elementary: $\phi(0)>39/100$ and
$\exp(-1/32)>31/32$ give $\phi(1/4)>3/8$; also $\pi e>8$ gives
$\phi(1)<1/4$.

\paragraph{The necessary outside probability is too large}

By \eqref{contact:two-triple-triangle-exclusion-2}, \eqref{contact:two-triple-triangle-exclusion-13}, and \eqref{contact:two-triple-triangle-exclusion-15},

\[
 b_j-a_i<c_1+c_2=z/K<1/4.
\]

Since $b_j>1$, this forces $a_i>3/4$ for each coordinate. Its selected
middle point also has $c_i>0$. Indeed its contact value is strictly less
than the left-minimum value one, whereas for $c\le 0$ the exact negative
branch difference is

\[
 f(Kc)-f(-Ka)=CK^3(c+a)^2(2a-3/2-c)>0
 \quad(a>3/4).
\]

Centering and the variance identity give

\[
 p_i=\frac{a_i+m_iv_i}{h_i}>\frac{a_i}{h_i}>\frac{39}{100},
 \qquad
 \operatorname{Var}(X_i)/K^2=a_ic_i+r_i h_iv_i.
\tag{C.20.16}\label{contact:two-triple-triangle-exclusion-16}
\]

For the first estimate, $a_i>3/4$ and
$b_i<(2+\sqrt{7})/4<7/6$ give $a_i/h_i>9/23>39/100$.
Moreover

\[
 h_iv_i<h_ib_i\leq\frac{17+7\sqrt7}{16}<\frac94.
\tag{C.20.17}\label{contact:two-triple-triangle-exclusion-17}
\]

To check the middle bound, differentiating along \eqref{contact:two-triple-triangle-exclusion-1} gives
$(h b)'=h(3-4a)/(2b-1)\le 0$ for $a\ge 3/4$.

Let $r=r_1+r_2$. From \eqref{contact:two-triple-triangle-exclusion-13}, \eqref{contact:two-triple-triangle-exclusion-16}--\eqref{contact:two-triple-triangle-exclusion-17}, $a_i<1$, and $c_i>0$,

\[
 \frac94r>
 \frac1{K^2}-\sum_i a_ic_i
 >\left(\frac{49}{60}\right)^2-z.
\]

Hence

\[
 r>r_0-4z/9=:r_*(z),\qquad r_0=\frac{2401}{8100}<3/10.
\tag{C.20.18}\label{contact:two-triple-triangle-exclusion-18}
\]

Also $0<r<1$, because $r_i<a_i/h_i<1/2$. Equation \eqref{contact:two-triple-triangle-exclusion-3}, \eqref{contact:two-triple-triangle-exclusion-16}, and
$r_1r_2\le r^2/4$ now imply

\[
 P_{\rm out}>f(r):=\frac{39}{100}r-r^2/4\geq f(r_*(z)).
\tag{C.20.19}\label{contact:two-triple-triangle-exclusion-19}
\]

For the last inequality, $f$ is concave, so its minimum on
$[r_*(z),1]$ is at an endpoint. The second endpoint has
$f(1)=7/50$, whereas $f(r_*(z))<(39/100)(3/10)<7/50$.

On the other hand, $\phi(z)>3/8$ gives
$N(z)\ge 1/2+(2/3)z \phi(z)>1/2+z/4$. The ratio identity and \eqref{contact:two-triple-triangle-exclusion-13} give

\[
 P_{\rm out}=1-N(z)-CK<\frac{11}{120}-\frac{z}{4}.
\tag{C.20.20}\label{contact:two-triple-triangle-exclusion-20}
\]

Equations \eqref{contact:two-triple-triangle-exclusion-19}--\eqref{contact:two-triple-triangle-exclusion-20} are incompatible. At zero their required gap is
exactly

\[
 f(r_0)-\frac{11}{120}=\frac{103447}{52488000}>0.
\]

For $0\le z\le 1/4$, the derivative of
$f(r_0-4z/9)-(11/120-z/4)$ is at least
$1/4-(4/9)(39/100)=23/300>0$. The gap therefore stays positive.
This is the final contradiction and proves the $321$ exclusion.
\subsection{Two triples: upper collisions}\label{contact:two-triple-collision-322-exclusion}
Normalize the total variance to one; factors of $K$ in the supports are displayed explicitly below.




\[
 K=R+z\phi(z)/(3C)>0,\qquad \phi_0=\phi(0).
\]

For shape $322$, the maximal included cells are exactly $(0,2)$ and
$(2,1)$. A genuine collision at the selected threshold therefore has
precisely these two representations. Let the second coordinate's high
atom have mass $r_2>0$, and put $q=1-r_2$.

\paragraph{The first triple forces a substantially negative threshold}

The first coordinate's deletion levels are $(1,q,q)$. If its successive
physical support gaps are $u,v>0$, its actual centered unit-variance
CDF derivative is

\[
 L=\frac{r_2}{u(u+v)}>0.
\]

The first and second variance derivatives are valid at collisions because
the support and closed threshold are fixed. Their elimination, derived
directly in
Section~\ref{contact:one-triple-bernoulli-collision-211},
gives

\[
 12L+3CR+(z^3+3z)\phi(z)\leq0.
\]

Consequently $z=-s<0$, and substitution of $CR=CK+s \phi(s)/3$ yields

\[
 12L+3CK\leq s(s^2+2)\phi(s)\leq2s\phi_0.
\tag{C.21.1}\label{contact:two-triple-collision-322-exclusion-1}
\]

The last inequality follows by differentiating $(s^2+2)\phi(s)$.

\paragraph{The second triple's equal-level contact prevents that threshold}

The second coordinate's deletion levels are $(1,1,\ell_1)$, where $\ell_1$
is the positive low mass of the first triple. Its low support point is
nonselected, hence a smooth negative minimum of the global cubic dual.
Write its first two support points as $-Ka,Kc$.

The contact at $Kc$ has the same value as the negative minimum, and its
derivative is nonpositive. Therefore the elementary equal-level cubic
calculation gives

\[
 c\geq-1/2.
\tag{C.21.2}\label{contact:two-triple-collision-322-exclusion-2}
\]

For clarity, if $c\le 0$, factor the negative cubic at $-a$: its other
equal-height contact is $c=2a-3/2$ with $a>1/2$, so $c>-1/2$.
If $c>0$, inequality \eqref{contact:two-triple-collision-322-exclusion-2} is immediate. No unique-representation assumption is
used for this contact; the nonpositive derivative follows from the
global majorant at any deletion-CDF jump.

The selected representation $(2,1)$ uses the first coordinate's highest
support point, which is strictly positive by centering and nondegeneracy,
and the second coordinate's middle point $Kc$. Thus

\[
 z>-K/2,\qquad 0<s<K/2.
\tag{C.21.3}\label{contact:two-triple-collision-322-exclusion-3}
\]

Combining \eqref{contact:two-triple-collision-322-exclusion-1} and \eqref{contact:two-triple-collision-322-exclusion-3} gives

\[
 3CK<2s\phi_0<K\phi_0,
 \qquad C<\phi_0/3,
\]

contrary to $C\ge49/120>\phi_0$. This excludes $322$. Exchanging the two
coordinates transposes its grid to $331$ and proves that exclusion too.
\subsection{Unit-span divided differences}\label{contact:unit-span-third-moment slope}


For $s,\mu \in [0,1]$, put $T=s(1-s)$ and

\[
 \rho(s,\mu)=(1-s)\mu^3+s(1-\mu)^3-|s-\mu|^3.
\]

This is exactly the third-moment slope in the three-representation $321$ mass-elimination
chart. One has the unconditional bounds

\[
 \boxed{T/2\le\rho(s,\mu)\le T\max\{1+s,2-s\}\le2T.}
\tag{C.22.1}\label{contact:unit-span-third-moment slope-1}
\]

At $0<s<1$, $\rho/T$ is the second divided difference of $|x|^3$ on
$(-\mu,s-\mu,1-\mu)$. The elementary three-point divided-difference bound in
Section~\ref{contact:three-point-variance-curvature} is half
the support diameter, giving the lower bound $\rho/T\ge 1/2$. The endpoint
cases $s=0,1$ satisfy $\rho=T=0$ directly.

For the upper bound, direct differentiation in $\mu$ gives

\[
 \partial_\mu^2\rho
 =6\{(1-s)\mu+s(1-\mu)-|s-\mu|\}
 =\begin{cases}12\mu(1-s),&\mu\le s,\\
 12s(1-\mu),&\mu\ge s.
 \end{cases}
\]

Thus $\rho$ is convex in $\mu$ on $[0,1]$, including its common derivative at
$\mu=s$, and its maximum is at an endpoint. The endpoint values are exactly
$\rho(s,0)=T(1+s)$ and $\rho(s,1)=T(2-s)$, proving \eqref{contact:unit-span-third-moment slope-1}.

In the complementary-gap two-coordinate chart, both third-moment slopes have
this same $T$. Their half difference $a=(\rho_1-\rho_2)/2$ therefore satisfies

\[
 \boxed{|a|\le3T/4.}
\tag{C.22.2}\label{contact:unit-span-third-moment slope-2}
\]
\subsection{Three representations: contact and variance equations}\label{contact:three-owner-321-contact-reduction}


\paragraph{One exact five-cube and a uniform variance bound}

Translations and a common positive scale put the raw supports and selected
raw threshold into the form

\[
 X_{\rm raw}=(0,s,1),\qquad Y_{\rm raw}=(0,1-s,1),
 \qquad y_{\rm raw}=1,\qquad 0<s<1.
\tag{C.23.1}\label{contact:three-owner-321-contact-reduction-1}
\]

Parameterize the two independent probability vectors by

\[
 p_1=m,\quad p_2=(1-m)r,\quad p_0=(1-m)(1-r),
\]
\[
 q_1=n,\quad q_2=(1-n)t,\quad q_0=(1-n)(1-t).
\tag{C.23.2}\label{contact:three-owner-321-contact-reduction-2}
\]

The single parameter cube is $(s,m,r,n,t) \in [0,1]^5$. Strict positivity of
all six probabilities is a necessary condition. The two means are $\mu,\nu$;
write centered raw support vectors $x,y$, variances $v_1,v_2$, third absolute moments
$r_1,r_2$, and signed second moments
$M_1=\sum_jp_jx_j|x_j|$, $M_2=\sum_jq_jy_j|y_j|$. Set

\[
 V=v_1+v_2,\qquad R=r_1+r_2,\qquad o=1-\mu-\nu,
 \qquad z=o/\sqrt V.
\tag{C.23.3}\label{contact:three-owner-321-contact-reduction-3}
\]

The fixed-two-coordinate bounds imply every normalized support point has
absolute value less than four. For this bound, let $\widehat X_1,\widehat X_2$
be the centered raw laws rescaled to total variance one. Their global
majorants are

\[
 f_i(u)=C|u|^3-Au^2-B_i u+D_i,
 \quad A=(3D+z\phi(z))/2,
 \quad B_i=3C\E[\widehat X_i|\widehat X_i|]+\phi(z).
\]

The already proved bounds $D<11/20$, $|z \phi(z)|<49/200$, and
$C<227/500$ give $A<19/20$ and $|B_i|<881/500$. Global majorization
gives $D_i=f_i(0)\ge 0$; a support contact has $f_i(u)\le 1$. Since $C\ge49/120$,
at $w=|u|\ge 4$ it would follow that

\[
 1>\frac{49}{120}w^3-\frac{19}{20}w^2-\frac{881}{500}w>1.
\]

The last cubic exceeds one at four and increases thereafter. This contradiction
proves the support bound. The common normalized span in \eqref{contact:three-owner-321-contact-reduction-1} is $1/\sqrt{V}<8$.
Consequently the raw quantities satisfy

\[
 \boxed{V>1/64,\qquad R>1/725.}
\tag{C.23.4}\label{contact:three-owner-321-contact-reduction-4}
\]

The third-moment bound follows from $R\ge V^{3/2}/\sqrt{2}$ and
$512 \sqrt{2}<725$. These bounds avoid any division by a vanishing raw variance
or sum of third absolute moments in the compact certificate. Both coordinate variance fractions
also exceed $13/200$, by the dominant-coordinate theorem.

\paragraph{Actual event, ratio, and global contact equations}

The event probability is the actual six-cell sum

\[
 F=1-p_1q_2-p_2(q_1+q_2),\qquad D=F-\Phi(z)>0,
 \qquad C=D V^{3/2}/R.
\tag{C.23.5}\label{contact:three-owner-321-contact-reduction-5}
\]

Thus require

\[
 (49/120)^2R^2\le D^2V^3<(227/500)^2R^2.
\tag{C.23.6}\label{contact:three-owner-321-contact-reduction-6}
\]

For the two coordinates the actual deletion levels are respectively

\[
 H_1=(1,q_0+q_1,q_0),\qquad
 H_2=(1,p_0+p_1,p_0).
\tag{C.23.7}\label{contact:three-owner-321-contact-reduction-7}
\]

Write $L=3D+z \phi(z)$. After multiplication by the positive factor $2VR$,
the global dual contact equation at a raw centered support point $u$ of
coordinate $i$ is

\[
 \boxed{
 2VD(|u|^3-r_i)-LR(u^2-v_i)
 -(6VDM_i+2\sqrt V\phi(z)R)u+2VR(F-H_i(u))=0.
 }
\tag{C.23.8}\label{contact:three-owner-321-contact-reduction-8}
\]

Indeed the unscaled dual price is $D/R$, its quadratic coefficient is
$L/(2V)$, its linear coefficient is $3DM_i/R+\phi(z)/\sqrt{V}$, and averaging
the majorant over the coordinate determines its constant. This proves \eqref{contact:three-owner-321-contact-reduction-8}
directly and preserves the actual centered moments.

The same positive multiplier applied to its raw support derivative gives

\[
 S_i(u)=6VD(u|u|-M_i)-2LRu-2\sqrt V\phi(z)R.
\tag{C.23.9}\label{contact:three-owner-321-contact-reduction-9}
\]

Each selected contact satisfies $S_i(u)\le 0$. Moving the two support points
of each representation by opposite amounts preserves all three raw collision sums.
The resulting gradient pairings are

\[
 \boxed{p_i S_1(x_i)=q_{2-i}S_2(y_{2-i}),\qquad i=0,1,2.}
\tag{C.23.10}\label{contact:three-owner-321-contact-reduction-10}
\]

These are necessary even though none of the six support contacts is a smooth
minimum. Equations \eqref{contact:three-owner-321-contact-reduction-8}--\eqref{contact:three-owner-321-contact-reduction-10} are deliberately redundant; this improves
outward rejection without adding assumptions.

The nonnegative constants of the two global majorants give two further tests:

\[
 2VRF-2VD r_i+LRv_i\geq0.
\tag{C.23.11}\label{contact:three-owner-321-contact-reduction-11}
\]

\paragraph{The original full variance Hessian}

For the two centered signed mass vectors, use

\[
 w_1=(1-s,-1,s),\qquad w_2=(s,-1,1-s),\qquad T=s(1-s).
\]

Both raw variance derivatives equal $T$. Their raw cubic derivatives are
$\rho_i=\sum_j w_{i,j}|x_{i,j}|^3$, where
$x_{1,j}=x_j$ and $x_{2,j}=y_j$ are the centered raw supports.
The actual $321$ mixed event
coefficient is

\[
 M_0=1-s+s^2=1-T.
\]

Put

\[
 W=R\{15D+\phi'''(z)\}/4,
\]
\[
 N_i=WT^2-3DVT\rho_i,
\]
\[
 H_{12}=RV^2M_0+\frac32DVT(\rho_1+\rho_2)-WT^2.
\tag{C.23.12}\label{contact:three-owner-321-contact-reduction-12}
\]

The necessary conditions from the original variance Hessian are

\[
 \boxed{N_1\geq0,\quad N_2\geq0,\quad
 N_1N_2-H_{12}^2\geq0.}
\tag{C.23.13}\label{contact:three-owner-321-contact-reduction-13}
\]

To verify the normalization, start with the independent unnormalized Hessian
derived at the end of
Section~\ref{contact:two-triple-collision-221-reduction},
replace its $G$ by $R/D$, and multiply all matrix entries by $D>0$.
The raw support scale need not equal $K$ for that calculation. Formula \eqref{contact:three-owner-321-contact-reduction-12}
uses the actual $\phi'''(z)=(3z-z^3)\phi(z)$.

In particular \eqref{contact:three-owner-321-contact-reduction-13} is effective at coalescing-support faces: as $s$ tends to
zero or one, $T$ tends to zero while $M_0$ tends to one. The positive bounds
\eqref{contact:three-owner-321-contact-reduction-4} prevent the off-diagonal term $RV^2M_0$ from disappearing. No division by
$s$ or $1-s$ has been introduced.

\paragraph{Further unconditional necessary tests}

The actual canonical variance widths are

\[
 W_1=\min\{(1-\mu)s,\mu(1-s)\},\qquad
 W_2=\min\{(1-\nu)(1-s),\nu s\}.
\]

Lemma~\ref{lem:mixture-width} requires $3(W_1+W_2)>V$ at a violating
candidate. The three excluded cells $(1,2),(2,1),(2,2)$ have raw gaps
$s,1-s,1$ and masses $p_1q_2,p_2q_1,p_2q_2$. Maximizing the threshold gives

\[
 25(p_iq_j)^2V<4\Delta_{ij}^2
\tag{C.23.14}\label{contact:three-owner-321-contact-reduction-14}
\]

for those three pairs.
\subsection{Three representations: the gap product}\label{contact:three-owner-321-gap-collar}
Use the unit-span scale of Section~\ref{contact:three-owner-321-contact-reduction}.



At a genuine attained violating fixed-two-coordinate three-representation $321$
maximum, retain the notation

\[
 T=s(1-s),\quad L=1-T,\quad
 W={15D+\phi'''(z)\over4},\quad C=\kappa V^{3/2}.
\]

The original variance Hessian has diagonal necessary quantities

\[
 N_i=WT^2-3\kappa VT\rho_i,
\]

and its determinant implies

\[
 N_1+N_2\ge V^2L.
\tag{C.24.1}\label{contact:three-owner-321-gap-collar-1}
\]

The pointwise third-moment slope bound
Section~\ref{contact:unit-span-third-moment slope} gives
$\rho_1+\rho_2\ge T$. Hence

\[
 V^2L\le N_1+N_2\le(2W-3\kappa V)T^2.
\tag{C.24.2}\label{contact:three-owner-321-gap-collar-2}
\]

The bound \eqref{eq:joint-normal-bound} gives $W<79/40$. Also $V\le 1/2$
and $C\ge49/120$ yield

\[
 \kappa V={C\over\sqrt V}>{49\sqrt2\over120}>{343\over600},
\]

where the final step uses $\sqrt{2}>7/5$. Thus

\[
 2W-3\kappa V<79/20-343/200=447/200.
\]

Since $L\ge 3/4$, \eqref{contact:three-owner-321-gap-collar-2} proves

\[
 \boxed{T^2>{50\over149}V^2.}
\tag{C.24.3}\label{contact:three-owner-321-gap-collar-3}
\]

The unit-span diameter bound is $V>2401/24964$. Substituting
it into \eqref{contact:three-owner-321-gap-collar-3} gives

\[
 T^2>{144120025\over46428496552}>{1\over324},
 \qquad\boxed{s(1-s)>1/18.}
\tag{C.24.4}\label{contact:three-owner-321-gap-collar-4}
\]

The exact final slack is $66597887/3760708220712>0$, saved in
the computational supplement. Equivalently,

\[
 {3-\sqrt7\over6}<s<{3+\sqrt7\over6}.
\]

In particular $1/18<s<17/18$.
\subsection{Three representations: an upper variance bound}\label{contact:three-owner-321-variance-upper}
The scale is the unit-span scale of Section~\ref{contact:three-owner-321-contact-reduction}.



At an attained violating fixed-two-coordinate maximum with two genuine
three-point laws and all three representations of the closed $321$ event,

\[
 \boxed{V<3/10.}
\tag{C.25.1}\label{contact:three-owner-321-variance-upper-1}
\]

\paragraph{A scalar lower bound}

Retain $D>49/(120\sqrt{2})$ and $-5/4<z<3/5$ from
Section~\ref{contact:fixed-c2-threshold-collar}. Define

\[
 J=3D+(z^3+3z)\phi(z).
\]

Then

\[
 \boxed{J>-9/50.}
\tag{C.25.2}\label{contact:three-owner-321-variance-upper-2}
\]

For $z\ge 0$ this is immediate. For $z=-t$, $0<t<5/4$, the derivative of
$(t^3+3t)\phi(t)$ is $(3-t^4)\phi(t)>0$. Thus its maximum on this interval
is less than $(365/64)\phi(5/4)$. The elementary bound $\phi(0)<399/1000$,
used in Section~\ref{contact:fixed-c2-threshold-collar}, and the alternating exponential
upper polynomial give

\[
 \phi(5/4)<{399\over1000}
 \sum_{j=0}^6{(-25/32)^j\over j!}
 ={9415683872273\over51539607552000}<{183\over1000}.
\]

Also $\sqrt{2}<99/70$, so $3D>343/396$. Consequently

\[
 J>{343\over396}-{365\over64}{183\over1000}
 =-{9\over50}+{631\over253440}>-{9\over50}.
\]

\paragraph{Eliminate the cubic slopes from the mass Hessian}

Write

\[
 T=s(1-s),\quad L=1-T,\quad
 u=(1-s)\nu,\quad v=s\mu,\quad w=m+n,
\]
\[
 E_0=\mu(1-\mu)+\nu(1-\nu),\quad V=E_0-Tw,
\]
\[
 A={3D+z\phi(z)\over2V},\quad
 W={15D+\phi'''(z)\over4},\quad \kappa=D/R.
\]

Applying the centered mass direction $(1-s,-1,s)$ to the first coordinate's
three actual contact equations cancels its constant and linear terms.
The deletion levels $(1,q_0+q_1,q_0)$ give

\[
 \kappa\rho_1-AT=(1-s)q_2-sq_1=u-Ln=:\eta_1.
\]

Similarly, the second coordinate's direction $(s,-1,1-s)$ gives

\[
 \kappa\rho_2-AT=s p_2-(1-s)p_1=v-Lm=:\eta_2.
\tag{C.25.3}\label{contact:three-owner-321-variance-upper-3}
\]

The regularized diagonal Hessian quantities from
Section~\ref{contact:three-owner-321-contact-reduction}
are therefore

\[
 H_i=W-3\kappa V\rho_i/T
     =-{J\over4}-{3V\eta_i\over T},
\tag{C.25.4}\label{contact:three-owner-321-variance-upper-4}
\]

since $W-3VA=-J/4$. The determinant condition implies
$H_1+H_2\ge V^2L/T^2$. Set

\[
 B_0=LE_0-T(u+v).
\]

As $\eta_1+\eta_2=(LV-B_0)/T$, \eqref{contact:three-owner-321-variance-upper-4} yields the scalar necessary inequality

\[
 \boxed{3VB_0-4LV^2-JT^2/2\ge0.}
\tag{C.25.5}\label{contact:three-owner-321-variance-upper-5}
\]

This eliminates both third-moment slopes and their difference from this part of
the Hessian. The identity uses the true independent convolution levels.

\paragraph{Maximize the remaining quadratic in the two means}

For fixed $s$, independently maximize the two concave quadratics in
$\mu,\nu$ appearing in $B_0$. This gives

\[
 {B_0\over L}\le
 { (L-Ts)^2+(L-T(1-s))^2\over4L^2}
 ={(1-2T)(T^2-2T+2)\over4(1-T)^2}=:M(T).
\tag{C.25.6}\label{contact:three-owner-321-variance-upper-6}
\]

The maximizing means lie in $[0,1]$; the displayed inequality also follows
by completing the two squares and does not need this observation.
The function $M$ strictly decreases on $0\le T\le 1/4$. Indeed, with
$L=1-T$, it is

\[
 M=L/2-1/4+1/(2L)-1/(4L^2),\qquad
 {dM\over dL}={L^3-L+1\over2L^3}>0
\]

on $3/4\le L\le 1$. Its useful rational value is

\[
 M(19/100)={513391\over1312200}
 ={47\over120}-{277\over656100}<{47\over120}.
\tag{C.25.7}\label{contact:three-owner-321-variance-upper-7}
\]

Combining \eqref{contact:three-owner-321-variance-upper-2}, \eqref{contact:three-owner-321-variance-upper-5}, \eqref{contact:three-owner-321-variance-upper-6}, and $T^2/L\le 1/12$ gives

\[
 4V^2-3M(T)V-3/400<0.
\tag{C.25.8}\label{contact:three-owner-321-variance-upper-8}
\]

In particular, whenever $T>19/100$, \eqref{contact:three-owner-321-variance-upper-7} implies

\[
 f(V):=4V^2-(47/40)V-3/400<0.
\tag{C.25.9}\label{contact:three-owner-321-variance-upper-9}
\]

But $f(3/10)=0$ exactly and $f$ increases for $V\ge 3/10$.

\paragraph{Two finite bootstrap steps}

The gap bound is $T^2>(50/149)V^2$. If $V\ge 1/3$,
then

\[
 T^2>{50\over1341}>{361\over10000},
\]

so \eqref{contact:three-owner-321-variance-upper-9} contradicts $V\ge 1/3$. Thus first $V<1/3$.

It follows that

\[
 \kappa V={C\over\sqrt V}>{49\sqrt3\over120}>{7\over10},
\]

using $\sqrt{3}>12/7$. The original gap-Hessian derivation now improves to

\[
 V^2L\le(2W-3\kappa V)T^2
 <{37\over20}T^2,
 \qquad T^2>{15\over37}V^2.
\tag{C.25.10}\label{contact:three-owner-321-variance-upper-10}
\]

If $V\ge 3/10$, \eqref{contact:three-owner-321-variance-upper-10} gives

\[
 T^2>{27\over740}>{361\over10000}.
\]

As a convenient weaker consequence of \eqref{contact:three-owner-321-variance-upper-1},
$\kappa\,V>37/50$, since $(37/50)^2\,(3/10)<(49/120)^2$.
Repeating the same gap inequality gives $T^2>(75/173)V^2$.
\subsection{Three representations: joint mean and variance bounds}\label{contact:three-owner-321-joint-mean-bound}
Keep the unit-span scale of Section~\ref{contact:three-owner-321-contact-reduction}.



In the three-representation $321$ case, coupling the two means to the threshold improves the variance bound to

\[
 \boxed{V<11/40.}
\tag{C.26.1}\label{contact:three-owner-321-joint-mean-bound-1}
\]

\paragraph{Couple the quadratic mean bound with the scalar J}

Retain $T,L,B_0,M(T)$ from Section~\ref{contact:three-owner-321-variance-upper} and
$J=3D+(z^3+3z)\phi(z)$, and write $g=T^2/L$ and
$d_{\min}=49/(120\sqrt{2})$. Set

\[
 y=z\sqrt V=1-\mu-\nu,\qquad d=\mu-\nu.
\]

Completing the square in $d$ while retaining $y$ gives

\[
 {B_0\over L}\le N_0(T)+{T\over2L}y-{y^2\over2},
\]
\[
 N_0(T)={1\over2}-{T\over2L}+{T^2(1-4T)\over8L^2}
 =M(T)-{T^2\over8L^2}\le M(T).
\tag{C.26.2}\label{contact:three-owner-321-joint-mean-bound-2}
\]

For $z\le 0$, monotonicity of $(t^2+3)\phi(t)$ for $t\ge 0$ gives

\[
 (z^3+3z)\phi(z)\ge(6/5)z,
 \qquad J>3d_{\min}+(6/5)z.
\tag{C.26.3}\label{contact:three-owner-321-joint-mean-bound-3}
\]

Indeed the derivative of $(t^2+3)\phi(t)$ is
$-t(t^2+1)\phi(t)\le 0$, and its value at zero is $3\phi(0)<6/5$.
In the actual contact/Hessian inequality

\[
 4V\le3B_0/L-Jg/(2V),
\]

substitute \eqref{contact:three-owner-321-joint-mean-bound-2}--\eqref{contact:three-owner-321-joint-mean-bound-3}. The coefficient of $y\le 0$ is

\[
 {3T\over2L}-{3g\over5V^{3/2}}
 ={3T\over2L}\left(1-{2T\over5V^{3/2}}\right)>0
\]

whenever $V\ge 1/4$, because then $T/V^{3/2}\le 2$. Hence that linear term
and $-3y^2/2$ are nonpositive. It follows that

\[
 4V^2-3M(T)V+{3d_{\min}\over2}g<0.
\tag{C.26.4}\label{contact:three-owner-321-joint-mean-bound-4}
\]

For $z\ge 0$, the same inequality follows directly from $J\ge 3D>3dmin$
and $B_0/L\le M(T)$. Thus \eqref{contact:three-owner-321-joint-mean-bound-4} holds for every candidate with $V\ge 1/4$.
In particular, since $3dmin/2>3/7$,

\[
 4V^2-3M(T)V+(3/7)g<0.
\tag{C.26.5}\label{contact:three-owner-321-joint-mean-bound-5}
\]

\paragraph{Two rational bootstrap steps}

The already proved $V<3/10$ gives $\kappa V>37/50$. Therefore the
gap-Hessian inequality gives

\[
 g>{100\over173}V^2.
\tag{C.26.6}\label{contact:three-owner-321-joint-mean-bound-6}
\]

If $V\ge 7/25$, \eqref{contact:three-owner-321-joint-mean-bound-6} implies
$g>196/4325>361/8100=g(19/100)$. Since $g(T)$ increases, $T>19/100$,
and Section~\ref{contact:three-owner-321-variance-upper} gives $M(T)<47/120$. Equation \eqref{contact:three-owner-321-joint-mean-bound-5} would then force

\[
 V<{47/40\over4+300/1211}
 ={56917\over205760}<{7\over25}.
\]

This contradiction proves $V<7/25$. Consequently $\kappa V>77/100$, so

\[
 g>{25\over41}V^2.
\tag{C.26.7}\label{contact:three-owner-321-joint-mean-bound-7}
\]

If $V\ge 11/40$, \eqref{contact:three-owner-321-joint-mean-bound-7} gives

\[
 g>{121\over2624}>{576\over12625}=g(24/125).
\]

Thus $T>24/125$, and exact evaluation gives

\[
 M(24/125)={994301\over2550250}
 ={39\over100}-{593\over5100500}<{39\over100}.
\]

Equation \eqref{contact:three-owner-321-joint-mean-bound-5} now forces

\[
 V<{117/100\over4+75/287}
 ={33579\over122300}<{11\over40},
\]

again a contradiction. This proves \eqref{contact:three-owner-321-joint-mean-bound-1}.
\subsection{Three representations: scalar exclusion}\label{contact:three-owner-321-scalar-closure}
Keep the unit-span scale of Section~\ref{contact:three-owner-321-contact-reduction}.



Use the notation $V,z,T=s(1-s),L=1-T,m,n,w=m+n$, and
$v_0=2401/24964$. The range is

\[
 v_0<V<11/40,\qquad -5/4<z<3/5.
\tag{C.27.1}\label{contact:three-owner-321-scalar-closure-1}
\]

\paragraph{The actual bounded sum and its exact slack}

The raw sum $U$ is in $[0,2]$, has mean $1-o$, variance $V$, and the
selected event is $U\le 1$, where $o=z\sqrt V$. Pointwise
$U(U-1)\le 2\ind_{\{U>1\}}$. Therefore

\[
 F\le1-\{V-o+o^2\}/2.
\tag{C.27.2}\label{contact:three-owner-321-scalar-closure-2}
\]

The precise nonnegative slack, using the actual independent $321$ event,
is

\[
 S:=2(1-F)-(V-o+o^2)
 =T(mq_0+p_0n)+(2-s-s^2)mq_2+(3s-s^2)p_2n.
\]

Both outside coefficients are at least $T$, so

\[
 S\ge T\{m(1-n)+n(1-m)\}
 \ge T\{w-w^2/2\}.
\tag{C.27.3}\label{contact:three-owner-321-scalar-closure-3}
\]

Let $H_{\Sigma}=H_1+H_2$, with $H_i$ defined in
\eqref{contact:three-owner-321-variance-upper-4}. The mass-Hessian
condition gives
\[
 H_{\Sigma}\ge\frac{V^2L}{T^2}\ge12V^2,
 \qquad T=s(1-s)\le\frac14,\quad L=1-T.
\]
The contact identities yield
\[
 \eta_{\mathrm{sum}}=u+v-Lw
 =-\frac{T(H_{\Sigma}+J/2)}{3V}.
\]
Since $J>-9/50$ and $12v_0^2>9/100$, this is negative. Therefore

\[
 w>{u+v\over L}.
\tag{C.27.4}\label{contact:three-owner-321-scalar-closure-4}
\]

If $T>1/5$, both $s$ and $1-s$ belong to $(1/4,3/4)$ and $L<4/5$.
Thus \eqref{contact:three-owner-321-scalar-closure-4} gives $w>5(\mu+\nu)/16$. The range \eqref{contact:three-owner-321-scalar-closure-1} gives
$\mu+\nu=1-z\sqrt V>137/200$, using $\sqrt{11/40}<21/40$. Hence

\[
 w>{137\over640}=:a.
\]

Let $X_1,X_2$ denote the raw marginals, and set $a_1=s$, $a_2=1-s$,
$m_1=m$, $m_2=n$. Since
$a_i\in(1/4,3/4)$, the raw marginal variances satisfy
\[
 v_i\le\E(X_i-a_i)^2<9(1-m_i)/16.
\]
Summing gives

\[
 w\le2-16V/9<2-16v_0/9=:b.
\]

The concave function $f(w)=w-w^2/2$ has its minimum on $[a,b]$ at an
endpoint. Both exact endpoint values exceed $3/20$. Consequently

\[
 \boxed{T>1/5\quad\Longrightarrow\quad S>3/100.}
\tag{C.27.5}\label{contact:three-owner-321-scalar-closure-5}
\]

\paragraph{Combine with the mass Hessian}

The mass Hessian and $\rho_1+\rho_2\ge T$ give

\[
 V^2\le\left\{{15D+\phi'''(z)\over2}-{3C\over\sqrt V}\right\}{T^2\over L}.
\tag{C.27.6}\label{contact:three-owner-321-scalar-closure-6}
\]

If $T\le 1/5$, then $T^2/L\le 1/20$, so
$15D+\phi'''(z)\ge 40V^2+6C/\sqrt{V}$. If $T>1/5$, use $T^2/L\le 1/12$
instead, giving $15D+\phi'''(z)\ge 24V^2+6C/\sqrt{V}$, together with \eqref{contact:three-owner-321-scalar-closure-5}.

The two cases are therefore enclosed by $(A,S_0)=(40,0)$ and
$(A,S_0)=(24,3/100)$. Put $c_0=49/120$ and

\[
 B(V,z;A,S_0)=15(1-\Phi(z))-{15\over2}V(1+z^2)
 +{15\over2}z\sqrt V-{15\over2}S_0
 +(3z-z^3)\phi(z)-AV^2-{6c_0\over\sqrt V}.
\tag{C.27.7}\label{contact:three-owner-321-scalar-closure-7}
\]

The plain probability bound also gives

\[
 P(V,z;A)=15(1-\Phi(z))+(3z-z^3)\phi(z)-AV^2-6c_0/\sqrt V.
\]

When $z<0$, Cantelli gives the additional bound

\[
 Q(V,z;A)=15\{(1+z^2)^{-1}-\Phi(z)\}
 +(3z-z^3)\phi(z)-AV^2-6c_0/\sqrt V.
\tag{C.27.8}\label{contact:three-owner-321-scalar-closure-8}
\]

Every genuine candidate requires each applicable upper expression to
be strictly positive, including when $C=c_0$. The proper selected
downset has probability less than one. The bounded-sum comparison has
positive slack at an interior support point; when the prescribed slack
is positive, \eqref{contact:three-owner-321-scalar-closure-5} gives a
strict improvement over it. For negative thresholds, equality in
Cantelli's inequality would force a two-point sum law, whereas the sum
of two genuine triples has at least five support points. At $z=0$ the
Cantelli expression equals the probability expression. Thus a
nonpositive enclosure excludes the candidate also at $C=c_0$.
It suffices to cover \eqref{contact:three-owner-321-scalar-closure-1},
separately for both $(A,S_0)$ cases, by boxes on which at least one of
$B,P$, or $Q$ (when the whole box has $z\le0$) is nonpositive.

\paragraph{Whole-box derivative enclosures}

For \eqref{contact:three-owner-321-scalar-closure-7}, the exact first derivatives are

\[
 B_V=-{15\over2}(1+z^2)+{15z\over4\sqrt V}-2AV+3c_0/V^{3/2},
\]
\[
 B_z=(z^4-6z^2-12)\phi(z)-15Vz+{15\over2}\sqrt V.
\]

For $P,Q$, their $V$ derivative is $-2AV+3c_0/V^{3/2}$; their $z$
derivatives are respectively $(z^4-6z^2-12)\phi(z)$ and that quantity
minus $30z/(1+z^2)^2$. These formulas apply throughout the original
rectangle, independently of candidate constraints. A midpoint value plus
these whole-box gradient enclosures gives a valid enclosure by the mean
value integral. The numerical domain has $V\ge v_0>0$.
The fixed certificate \texttt{three\_owner321} covers both cases in
\eqref{contact:three-owner-321-scalar-closure-1} and proves that every
point has a nonpositive applicable exclusion expression. Therefore
no maximizing pair can have all three representations of the $321$
threshold. Section~\ref{sec:verification} gives the enclosure and
partition rules.

\subsection{Two representations: contact geometry}\label{contact:two-owner-321-contact-reduction}


Use the contact scale and the strict bounds in
\eqref{eq:contact-basic-collars}:
$K=R_{\mathrm{normalized}}+z \phi(z)/(3C)$, scaled support coordinates equal the
normalized coordinates divided by $K$, and $\tau=C K^3$. Raw moments in
this calculation refer to those scaled supports. Thus

\[
 \frac1{20}<\tau<2,\qquad \frac{25}{64}<V<4,
 \qquad C=\tau V^{3/2}.
\tag{C.28.1}\label{contact:two-owner-321-contact-reduction-1}
\]

The two coordinate variance fractions exceed $13/200$. Define the cubic
potential

\[
 P_\beta(x)=|x|^3-\frac32x^2-3\beta x.
\]

Every global majorant in this scale is a constant plus
$\tau P_\beta$. A nonselected positive contact $b$ has $\beta=b(b-1)$; a
nonselected negative contact $-a$ has $\beta=a(1-a)$.

\paragraph{Adjacent representations}

Write the supports and the selected threshold as

\[
 X=(-t,c,b),\quad Y=(-a,d,e),\quad
 o=e-t=c+d,\qquad d=e-t-c.
\tag{C.28.2}\label{contact:two-owner-321-contact-reduction-2}
\]

Here the representations are $(0,2),(1,1)$. The unused maximal cell $(2,0)$ is
strictly lower: $b-a<o$. The high point $b$ and the low point $-a$ are
nonselected smooth minima, so $b>1/2$ and $a>1/2$.
In fact $b\ge1$. If $1/2<b<1$, then $\beta=b(b-1)<0$
and $P_\beta$ has a negative minimum at $-a_*$ with $a_*>1$.
Its value is less than $-1/2$, whereas
$P_\beta(b)=b^2(3/2-2b)>-1/2$.
The deletion CDF at $-a_*$ is at least its value at $b$;
global majorization therefore rules out this lower minimum.

Also $a<1$. Otherwise $-a$ would be a global minimum of
$P_{a(1-a)}$. Indeed, for $x\ge0$,
\[
 P_{a(1-a)}(x)=x^3-\frac32x^2+3a(a-1)x
 \ge-\frac12\ge P_{a(1-a)}(-a),
\]
and for $x=-t\le0$ the difference from the value at $-a$ is
$(t-a)^2(t+2a-3/2)\ge0$.
The contact level at $-a$ is one, so a global minimum there would
preclude the smaller contact levels at the other support points.

The nonpositive derivative at the selected high point gives
$e\le B(a)$, where $B$ is defined below. The unused-corner inequality
now gives
\[
 t+b<a+e\le a+B(a)\le2.
\]
Since $b\ge1$, we have $t<1$. The nonpositive derivative at $-t$
gives $b(b-1)\ge t(1-t)$, hence $b\ge B(t)$.
The function $x+B(x)$ is strictly increasing on $(0,1)$, because
its derivative is $1+(1-2x)/\sqrt{1+4x(1-x)}>0$.
Thus $t\ge a$ would contradict the strict span inequality.
Finally $e>t+b-a>t$, so $o=e-t>0$. We have proved

\[
 \frac12<a<1,\quad 0<t<a,\quad
 B(t)\leq b<2-t,\quad 0<e\leq B(a),\quad -t<c<b,
\]
\[
 B(v)=\frac{1+\sqrt{1+4v(1-v)}}2,\qquad
 A_2:=a+e>A_1:=t+b,\qquad o>0.
\tag{C.28.3}\label{contact:two-owner-321-contact-reduction-3}
\]

The strict span ordering also implies $-a<d<e$; it need not be imposed by
a separate chart. Set $h=t+c=e-d>0$, $v=b-c>0$.

The first high contact has actual deletion level $\ell_2$; the second low
contact has level one. Define

\[
 g_b(x)=P_{b(b-1)}(x)-P_{b(b-1)}(b),\qquad
 G=g_b(-t),\quad g=g_b(c),
\]
\[
 U_a(x)=P_{a(1-a)}(-a)-P_{a(1-a)}(x)
 =a^2(3/2-2a)-|x|^3+\frac32x^2+3a(1-a)x.
\tag{C.28.4}\label{contact:two-owner-321-contact-reduction-4}
\]

The actual deletion levels of the closed $321$ event are
$(1,1-r_2,\ell_2)$ and $(1,1-r_1,\ell_1)$. Subtracting the contacts at the
respective nonselected minima determines all probabilities:

\[
 (\ell_1,m_1,r_1)
 =(1-\tau U_a(e),\ \tau[U_a(e)-U_a(d)],\ \tau U_a(d)),
\]
\[
 (\ell_2,m_2,r_2)
 =(1-\tau G,\ \tau g,\ \tau[G-g]).
\tag{C.28.5}\label{contact:two-owner-321-contact-reduction-5}
\]

Consequently a genuine candidate requires
$U_a(e)>U_a(d)>0$ and $G>g>0$. First-coordinate centering is exactly

\[
 t=\tau\{hU_a(e)+vU_a(d)\}.
\]

Its denominator is strictly positive, so set

\[
 \boxed{\tau=\frac{t}{hU_a(e)+vU_a(d)}.}
\tag{C.28.6}\label{contact:two-owner-321-contact-reduction-6}
\]

The second-coordinate centering equation is retained as the first residual:

\[
 \boxed{\tau\{(a+e)G-hg\}-a=0.}
\tag{C.28.7}\label{contact:two-owner-321-contact-reduction-7}
\]

Define the nonnegative selected slope magnitudes

\[
 H_t=3[t^2-t+b(b-1)],\qquad
 H_c=3[b(b-1)+c-c|c|],
\]
\[
 H_d=3[a(1-a)+d-d|d|],\qquad
 H_e=3[a(1-a)+e-e^2].
\]

Collision-preserving support variations give the two further residuals

\[
 \boxed{\ell_1H_t-r_2H_e=0,\qquad m_1H_c-m_2H_d=0.}
\tag{C.28.8}\label{contact:two-owner-321-contact-reduction-8}
\]

Put $\sigma=\ell_1 H_t+m_1 H_c>0$. By \eqref{contact:two-owner-321-contact-reduction-8} it also equals
$m_2 H_d+r_2 H_e$.

There is a useful additional global-majorant constraint. The unused maximal
cell $(2,0)$ creates a deletion jump at $y_*=o-b$, with
$-a<y_*<d$. At that point the second deletion CDF is still one. Its majorant
is $1-\tau U_a$, so

\[
 \boxed{U_a(o-b)\leq0.}
\tag{C.28.9}\label{contact:two-owner-321-contact-reduction-9}
\]

The corresponding first-coordinate deletion jump lies above $b$. Its
majorant bound follows automatically from increase of $P_{b(b-1)}$ to the
right of the positive minimum.

A single closed five-cube covers \eqref{contact:two-owner-321-contact-reduction-3}, with strict conditions retained as
rejection tests. For $(k_1,\ldots ,k_5) \in [0,1]^5$, take

\[
 a=\frac{1+k_1}{2},\quad t=ak_2,\quad
 b=B(t)+(2-t-B(t))k_3,\quad e=B(a)k_4,
 \quad c=-t+(t+b)k_5.
\tag{C.28.10}\label{contact:two-owner-321-contact-reduction-10}
\]

All valid adjacent-representation parameters have a preimage. The cube contains
extraneous parameters, including failures of $a+e>t+b$; they are rejected
by necessary tests. Transposition covers the other adjacent representation pair.

\paragraph{Endpoint representations}

Now the representations are $(0,2),(2,0)$ and the unused middle-middle cell lies
strictly below their common threshold. Both nonselected middle contacts
are negative smooth minima. A positive smooth minimum would force a positive derivative at the larger contact, contrary to Lemma~\ref{lem:full-law-contacts}; zero cannot be a minimum because the quadratic coefficient is negative. Write

\[
 X=(-t_1,-a_1,b_1),\qquad Y=(-t_2,-a_2,b_2),
\]
\[
 \frac12<a_i<1,\quad a_i<t_i<4,\quad 0<b_i\leq B(a_i),
 \quad A=t_1+b_1=t_2+b_2,
\]
\[
 o=b_2-t_1=b_1-t_2>-a_1-a_2.
\tag{C.28.11}\label{contact:two-owner-321-contact-reduction-11}
\]

Put $u_i=t_i-a_i$, $v_i=a_i+b_i$, and define

\[
 E_i=u_i^2(t_i+2a_i-3/2)>0,\qquad U_i=U_{a_i}(b_i)>0.
\tag{C.28.12}\label{contact:two-owner-321-contact-reduction-12}
\]

Subtraction of the actual contact levels gives

\[
 (\ell_1,m_1,r_1)
 =(1-\tau(E_2+U_2),\ \tau U_2,\ \tau E_2),
\]
\[
 (\ell_2,m_2,r_2)
 =(1-\tau(E_1+U_1),\ \tau U_1,\ \tau E_1).
\tag{C.28.13}\label{contact:two-owner-321-contact-reduction-13}
\]

The identity defining $E_i$ is the exact factorization
$P_{a_i(1-a_i)}(-t_i)-P_{a_i(1-a_i)}(-a_i)=E_i$.
First-coordinate centering therefore eliminates the positive parameter

\[
 \boxed{\tau=\frac{t_1}{U_2u_1+E_2A},}
\tag{C.28.14}\label{contact:two-owner-321-contact-reduction-14}
\]

and the retained second centering residual is

\[
 \boxed{\tau(U_1u_2+E_1A)-t_2=0.}
\tag{C.28.15}\label{contact:two-owner-321-contact-reduction-15}
\]

Let

\[
 H_{t_i}=3u_i(t_i+a_i-1)>0,\qquad
 H_{b_i}=3[a_i(1-a_i)+b_i-b_i^2].
\]

The two actual representation pairings are

\[
 \boxed{\ell_1H_{t_1}-r_2H_{b_2}=0,\qquad
 r_1H_{b_1}-\ell_2H_{t_2}=0.}
\tag{C.28.16}\label{contact:two-owner-321-contact-reduction-16}
\]

Set $\sigma=\ell_1 H_{t_1}+r_1 H_{b_1}>0$. Equations \eqref{contact:two-owner-321-contact-reduction-16} equate it with
the corresponding sum for the second coordinate. They also imply the two
high slope magnitudes are strictly positive.

The unused middle-middle cell produces deletion jumps at $o+a_2$ and
$o+a_1$, strictly between the respective middle and high support points.
Immediately at each jump the relevant deletion level still equals its
nonselected middle contact level. Global majorization therefore requires

\[
 \boxed{U_{a_1}(o+a_2)\leq0,\qquad U_{a_2}(o+a_1)\leq0.}
\tag{C.28.17}\label{contact:two-owner-321-contact-reduction-17}
\]

A single closed five-cube covers all of \eqref{contact:two-owner-321-contact-reduction-11}:

\[
 a_1=(1+k_1)/2,\quad a_2=(1+k_2)/2,
\quad t_1=a_1+(4-a_1)k_3,\quad t_2=a_2+(4-a_2)k_4,
\]
\[
 b_1=B(a_1)k_5,\qquad b_2=t_1+b_1-t_2.
\tag{C.28.18}\label{contact:two-owner-321-contact-reduction-18}
\]

The necessary tests retain $0<b_2\le B(a_2)$ and the unused middle-cell
inequality in \eqref{contact:two-owner-321-contact-reduction-11}. The denominator in \eqref{contact:two-owner-321-contact-reduction-14} is positive at a genuine
candidate. A numerical evaluator may divide by it only when its enclosure
is positive on the complete box.

\paragraph{Three common residuals and the original full Hessian}

For either family, use the probabilities just derived to sum raw centered
variances $v_i$, third absolute moments $r_i$, and totals $V,R$. Let

\[
 F=\ell_1+m_1(\ell_2+m_2)+r_1\ell_2
   =1-m_1r_2-r_1(m_2+r_2),\qquad z=o/\sqrt V.
\]

The common scale, density and actual-CDF equations are

\[
 \boxed{3(V-R)-\sigma o=0,\quad
 \tau\sigma\sqrt V-\phi(z)=0,\quad
 F-\Phi(z)-\tau R=0.}
\tag{C.28.19}\label{contact:two-owner-321-contact-reduction-19}
\]

Together with \eqref{contact:two-owner-321-contact-reduction-7},\eqref{contact:two-owner-321-contact-reduction-8}, or respectively \eqref{contact:two-owner-321-contact-reduction-15},\eqref{contact:two-owner-321-contact-reduction-16}, these give six necessary
equations on each five-cube. The first equation follows from the definition
of $K$, the second from $\E f_i'(X_i)=-\phi(z)$, and the third from the actual
ratio $C=\tau V^{3/2}$.

The strict ratio tests are

\[
 (49/120)^2\le \tau^2V^3<(227/500)^2.
\tag{C.28.20}\label{contact:two-owner-321-contact-reduction-20}
\]

Also $\tau R>49/(120 \sqrt{2})$ follows from \eqref{contact:two-owner-321-contact-reduction-20} and the two-coordinate
moment bound $R\ge V^{3/2}/\sqrt{2}$. This is a necessary discrepancy lower
bound, not an additional variational assumption.

The full variance Hessian can be evaluated in the original signed-mass
basis. For each support vector $(x_0,x_1,x_2)$, set

\[
 w_i=\left(\frac{x_2-x_1}{x_2-x_0},-1,
             \frac{x_1-x_0}{x_2-x_0}\right),\quad
 T_i=(x_1-x_0)(x_2-x_1),\quad
 \rho_i=\sum_jw_{i,j}|x_j|^3.
\]

Let $M_0=\sum_{j+k\le 2}w_{1,j}w_{2,k}$, the actual $321$ mixed event
coefficient, and put

\[
 W=\{15\tau R+\phi'''(z)\}/4,
\quad N_i=WT_i^2-3\tau VT_i\rho_i,
\]
\[
 Q=V^2M_0+\frac32\tau V(T_1\rho_2+T_2\rho_1)-WT_1T_2.
\]

Necessary second-order conditions are

\[
 N_1\geq0,\qquad N_2\geq0,\qquad N_1N_2-Q^2\geq0.
\tag{C.28.21}\label{contact:two-owner-321-contact-reduction-21}
\]

These are the already derived original Hessian entries multiplied by the
positive scalar $\tau$, with the actual normal derivative
$\phi'''(z)=(3z-z^3)\phi(z)$.

Finally retain the canonical mixture-width exclusion, both variance-fraction
bounds, and the three excluded-cell atom-gap inequalities. In the adjacent
case the excluded cells $(1,2),(2,1),(2,2)$ have gaps $h,v,A_1$; in the
endpoint case their gaps are $u_1,u_2,A$. Each satisfies
the inequality $25p_*^2V<4\delta_*^2$, where $p_*$ is the atom's
probability and $\delta_*$ its raw gap. These tests and
\eqref{contact:two-owner-321-contact-reduction-9},\eqref{contact:two-owner-321-contact-reduction-17} are
inequalities; they need not enter the full-box derivative system.

\subsection{Adjacent representations: scalar exclusion}\label{contact:adjacent-owner-321-scalar-reduction}


\paragraph{Normalize by the larger span}

Translate each support to start at zero and divide by the second span. The
supports and selected threshold become

\[
 X\in\{0,a,r\},\quad Y\in\{0,1-a,1\},\quad
 0<a<r<1,\quad X+Y\le1.
\tag{C.29.1}\label{contact:adjacent-owner-321-scalar-reduction-1}
\]

Let $V$ denote total variance in this raw scale, $z$ the standardized
selected threshold, $D=F-\Phi(z)$, and $C=D V^{3/2}/R$. The uniform
diameter and threshold bounds give

\[
 v_0:=2401/24964<V\le(1+r^2)/4\le1/2,\qquad 0<z<3/5.
\tag{C.29.2}\label{contact:adjacent-owner-321-scalar-reduction-2}
\]

Put $s=a/r$, $T=s(1-s)$, and use the centered mass directions
$(1-s,-1,s)$ and $(a,-1,1-a)$. Their raw variance derivatives are
$T_1=a(r-a)$ and $T_2=a(1-a)$. Summation over the actual $321$ cells gives
the mixed event derivative

\[
 M=1-a+a^2/r=1-rT.
\tag{C.29.3}\label{contact:adjacent-owner-321-scalar-reduction-3}
\]

The cubic third-moment derivatives divided by $T_i$ are the centered second
divided differences $\gamma_i$. The elementary divided-difference diameter
bound gives $\gamma_1\ge r/2$ and $\gamma_2\ge 1/2$.

\paragraph{Retain the dependence between the mixed coefficient and span}

Write $W=(15D+\phi'''(z))/4$ and
$H_i=W-3C \gamma_i/\sqrt{V}$. The mass-Hessian condition
\eqref{contact:two-owner-321-contact-reduction-21}, divided by the positive variance-direction factors,
implies

\[
 H_1+H_2\ge V^2\Lambda,\qquad \Lambda=M/(T_1T_2).
\tag{C.29.4}\label{contact:adjacent-owner-321-scalar-reduction-4}
\]

This follows directly from its determinant: the off-diagonal entry is
$V^2 \Lambda-(H_1+H_2)/2$, whose absolute value is at most
$\sqrt{H_1H_2}\le (H_1+H_2)/2$.

Now $a(1-a)\le 1/4$, and
$T_1\,T_2=r^2 T\,a(1-a)$. Therefore

\[
 \Lambda\ge {4(1-rT)\over r^2T}
 \ge {16-4r\over r^2}\ge{12\over r^2}.
\tag{C.29.5}\label{contact:adjacent-owner-321-scalar-reduction-5}
\]

Combining \eqref{contact:adjacent-owner-321-scalar-reduction-4}, \eqref{contact:adjacent-owner-321-scalar-reduction-5}, and the two divided-difference bounds proves

\[
 \boxed{15D+\phi'''(z)\ge
 {32-8r\over r^2}V^2+{3C(1+r)\over\sqrt V}.}
\tag{C.29.6}\label{contact:adjacent-owner-321-scalar-reduction-6}
\]

The universal bound $15D+\phi'''(z)<79/10$ also gives
$V^2 \Lambda<79/20$. Hence
$r^2>(240/79)V^2>(240/79)v_0^2>1/36$, proving the convenient
strict bound $r>1/6$. The final comparison is exact rational arithmetic.

\paragraph{The bounded sum supplies the scalar upper bounds}

Set $U=X+Y$. It belongs to $[0,1+r]$, has mean $1-z\sqrt V$, and its
selected event is $U\le 1$. Pointwise,

\[
 U(U-1)\le r(1+r)\ind_{\{U>1\}}.
\]

Consequently

\[
 F\le1-{V(1+z^2)-z\sqrt V\over r(1+r)},\qquad F\le1.
\tag{C.29.7}\label{contact:adjacent-owner-321-scalar-reduction-7}
\]

Put $c_0=49/120$, $A(r)=32/r^2-8/r$, and

\[
 P(r,V,z)=15(1-\Phi(z))+(3z-z^3)\phi(z)
 -A(r)V^2-{3c_0(1+r)\over\sqrt V},
\]
\[
 B(r,V,z)=P(r,V,z)
 -{15\{V(1+z^2)-z\sqrt V\}\over r(1+r)}.
\tag{C.29.8}\label{contact:adjacent-owner-321-scalar-reduction-8}
\]

Every genuine candidate requires both $P>0$ and $B>0$, including when
$C=c_0$. The selected downset is proper, so $F<1$. The bounded-sum
comparison is also strict: the sum has an interior support point at
which its pointwise quadratic bound has positive slack. Together with
\eqref{contact:adjacent-owner-321-scalar-reduction-6}, these strict
comparisons show that a nonpositive enclosure excludes the candidate
at the equality floor too. It suffices to cover the closed box

\[
 [1/6,1]\times[v_0,1/2]\times[0,3/5]
\tag{C.29.9}\label{contact:adjacent-owner-321-scalar-reduction-9}
\]

by boxes on each of which one expression in
\eqref{contact:adjacent-owner-321-scalar-reduction-8} is nonpositive.

\paragraph{Whole-box derivatives for a fixed outward cover}

Let $N=V(1+z^2)-z\sqrt V$. On the whole domain
\eqref{contact:adjacent-owner-321-scalar-reduction-9}, the exact derivatives are

\[
 P_r=(64/r^3-8/r^2)V^2-3c_0/\sqrt V,
\]
\[
 P_V=-2A(r)V+{3c_0(1+r)\over2V^{3/2}},\qquad
 P_z=(z^4-6z^2-12)\phi(z),
\]
\[
 B_r=P_r+{15N(1+2r)\over r^2(1+r)^2},
\]
\[
 B_V=P_V-{15\{1+z^2-z/(2\sqrt V)\}\over r(1+r)},\qquad
 B_z=P_z-{15(2Vz-\sqrt V)\over r(1+r)}.
\]

Thus midpoint values plus whole-box derivative enclosures give a valid
mean-value enclosure. Both denominators $r$ and $V$ are uniformly positive.
The fixed records \texttt{adjacent\_owner321} and
\texttt{adjacent\_owner321\_completion} cover
\eqref{contact:adjacent-owner-321-scalar-reduction-9}; at every point
they prove $P\le0$ or $B\le0$. Thus the adjacent-representation
case has no maximizing pair. Section~\ref{sec:verification} gives
the enclosure and partition rules, including the subdivision of the
four cells treated by the completion record.

\subsection{Endpoint representations: scalar exclusion}\label{contact:endpoint-owner-321-scalar-reduction}


\paragraph{Geometry, variance directions, and the first scalar inequality}

Translate both supports to start at zero and divide by their common span:

\[
 X\in\{0,a,1\},\quad Y\in\{0,b,1\},\quad
 a,b>0,\quad s:=a+b<1,
 \qquad U=X+Y\le1.
\tag{C.30.1}\label{contact:endpoint-owner-321-scalar-reduction-1}
\]

Let the respective means be $\mu,\nu$, middle masses $m,n$, and total raw
variance and sum of third absolute moments be $V,R$. Write

\[
 D=F-\Phi(z),\quad C=DV^{3/2}/R,\quad
 T_1=a(1-a),\quad T_2=b(1-b),\quad
 M=a+b-ab,\quad g=T_1T_2/M.
\]

The diameter and threshold results give

\[
 v_0:=2401/24964<V\le1/2,\qquad -5/4<z<3/5,
 \qquad D>{49\over120\sqrt2},\quad C\ge c_0:=49/120.
\tag{C.30.2}\label{contact:endpoint-owner-321-scalar-reduction-2}
\]

The actual mixed derivative of the $321$ event in the centered mass
directions $(1-a,-1,a)$ and $(1-b,-1,b)$ is $M$. The original mass Hessian,
as in the two-representation reduction, implies

\[
 15D+\phi'''(z)\ge {2V^2\over g}+{6C\over\sqrt V},
\tag{C.30.3}\label{contact:endpoint-owner-321-scalar-reduction-3}
\]

using the lower bound $1/2$ for each centered cubic divided difference.

We will use the uniform bound

\[
 \boxed{g<1/11.}
\tag{C.30.4}\label{contact:endpoint-owner-321-scalar-reduction-4}
\]

Indeed put $p=ab$, $d=(1-a)(1-b)=1-s+p$; then $g=pd/(s-p)$.
For fixed $s$, its derivative in $p$ has positive numerator
$s(1-s+2p)-p^2$. Thus its maximum occurs at $p=s^2/4$, or $a=b=x$.
The resulting function is $x(1-x)^2/(2-x)$, $0<x\le 1/2$, whose derivative
has the sign of $x^2-3x+1$. Its maximum is
$(5\sqrt{5}-11)/2<1/11$; the last comparison follows from
$\sqrt{5}<123/55$, since $123^2-5\,55^2=4>0$.

\paragraph{At moderately negative thresholds the actual contacts give mass bounds}

For $z\ge -3/4$, define $J=3D+(z^3+3z)\phi(z)$. Then

\[
 J>{343\over396}-{171\over64}{31\over100}
 ={24001\over633600}>0.
\tag{C.30.5}\label{contact:endpoint-owner-321-scalar-reduction-5}
\]

Here $\sqrt{2}<99/70$, and $\phi(3/4)<31/100$. The latter follows from
$\phi(0)<399/1000$ and the degree-four upper Taylor polynomial for
$\exp(-9/32)$. On $[0,3/4]$, both
$t(t^2+3)\phi(t)$ and $t(t^2+2)\phi(t)$ increase; their derivatives are
respectively $(3-t^4)\phi(t)$ and $(2+t^2-t^4)\phi(t)$.

Let $\kappa=D/R$, $A=(3D+z\phi(z))/(2V)$, and
$W=(15D+\phi'''(z))/4$. Applying each centered mass direction to its three
actual contact equations gives

\[
 \eta_1:=\kappa\rho_1-AT_1=(1-a)\nu-Mn,
 \qquad
 \eta_2:=\kappa\rho_2-AT_2=(1-b)\mu-Mm.
\tag{C.30.6}\label{contact:endpoint-owner-321-scalar-reduction-6}
\]

The diagonal Hessian quantities satisfy
$H_i=W-3\kappa V \rho_i/T_i=-J/4-3V \eta_i/T_i\ge 0$.
Thus \eqref{contact:endpoint-owner-321-scalar-reduction-5} gives $\eta_i<0$. Both middle contacts are negative after centering,
so $\mu>a$ and $\nu>b$. Consequently

\[
 \boxed{m>{a(1-b)\over M},\qquad n>{b(1-a)\over M}.}
\tag{C.30.7}\label{contact:endpoint-owner-321-scalar-reduction-7}
\]

Define the cubic contact scale in the present unit-span coordinates by

\[
 k={\sqrt V\{D+z\phi(z)/3\}\over C}.
\tag{C.30.8}\label{contact:endpoint-owner-321-scalar-reduction-8}
\]

A nonselected negative minimum of the cubic majorant has distance strictly
larger than $k/2$ from zero: its second derivative is positive, and the
quadratic coefficient is $-3k/2$ relative to the cubic coefficient. Hence
$\mu-a>k/2$ and $\nu-b>k/2$. Positivity of the low masses now gives

\[
 \boxed{1-m>{k\over2(1-a)},\qquad
 1-n>{k\over2(1-b)}.}
\tag{C.30.9}\label{contact:endpoint-owner-321-scalar-reduction-9}
\]

\paragraph{A uniform positive lower bound on this raw contact scale}

Combining \eqref{contact:endpoint-owner-321-scalar-reduction-3}, \eqref{contact:endpoint-owner-321-scalar-reduction-4}, and \eqref{contact:endpoint-owner-321-scalar-reduction-8} yields

\[
 k\ge{2\over5}+{\sqrt V\over15C}
       \{22V^2+(z^3+2z)\phi(z)\}.
\tag{C.30.10}\label{contact:endpoint-owner-321-scalar-reduction-10}
\]

For $z\ge 0$, the braces are positive. For $-3/4\le z\le 0$, their value exceeds
$1/5-(123/64)(31/100)=-2533/6400$, because $22v_0^2>1/5$.
Moreover $\sqrt{V}/(15C)<(5/7)/(49/8)=40/343$. In either sign case,

\[
 \boxed{k>{7\over20}.}
\tag{C.30.11}\label{contact:endpoint-owner-321-scalar-reduction-11}
\]

The exact margin in the negative sign case is
$2/5-(40/343)(2533/6400)-7/20=211/54880>0$.

\paragraph{Exact event slack, retaining the independent probabilities}

Let $p_0,m,p_2$ and $q_0,n,q_2$ be the actual probabilities. The nonnegative
bounded-sum slack is

\[
 S:=2(1-F)-\{V-z\sqrt V+z^2V\}
\]
\[
 =T_1mq_0+T_2p_0n+s(1-s)mn
 +(1-a)(2+a)mq_2+(1-b)(2+b)p_2n.
\tag{C.30.12}\label{contact:endpoint-owner-321-scalar-reduction-12}
\]

Since $\mu>a$, $\nu>b$, the high masses satisfy
$p_2>a(1-m)$ and $q_2>b(1-n)$. After grouping each low/high pair in \eqref{contact:endpoint-owner-321-scalar-reduction-12},

\[
 S>A_0m(1-n)+B_0n(1-m)+s(1-s)mn,
\]

where $A_0=(1-a)(a+2b)$ and $B_0=(1-b)(b+2a)$. Inserting \eqref{contact:endpoint-owner-321-scalar-reduction-7}, \eqref{contact:endpoint-owner-321-scalar-reduction-9} into
the first two positive terms proves

\[
 S>{k\over2M}\{aA_0+bB_0\}
 ={k\over2M}\{s^2(1-s)+p(s+2)\}.
\tag{C.30.13}\label{contact:endpoint-owner-321-scalar-reduction-13}
\]

The expression in braces is at least $4pd$. To see this, subtract $4pd$:
$s^2(1-s)+p(5s-2)-4p^2$ is concave in $p \in [0,s^2/4]$. Its endpoint
values are $s^2(1-s)\ge 0$ and $s^2(2+s-s^2)/4>0$.
Using $g=pd/M$, \eqref{contact:endpoint-owner-321-scalar-reduction-13} therefore yields the simple exact inequality

\[
 \boxed{S>2kg\qquad(z\ge-3/4).}
\tag{C.30.14}\label{contact:endpoint-owner-321-scalar-reduction-14}
\]

In particular, \eqref{contact:endpoint-owner-321-scalar-reduction-11} proves

\[
 \boxed{g>1/20,\ z\ge-3/4\quad\Longrightarrow\quad S>7/200.}
\tag{C.30.15}\label{contact:endpoint-owner-321-scalar-reduction-15}
\]

\paragraph{Three fixed scalar rectangles suffice}

The expressions $P,B,Q$ below are upper bounds obtained respectively from $F\le 1$,
the bounded sum with slack $S_0$, and Cantelli for $z\le 0$:

\[
 P=15(1-\Phi(z))+(3z-z^3)\phi(z)-A_0V^2-6c_0/\sqrt V,
\]
\[
 B=P-{15\over2}\{V(1+z^2)-z\sqrt V+S_0\},\qquad
 Q=P+15\{(1+z^2)^{-1}-1\}.
\tag{C.30.16}\label{contact:endpoint-owner-321-scalar-reduction-16}
\]

Every candidate requires each applicable expression to be strictly positive,
including at $C=c_0$, by the strict probability and bounded-sum
comparisons used in Section~\ref{contact:three-owner-321-scalar-closure}.
For the positive prescribed slack use
\eqref{contact:endpoint-owner-321-scalar-reduction-15}; for negative
thresholds Cantelli equality is impossible for a sum with at least five
support points.
The following three closed rectangles cover all possibilities:

\begin{center}\small
\begin{tabular}{llll}\toprule
Case & Coefficient $A_0$ & Slack $S_0$ & $z$ interval\\\midrule
$g\le 1/20$ & 40 & 0 & $[-5/4,3/5]$\\
$g>1/20$, negative tail & 22 & 0 & $[-5/4,-3/4]$\\
$g>1/20$, remaining thresholds & 22 & $7/200$ & $[-3/4,3/5]$\\
\bottomrule\end{tabular}\end{center}

Every row has the same $V$ interval $[2401/24964,1/2]$. The first row uses
\eqref{contact:endpoint-owner-321-scalar-reduction-3} directly; the others use \eqref{contact:endpoint-owner-321-scalar-reduction-4}, and the last also uses \eqref{contact:endpoint-owner-321-scalar-reduction-15}.
Closed shared faces do not omit any genuine candidate.

The whole-box derivatives are the ones in
Section~\ref{contact:three-owner-321-scalar-closure}, with
the displayed coefficient and slack. They hold independently of all
candidate restrictions throughout these rectangles.
The fixed certificate \texttt{endpoint\_owner321} covers all three
rectangles and proves that every point has a nonpositive applicable
expression in \eqref{contact:endpoint-owner-321-scalar-reduction-16}.
Hence the endpoint-representation case has no maximizing pair.
Section~\ref{sec:verification} gives the enclosure and partition rules.

\subsection{Two triples: the collided asymmetric event}\label{contact:two-triple-collision-221-reduction}


\[
 K=R+z\phi(z)/(3C),\qquad \tau=CK^3,
 \qquad \sigma=\phi(z)/(CK^2).
\]

The elementary fixed-two-coordinate bounds give
$1/2<K<8/5$, $1/20<\tau<2$, and $49/120\le C<227/500$.
All support quantities below are divided by $K$.

\paragraph{Supports, actual levels, and exact compact charts}

The maximal included cells are $(1,1)$ and $(2,0)$. A genuine selected
collision therefore has these two representations. Write

\[
 T_1=(-a,c,b),\qquad T_2=(-t,d,e),
 \qquad h=b-c=t+d>0,
\]
\[
 u=a+c>0,\qquad A_1=a+b=u+h,
 \qquad v=e-d>0,\qquad A_2=t+e=h+v.
\tag{C.31.1}\label{contact:two-triple-collision-221-reduction-1}
\]

The unused corner $(0,2)$ is excluded, giving the strict geometric condition

\[
 A_2>A_1,\qquad\text{equivalently }v>u.
\tag{C.31.2}\label{contact:two-triple-collision-221-reduction-2}
\]

Denote low/middle/high probabilities by $\ell_i,m_i,r_i$, and put
$q_i=\ell_i+m_i=1-r_i$. The first triple's deletion levels are
$(q_2,q_2,\ell_2)$; the second triple's levels are $(1,q_1,0)$.
The first low support and the second high support are nonselected. They are
smooth minima of their global cubic majorants.

In the first triple the negative minimum $-a$ and the selected middle contact
$c$ have equal height and nonpositive derivatives. The complete equal-height
charts from the contact formulas above are

\[
 \begin{array}{c|c|c}
 &a(k)&c(k)\\ \hline
 \text{negative}&\frac12+\frac{k}{4}&\frac{k-1}{2}\\[2mm]
 \text{positive}&\frac{3(1+k)^2}{2(k^3+3k+2)}&k a(k)
 \end{array}
 \quad 0\leq k\leq1.
\tag{C.31.3}\label{contact:two-triple-collision-221-reduction-3}
\]

The next positive minimum is

\[
 B(a)=\frac{1+\sqrt{1+4a(1-a)}}2.
\]

The selected high contact satisfies $0<b\le B(a)$. Set $b=wB(a)$,
$0\le w\le 1$, retaining $b>c$ as a strict necessary condition. The genuine
level drop below the negative minimum ensures $a<1$; including the boundary
$a=1$ in \eqref{contact:two-triple-collision-221-reduction-3} is harmless because the level drop is tested separately.

For the second triple its nonselected high contact has level zero. Its scaled
cubic majorant is exactly

\[
 g_e(x)=|x|^3-\frac32x^2-3e(e-1)x+B_0,
 \qquad B_0=e^2(2e-3/2).
\tag{C.31.4}\label{contact:two-triple-collision-221-reduction-4}
\]

Necessarily $e\ge 1$. Indeed a smooth positive minimum has $e>1/2$; if
$1/2<e<1$, the corresponding negative minimum has smaller potential value.
That contradicts global majorization of a nonnegative deletion CDF when
$g_e(e)=0$. More generally the same argument applies to a nonselected positive
minimum at any deletion level, since the deletion CDF is nonincreasing.

At the low contact set

\[
 G=g_e(-t)=t^3-\frac32t^2+3e(e-1)t+B_0=1/\tau.
\tag{C.31.5}\label{contact:two-triple-collision-221-reduction-5}
\]

Thus $1/2<G<20$. It follows that $1\le e<5/2$ and $0<t<7/2$.
For $e\ge 5/2$, the expression in \eqref{contact:two-triple-collision-221-reduction-5} increases with $t\ge 0$ and is at
least $B_0\ge 175/8>20$. For $t\ge 7/2$, it is at least
$t^3-(3/2)t^2+1/2\ge 25>20$.

Use $e=1+(3/2)E$, $t=(7/2)T$, and $d=h-t$. Together with \eqref{contact:two-triple-collision-221-reduction-3}, the
two parameter cubes are $(k,w,E,T,m_1) \in [0,1]^5$.

\paragraph{Contact values determine the remaining masses}

The first triple's scaled level difference is

\[
 U=a^2(3/2-2a)-b^3+\frac32b^2+3a(1-a)b>0.
\tag{C.31.6}\label{contact:two-triple-collision-221-reduction-6}
\]

Actual deletion levels give

\[
 \boxed{m_2=U/G.}
\tag{C.31.7}\label{contact:two-triple-collision-221-reduction-7}
\]

Centering determines all remaining probabilities:

\[
 \ell_1=\frac{b-m_1h}{A_1},\quad
 r_1=\frac{a-m_1u}{A_1},\qquad
 \ell_2=\frac{e-m_2v}{A_2},\quad
 r_2=\frac{t-m_2h}{A_2}.
\tag{C.31.8}\label{contact:two-triple-collision-221-reduction-8}
\]

Require all six masses strictly positive. The second middle contact supplies
the exact equation

\[
 \boxed{(1-r_1)G-g_e(d)=0.}
\tag{C.31.9}\label{contact:two-triple-collision-221-reduction-9}
\]

In the open region $u>0$, equation \eqref{contact:two-triple-collision-221-reduction-9} can eliminate the last free mass:
$m_1=(A_1 g_e(d)/G-b)/u$. The five-cube formulation retains $m_1$ and
\eqref{contact:two-triple-collision-221-reduction-9}, avoiding a division whose denominator vanishes on the coalescing-support
face. This is a numerical choice, not an extra mathematical degree of freedom
in an actual solution.

\paragraph{Collision-preserving gradients and actual moments}

Define the nonnegative contact slopes

\[
 H_c=3[a(1-a)+c-c|c|],\quad H_b=3[a(1-a)+b-b^2],
\]
\[
 H_t=3[t^2-t+e(e-1)],\quad H_d=3[e(e-1)+d-d|d|].
\tag{C.31.10}\label{contact:two-triple-collision-221-reduction-10}
\]

Moving the two support points in either selected representation by opposite
amounts, with the usual recentering of the two marginals, preserves both
selected raw sums. These two admissible variations give

\[
 \boxed{m_1H_c=m_2H_d,\qquad r_1H_b=\ell_2H_t.}
\tag{C.31.11}\label{contact:two-triple-collision-221-reduction-11}
\]

Consequently

\[
 \sigma=m_1H_c+r_1H_b=\ell_2H_t+m_2H_d>0.
\tag{C.31.12}\label{contact:two-triple-collision-221-reduction-12}
\]

The scaled variances and third absolute moments are

\[
 V_1=ab-m_1uh,\quad V_2=te-m_2hv,\quad V=V_1+V_2,
\]
\[
 R_1=\ell_1a^3+m_1|c|^3+r_1b^3,
 \quad R_2=\ell_2t^3+m_2|d|^3+r_2e^3,
 \quad R_0=R_1+R_2.
\tag{C.31.13}\label{contact:two-triple-collision-221-reduction-13}
\]

The selected scaled sum is $o=b-t=c+d$. Thus

\[
 K=V^{-1/2},\quad z=o/\sqrt V,\quad C=V^{3/2}/G,
\]
\[
 \boxed{3(V-R_0)-\sigma o=0,\qquad
 \sigma\sqrt V-G\phi(o/\sqrt V)=0.}
\tag{C.31.14}\label{contact:two-triple-collision-221-reduction-14}
\]

The actual closed CDF is

\[
 F=q_2-m_2r_1.
\]

Its exact normal equation and a useful necessary positive expression are

\[
 \boxed{G[F-\Phi(o/\sqrt V)]-R_0=0,}
\tag{C.31.15}\label{contact:two-triple-collision-221-reduction-15}
\]
\[
 \boxed{B_0-\tfrac32V_2-R_1>0.}
\tag{C.31.16}\label{contact:two-triple-collision-221-reduction-16}
\]

Equation \eqref{contact:two-triple-collision-221-reduction-16} follows by averaging the second dual \eqref{contact:two-triple-collision-221-reduction-4}, giving
$GF=R_2-(3/2)V_2+B_0$, and using $\Phi(z)=F-R_0/G>0$.

The fixed-two-coordinate and dominant-coordinate results impose

\[
 \frac{25}{64}<V<4,\qquad V_i>\frac{13}{200}V\quad(i=1,2),
\]
\[
 (49/120)^2G^2\le V^3<(227/500)^2G^2.
\tag{C.31.17}\label{contact:two-triple-collision-221-reduction-17}
\]

The unconditional two-triple mixture-window result additionally requires

\[
 3(W_1+W_2)>V,
\tag{C.31.18}\label{contact:two-triple-collision-221-reduction-18}
\]

where the actual canonical variance widths are
$W_1=\min(bu,ah)$ and $W_2=\min(eh,tv)$. Equivalently,
$W_1=bu$ for $c\le 0$ and $W_1=ah$ for $c\ge 0$, while
$W_2=(eh+tv-|d|A_2)/2$. The adjacent-gap products $uh,hv$ only
give upper bounds for these widths; they are not the canonical widths.

\paragraph{The full actual variance Hessian, with cleared denominators}

For the centered unit-variance directions of the two triples, the actual CDF
first and mixed derivatives are

\[
 L_1=-\frac{m_2V}{hA_1},\qquad
 L_2=\frac{(r_1A_2-h)V}{hvA_2},\qquad
 g=\frac{V^2}{A_1h^2v}>0.
\tag{C.31.19}\label{contact:two-triple-collision-221-reduction-19}
\]

These follow directly from the deletion levels and signed weights in
Proposition~\ref{prop:finite-contact-exclusion}.
In particular the mixed coefficient is that of the actual $221$ event.

Use centered mass parameters whose individual variations have raw
total-variance derivative $V$; these are unit-variance directions
after standardization. Write
\[
 J=F-\Phi(o/\sqrt V)-\frac{CR}{V^{3/2}},
\]
where $C$ is fixed, the raw quantities $F,V,R$ vary, and the raw
threshold $o$ stays fixed.
At the attained maximum, both first derivatives vanish and the Hessian is
negative semidefinite. If $\gamma_i$ is the scaled second divided difference
of $|x|^3$, put

\[
 B_i=\frac{15R_0+G\phi'''(z)}4-3V\gamma_i=-GJ_{ii}.
\]

The full condition is

\[
 B_1\geq0,\quad B_2\geq0,\quad
 B_1B_2-\left[Gg-\frac{B_1+B_2}{2}\right]^2\geq0.
\tag{C.31.20}\label{contact:two-triple-collision-221-reduction-20}
\]

Eliminating both divided differences and the normal derivative gives polynomial necessary inequalities. The variance identity
in Section~\ref{contact:one-triple-bernoulli-collision-211}
is

\[
 4J_{ii}=12L_i+3D+(z^3+3z)\phi(z),
 \qquad D=\frac{R_0}{G}=\frac{CR_0}{V^{3/2}}.
\]

For completeness, the raw third-moment derivative in this direction is
$V\gamma_i$. The first derivative gives
$C\gamma_i/\sqrt V=L_i+3D/2+z\phi(z)/2$, while the second gives
$J_{ii}=3C\gamma_i/\sqrt V-(15D+\phi'''(z))/4$.
Eliminating $\gamma_i$ proves the displayed identity.

Using \eqref{contact:two-triple-collision-221-reduction-14}, define

\[
 M=3R_0V+\sigma o(o^2+3V),\qquad
 D_*=4Vh^2A_1vA_2>0,
\]
\[
 P_1=12V^2Uh vA_2-Mh^2A_1vA_2,
\]
\[
 P_2=-12GV^2(r_1A_2-h)hA_1-Mh^2A_1vA_2,
 \qquad Q=4GV^3A_2.
\tag{C.31.21}\label{contact:two-triple-collision-221-reduction-21}
\]

Direct substitution gives $P_i=D_*B_i$ and $Q=D_*Gg$. Therefore the
equivalent denominator-cleared necessary conditions are

\[
 \boxed{P_1\geq0,\qquad P_2\geq0,\qquad
 4P_1P_2-(2Q-P_1-P_2)^2\geq0.}
\tag{C.31.22}\label{contact:two-triple-collision-221-reduction-22}
\]

Only division by $G>1/2$ and the positive outer spans is needed to evaluate
the necessary system. The Hessian tests remain finite when one of
the adjacent gaps approaches zero.

\paragraph{Additional global-dual and threshold-gap tests}

At a maximizing threshold, every excluded sum atom at $z+\delta$ with
$\delta>0$ has mass at most $\Phi(z+\delta)-\Phi(z)\le \phi(0)\delta$.
Since $\phi(0)<2/5$, a scaled gap $\Delta=\delta/K$ and atom mass $w_0$
must satisfy

\[
 25w_0^2V<4\Delta^2.
\tag{C.31.23}\label{contact:two-triple-collision-221-reduction-23}
\]

Four available excluded atoms give $(w_0,\Delta)$ respectively

\[
 (\ell_1r_2,A_2-A_1),\quad (m_1r_2,v),
 \quad(r_1m_2,h),\quad(r_1r_2,A_2).
\]

For the first pair, the strict sum-gap \eqref{contact:two-triple-collision-221-reduction-2} cannot approach zero with a fixed
positive atom mass. This controls the transition from $221$ to $321$ at
positive masses, rather than treating it as a missing boundary face.

Two further tests follow by evaluating the global majorants at unused jumps.
For the first dual, the point $o-e=b-A_2<-a$ has deletion level one. The
negative cubic factors at $-a$, giving

\[
 \boxed{(A_2-A_1)^2(A_2-b+2a-3/2)-r_2G\geq0.}
\tag{C.31.24}\label{contact:two-triple-collision-221-reduction-24}
\]

For the second dual, $o+a=A_1-t$ lies strictly between $d$ and $e$,
and its deletion level is $\ell_1$. Hence

\[
 \boxed{g_e(A_1-t)-\ell_1G\geq0.}
\tag{C.31.25}\label{contact:two-triple-collision-221-reduction-25}
\]

These tests use only global majorization and the actual convolution levels.
\subsection{Two triples: eliminating the masses}\label{contact:two-triple-221-mass-elimination}
Keep the contact scale of Section~\ref{contact:two-triple-collision-221-reduction}.



Use the notation and equal-height charts of Section~\ref{contact:two-triple-collision-221-reduction}:

\[
 T_1=(-a,c,b),\quad T_2=(-t,d,e),\quad
 u=a+c,\ h=b-c=t+d,\ v=e-d,
\]
\[
 A_1=a+b=u+h,\quad A_2=t+e=h+v,\quad A_2>A_1.
\]

The supports depend on only $(k,w,E,T) \in [0,1]^4$ in each chart:
$a,c$ are the equal-height functions, $b=w B(a)$, $e=1+3E/2$,
$t=7T/2$, and $d=b-c-t$. The eliminated variable was the first middle mass.

\paragraph{Contact height determines the useful mass products}

Let $G=g_e(-t)$, $g=g_e(d)$, and $U$ be the first triple's positive level
drop, as in the reduction. Put

\[
 N=A_1g-bG,\qquad N_0=bG-hg,\qquad Z=N/G.
\tag{C.32.1}\label{contact:two-triple-221-mass-elimination-1}
\]

The second middle deletion level gives $q_1=1-r_1=g/G$. Centering the first
coordinate then yields

\[
 \boxed{r_1=1-g/G,\quad m_1u=Z,\quad
 m_1={N\over uG},\quad \ell_1={N_0\over uG}.}
\tag{C.32.2}\label{contact:two-triple-221-mass-elimination-2}
\]

Thus all three masses are positive exactly when $N>0$, $N_0>0$, $G-g>0$
(with $G,u>0$). They sum to one and have mean zero by direct substitution.
Use $Z=m_1u$; the formulas below then require no division by $u$.
The second coordinate is as follows:

\[
 m_2=U/G,\qquad \ell_2=(e-m_2v)/A_2,\qquad
 r_2=(t-m_2h)/A_2.
\tag{C.32.3}\label{contact:two-triple-221-mass-elimination-3}
\]

\paragraph{Equal height removes the apparent cubic-moment singularity}

The first low and middle contacts have equal cubic-majorant height.
Subtracting those values gives the exact identity

\[
 |c|^3-a^3=(a+c)\left[\frac32(c-a)+3a(1-a)\right].
\]

Define its bracket as $W_c$. Since $\ell_1+m_1=q_1$, the first variance and
sum of third absolute moments become

\[
 \boxed{V_1=ab-hZ,\qquad
 R_1=q_1a^3+ZW_c+r_1b^3.}
\tag{C.32.4}\label{contact:two-triple-221-mass-elimination-4}
\]

These are exactly the centered raw moments. Neither contains a small-gap
denominator. The same holds for the selected derivative. The quotient
$J_c=H_c/u$ has the explicit closed-chart extensions

\[
 \boxed{
 J_c=\begin{cases}
 u=3k/4,&\text{negative chart},\\[1mm]
 \displaystyle{3(1-k)(k^2+4k+1)\over2(k^3+3k+2)},
       &\text{positive chart}.
 \end{cases}}
\tag{C.32.5}\label{contact:two-triple-221-mass-elimination-5}
\]

For the negative chart, $H_c=u^2$. For the positive chart, substitute
$c=ka$ into $H_c/u=3[1-a(1+k^2)/(1+k)]$, and then the exact rational formula
for $a(k)$. Hence

\[
 \boxed{\sigma=ZJ_c+r_1H_b.}
\tag{C.32.6}\label{contact:two-triple-221-mass-elimination-6}
\]

The two retained representation-gradient equations can be written

\[
 \boxed{N J_c-UH_d=0,\qquad
 A_2(G-g)H_b-(eG-Uv)H_t=0.}
\tag{C.32.7}\label{contact:two-triple-221-mass-elimination-7}
\]

They are just the original equations multiplied by candidate-positive
factors. The second-middle height equation has been used exactly, not
discarded.

\paragraph{Remaining equations, probability signs and curvature}

Set $V_2=te-m_2hv$ and compute $R_2$ from the second raw support and masses.
Then $V=V_1+V_2$, $R=R_1+R_2$, $o=b-t$, $z=o/\sqrt{V}$. The actual event is

\[
 F=\ell_2+m_2q_1.
\]

The remaining three zero equations are as follows:

\[
 3(V-R)-\sigma o=0,\qquad
 \sigma^2V-G^2\phi(z)^2=0,\qquad
 G[F-\Phi(z)]-R=0.
\tag{C.32.8}\label{contact:two-triple-221-mass-elimination-8}
\]

The positive slope requirement retains the sign of the unsquared density
equation. All ratio, raw moment, variance-fraction, width, diameter and
majorant tests remain necessary. In particular the useful CDF expression remains
$e^2(2e-3/2)-(3/2)V_2-R_1>0$.

The full variance Hessian in
\eqref{contact:two-triple-collision-221-reduction-21}--\eqref{contact:two-triple-collision-221-reduction-22} uses
only $U$, $V$, $R$, $\sigma$, $r_1$, $h$, $v$, $A_1$, $A_2$, $G$ and $o$.
All are now functions of the four retained parameters with no division by
$u$, so the identical three cleared Hessian inequalities can be retained.
For an independent original-Hessian formulation, its first divided
difference also has the regular expression

\[
 \gamma_1={b^3-a^3-A_1W_c\over hA_1}.
\tag{C.32.9}\label{contact:two-triple-221-mass-elimination-9}
\]

This follows by subtracting the two secant slopes and using the equal-height
identity. It likewise has no $u$ denominator.

\paragraph{Uniform bounds for the denominators that remain}

The canonical widths satisfy

\[
 W_1=\min(bu,ah),\quad W_2=\min(eh,tv),\quad
 W_1+W_2>V/3>25/192.
\]

Since $a<1$ and $e<5/2$, their sum is at most $(a+e)h<(7/2)h$.
Since $v>u$, $b<B(a)<5/4$, and $t<7/2$, it is also smaller than
$(b+t)v<(19/4)v$. Therefore

\[
 \boxed{h>25/672,\qquad v>25/912.}
\tag{C.32.10}\label{contact:two-triple-221-mass-elimination-10}
\]

Also $V_1>13/512$ and $V_1\le ab<b$ imply $b>13/512$.
Similarly $t>13/1280$ follows from $V_2>13/512$ and $e<5/2$.
The outer spans obey $A_1>1/2$, $A_2\ge 1$, and $G>1/2$ was already proved.
Thus every remaining rational moment/curvature denominator is uniformly
bounded away from zero at a genuine candidate. The first gap $u$ may tend
to zero, but no evaluation denominator uses it.

\paragraph{Global-dual and atom-gap tests without individual first masses}

The first unused dual jump is the inequality already stated for $g_a$. For the second unused jump,
multiply its necessary inequality by $u>0$ and use \eqref{contact:two-triple-221-mass-elimination-2}:

\[
 \boxed{u\,g_e(A_1-t)-N_0\geq0.}
\tag{C.32.11}\label{contact:two-triple-221-mass-elimination-11}
\]

For the excluded low/high and middle/high atoms, multiply the squared
atom-gap inequalities by $(uG)^2>0$:

\[
 4(A_2-A_1)^2u^2G^2-25N_0^2r_2^2V>0,
\]
\[
 4v^2u^2G^2-25N^2r_2^2V>0.
\tag{C.32.12}\label{contact:two-triple-221-mass-elimination-12}
\]

The other two atom-gap inequalities need only $r_1,m_2,r_2$, already
available without $u$ division. Two aggregate tests remain nontrivial on
the coalescing face. At the middle/high sum, both first-low/second-high and
first-middle/second-high atoms have entered the CDF, with total mass
$q_1r_2$ and maximal raw gap $v$. At the largest sum, the entire excluded
mass $1-F$ has entered, with raw gap $A_2$. Therefore

\[
 \boxed{4v^2-25(q_1r_2)^2V>0,\qquad
 4A_2^2-25(1-F)^2V>0.}
\tag{C.32.13}\label{contact:two-triple-221-mass-elimination-13}
\]

Both follow from threshold maximality and $\phi(0)<2/5$, exactly as the
individual atom-gap tests. Coincidences among later atoms cause no problem.

Each genuine collided maximum is covered by the support chart
and maps into one of the two original closed four-cubes. Conversely, on
$u,G>0$, the retained probability signs reconstruct a normalized centered
first triple through \eqref{contact:two-triple-221-mass-elimination-2}. The closed face $u=0$ is included by continuity.
\subsection{Support geometry inequalities}\label{contact:candidate-support-geometry-closure}
The following identities are used in the contact-scale charts of Sections~\ref{contact:collision-210-contact-reduction} and~\ref{contact:two-triple-collision-221-reduction}.



If the candidate values satisfy an affine equation
$\sum_j a_j x_j=b$, with fixed nonzero exact coefficients, then for every
$i$ they satisfy

\[
 x_i={b-\sum_{j\ne i}a_jx_j\over a_i}. \tag{C.33.1}\label{contact:candidate-support-geometry-closure-1}
\]

Evaluate the right side on current candidate intervals and intersect with
the interval for $x_i$. This preserves every candidate. Any fixed
sequence of such operations is valid by induction; stopping before a
fixed point only weakens the resulting test. No derivative is involved.
The coefficients used below are integers or exact half-integers.

For collided $210$, use all support equations

\[
 A=t+b=h+v,\quad h=t+c=d+e,\quad v=b-c,
 \quad o=c-d=e-t,
\]

and the probability, variance, and third-moment sums. For collided $221$, use

\[
 A_1=a+b=u+h,\quad A_2=t+e=h+v,\quad
 u=a+c,\quad h=b-c=t+d,\quad v=e-d,
\]
\[
 \delta=A_2-A_1=v-u,\qquad o=b-t=c+d.
\]

On its chart with nonpositive first middle contact also use $3a-u=3/2$, $3c-2u=-3/2$,
and $4u=3k$, where $k$ is that chart's parameter (the contact scale remains
one). These are the exact chart identities, including the retained $u=0$
face. On the positive chart $a>0$, so $k=c/a$ may be intersected before
re-evaluating its rational chart formulas.

After a support interval narrows, re-evaluate the original cubic duals and
their slopes on those narrower candidate values. For example in $221$,
$G=g_e(-t)$, $g=g_e(d)$, and $m_2=U/G$ remain exact candidate equalities.
These are not ambient derivatives of the support-chart residuals.

Two existing unused-jump majorizations are also useful in inverse form:

\[
 r_2\le {\delta^2(A_2-b+2a-3/2)\over G},\qquad
 Y\le {u\,g_e(A_1-t)\over G}. \tag{C.33.2}\label{contact:candidate-support-geometry-closure-2}
\]

Here $Y=u \ell_1$; \eqref{contact:candidate-support-geometry-closure-2} never divides by $u$. The factor in the first bound
is nonnegative at every candidate since $a\ge 1/2$ and
$A_2-b+2a-3/2=\delta+3a-3/2$. On the negative chart it is exactly
$\delta+u=v$, so the first bound simplifies to

\[
 \boxed{\quad r_2\le\delta^2v/G.\quad} \tag{C.33.3}\label{contact:candidate-support-geometry-closure-3}
\]

The potential $g_e(A_1-t)$ in the second bound is nonnegative because the
global cubic majorant dominates a nonnegative deletion CDF. Its factorized
form is $\delta^2(3e-\delta-3/2)$ when $A_1-t\ge 0$, and its absolute-cubic form
works without a sign restriction. Divisions by $G$, outer spans and $a$
use their already proved positive lower bounds.
\subsection{The support Hessian with fixed moments}\label{contact:two-triple-221-fixed-moment-support-hessian}
Keep the contact scale of Section~\ref{contact:two-triple-collision-221-reduction}.



Keep the scaled supports $(-a,c,b)$, $(-t,d,e)$, gaps
$u=a+c$, $h=b-c=t+d$, $v=e-d$, spans $A_1=u+h$, $A_2=h+v$, and
$A_2>A_1$. Let $G$, the six masses, and the contact slopes have their
previous meanings. Write $Z=u m_1$, $J_c=H_c/u$, and $L(x)=6|x|-3$.
The selected-support equations are

\[
 m_1H_c=m_2H_d,\qquad r_1H_b=\ell_2H_t. \tag{C.34.1}\label{contact:two-triple-221-fixed-moment-support-hessian-1}
\]

\paragraph{An actual two-parameter family}

Vary the supports as

\[
 (-a,c+\epsilon,b+\eta),\qquad (-t-\eta,d-\epsilon,e).
\tag{C.34.2}\label{contact:two-triple-221-fixed-moment-support-hessian-2}
\]

For each marginal separately choose its three masses to keep its total
mass, mean zero, and original variance fixed. The two $3\times3$ Vandermonde
systems have unique smooth solutions near every genuine candidate. The
base masses are positive, so both signs of small $\epsilon$ and $\eta$ are
admissible. Both selected raw sums remain exactly $o=c+d=b-t$ and all
other event inequalities remain strict. Thus the event matrix is the
actual $221$ matrix throughout this family. Its CDF is
$F=\ell_2+m_2(1-r_1)$.

Both the centered threshold and the total variance are fixed. The normal
CDF term is therefore constant, and the relevant objective is $F-R/G$.
Let $P_1,P_2$ be the scaled cubic dual potentials. Their values at the three
supports equal $G$ times the actual deletion probabilities. Their first
derivatives are respectively $(0,-H_c,-H_b)$ and $(-H_t,-H_d,0)$, and their
second derivatives are $L(x)$. Adding the fixed mean, variance and mass
constraints to the objective changes it only by a constant on \eqref{contact:two-triple-221-fixed-moment-support-hessian-2}.
The resulting Lagrangian has zero derivatives in all mass variables at
the base point.

Consequently, for two tangent support motions $s,t$, its Hessian multiplied
by $G$ is exactly

\[
 G\sum_{ij} I_{ij}(w_{1i}^{s}w_{2j}^{t}+w_{1i}^{t}w_{2j}^{s})
 -\sum_{ij}p_{ij}L(x_{ij})\dot x_{ij}^{s}\dot x_{ij}^{t}
 -\sum_{ij}P_i'(x_{ij})(w_{ij}^{s}\dot x_{ij}^{t}
                         +w_{ij}^{t}\dot x_{ij}^{s}). \tag{C.34.3}\label{contact:two-triple-221-fixed-moment-support-hessian-3}
\]

There are no second derivatives of the masses in \eqref{contact:two-triple-221-fixed-moment-support-hessian-3}: their coefficients
are the zero first mass derivatives of this Lagrangian. The support paths
are affine. Equation \eqref{contact:two-triple-221-fixed-moment-support-hessian-1} makes both first support derivatives zero.

\paragraph{Exact weights and regular Hessian}

For a marginal with support points $x_0,x_1,x_2$, let $L_i$ be its
quadratic Lagrange basis, defined by $L_i(x_j)=\ind_{\{i=j\}}$.
If $x_j$ moves with velocity $s$, preserving mass, mean, and variance
requires the probability derivative $w_i=-p_j s L_i'(x_j)$.
The entries needed in the actual event are

\[
\begin{array}{c|cc}
 &\epsilon&\eta\\ \hline
 w_{1r}&-Z/(A_1h)&-r_1(A_1+h)/(A_1h)\\
 w_{1m}&-m_1(h-u)/(uh)&r_1A_1/(uh)\\
 w_{2m}&m_2(v-h)/(hv)&\ell_2A_2/(hv)\\
 w_{2\ell}&-m_2v/(hA_2)&-\ell_2(h+A_2)/(hA_2).
\end{array} \tag{C.34.4}\label{contact:two-triple-221-fixed-moment-support-hessian-4}
\]

The mixed event form is exactly $-w_{1r}w_{2m}$. Substitution in \eqref{contact:two-triple-221-fixed-moment-support-hessian-3}, followed
only by \eqref{contact:two-triple-221-fixed-moment-support-hessian-1}, gives the scaled Hessian

\[
 \begin{pmatrix}-C_m&B\\ B&-C_e\end{pmatrix}\preceq0, \tag{C.34.5}\label{contact:two-triple-221-fixed-moment-support-hessian-5}
\]

where

\[
 C_m=m_2L(d)+Z\left({L(c)+2J_c\over u}-{2J_c\over v}\right)
       -{2GZm_2(v-h)\over A_1h^2v}, \tag{C.34.6}\label{contact:two-triple-221-fixed-moment-support-hessian-6}
\]
\[
 C_e=r_1L(b)+\ell_2L(-t)
       +2r_1H_b\left({1\over A_1}-{1\over A_2}\right)
       -{2Gr_1\ell_2A_2(A_1+h)\over A_1h^2v}, \tag{C.34.7}\label{contact:two-triple-221-fixed-moment-support-hessian-7}
\]
\[
 \begin{split}
 B={G\{Z\ell_2A_2+r_1m_2(A_1+h)(v-h)\}\over A_1h^2v}
   -{ZH_b\over A_1h}+{r_1A_1J_c\over h}
   +{m_2vH_t\over hA_2}-{\ell_2A_2H_d\over hv}.
 \end{split} \tag{C.34.8}\label{contact:two-triple-221-fixed-moment-support-hessian-8}
\]

Thus every genuine candidate satisfies $C_m\ge 0$, $C_e\ge 0$, and
$C_m\,C_e-B^2\ge 0$. On the negative first-middle chart, $J_c=u$, $L(c)=-4u$,
so \eqref{contact:two-triple-221-fixed-moment-support-hessian-6} has the regular form

\[
 \boxed{C_m=m_2L(d)-2Z(1+u/v)
                 -{2GZm_2(v-h)\over A_1h^2v}.} \tag{C.34.9}\label{contact:two-triple-221-fixed-moment-support-hessian-9}
\]

\paragraph{Analytic restriction when $c\le 0$ and $v\ge h$}

In this genuine subcase $Z>0$, $m_2>0$, and $0<u\le 3/4$. Equation \eqref{contact:two-triple-221-fixed-moment-support-hessian-9}
first gives $L(d)>0$, hence $|d|>1/2$. If $d>1/2$, then $h=t+d>d$
and $e=d+v\ge d+h>2d\ge 1$. Since $x(x-1)$ increases for $x\ge 1$,

\[
 H_d=3\{e(e-1)+d-d^2\}>3d(3d-1).
\]

But \eqref{contact:two-triple-221-fixed-moment-support-hessian-1} and \eqref{contact:two-triple-221-fixed-moment-support-hessian-9} also give $L(d)\ge 2Z/m_2=2Hd/u$, so

\[
 6d-3>8d(3d-1).
\]

This is impossible because

\[
 24d^2-14d+3=24(d-7/24)^2+23/24>0.
\]

Therefore

\[
 \boxed{c\le0,\quad v\ge h\quad\Longrightarrow\quad d<-1/2.} \tag{C.34.10}\label{contact:two-triple-221-fixed-moment-support-hessian-10}
\]

There is also a useful upper-support bound. Write $w=-d>1/2$.
Equation \eqref{contact:two-triple-221-fixed-moment-support-hessian-9} is strict in the weaker form $m_2 L(d)>2Z$, since $u/v>0$.
Using \eqref{contact:two-triple-221-fixed-moment-support-hessian-1} yields

\[
 e(e-1)+w^2-w<u(w-1/2).
\]

Complete the square in $w$ to obtain

\[
 \boxed{e(e-1)<{1+u^2\over4},\qquad
 e<{1+\sqrt{2+u^2}\over2}\le {1\over2}+{\sqrt{41}\over8}.} \tag{C.34.11}\label{contact:two-triple-221-fixed-moment-support-hessian-11}
\]

The last two inequalities hold when $v\ge h$; the next calculation proves this condition on the negative chart.
\subsection{The negative middle contact}\label{contact:two-triple-221-negative-middle-location}
Keep the contact scale of Section~\ref{contact:two-triple-collision-221-reduction}.




On the chart with nonpositive first middle contact,

\[
 c=2a-3/2,\quad u=a+c\in(0,3/4],\quad
 h=b-c\le7/4,\quad v=e-d>u,\quad e\ge1.
\]

The actual first dual height and selected-flow identity give

\[
 U=A_1^2h-2b^3<A_1^2h,\qquad
 (A_1g-bG)u=U H_d,\quad g=g_e(d),\quad G>0.
\tag{C.35.1}\label{contact:two-triple-221-negative-middle-location-1}
\]

In particular $H_d>0$ and $U/A_1<u g/H_d$. The middle support
curvature, after using $Z u=m_2 H_d$, is

\[
 {C_m\over m_2}=L(d)-{2H_d\over u}Q,\quad
 Q=1+{u\over v}-{U(h-v)\over A_1h^2v},\quad L(d)=6|d|-3.
\tag{C.35.2}\label{contact:two-triple-221-negative-middle-location-2}
\]

A maximum requires $C_m\ge 0$.

First suppose $-1/2\le d\le 1/8$. If $v\ge h$, then $Q\ge 1+u/v>0$.
If $v<h$, use the first upper bound in \eqref{contact:two-triple-221-negative-middle-location-1}:

\[
 Q>2-h/v+u/h.
\]

Since $v=e-d\ge 7/8$ and $h\le 7/4$, the right side is at least $u/h>0$.
But $L(d)\le 0$, so \eqref{contact:two-triple-221-negative-middle-location-2} is strictly negative. This interval is impossible.

Now suppose $d\ge 1/8$. If $v\ge h$, \eqref{contact:two-triple-221-negative-middle-location-2} and $v>u$ give

\[
 C_m/m_2<L(d)-2H_d/v=3-6e<0.
\]

If $v<h$, write $r=v/h$ and $\ell=e+d-1\ge d\ge 1/8$.
The positive-side cubic potential factors exactly as

\[
 {g\over H_d}={v\over2}+{v^2\over6\ell}.
\]

Hence, using $0<r<1$, $h\le 7/4$,

\[
 {g(h-v)\over H_dh^2}
 ={r(1-r)\over2}+{h\over6\ell}r^2(1-r)
 \le {1\over8}+{7\over3}{4\over27}
 ={305\over648}<1.
\tag{C.35.3}\label{contact:two-triple-221-negative-middle-location-3}
\]

The second upper bound in \eqref{contact:two-triple-221-negative-middle-location-1} now implies

\[
 Q>1+{u\over v}-{u\over v}{305\over648}>1>{u\over v}.
\]

Again \eqref{contact:two-triple-221-negative-middle-location-2} is strictly below $L(d)-2Hd/v=3-6e<0$.
This excludes every $d\ge 1/8$. Combining the two ranges proves

\[
 \boxed{d<-1/2\quad\hbox{throughout the chart with nonpositive first middle contact}.}
\tag{C.35.4}\label{contact:two-triple-221-negative-middle-location-4}
\]

In fact the opposite-gap sector is impossible. Put $w=-d>1/2$.
The first high support satisfies the exact bound

\[
 b<B(a)={1+\sqrt{1+4a(1-a)}\over2}
 \le {1+\sqrt2\over2}.
\]

Consequently

\[
 h=b-c<1+1/\sqrt2-2u/3<1+1/\sqrt2.
\tag{C.35.5}\label{contact:two-triple-221-negative-middle-location-5}
\]

If $v\le h$, then $v=e+w$ lies in $(3/2,1+1/\sqrt{2})$.
For fixed $v$, convexity in $w$ on $[1/2,v-1]$ gives

\[
 {H_d\over3}=e^2-e+w^2-w
 \le\max\{v^2-2v+1/2,\ (v-1)(v-2)\}<0.
\]

Both endpoint expressions are strictly negative on the stated interval.
This contradicts $H_d>0$ from the selected flow. Thus

\[
 \boxed{v>h\quad\hbox{throughout the chart with nonpositive first middle contact}.}
\tag{C.35.6}\label{contact:two-triple-221-negative-middle-location-6}
\]

The conditional negative-chart support bound therefore
also holds unconditionally on this chart:
$e(e-1)<(1+u^2)/4$, hence $e<1/2+\sqrt{41}/8$.
\subsection{The signed divided-difference inequality}\label{contact:two-triple-221-corrected-dd-load}
In the two-triple $221$ contact scale let $r=\sqrt V$. The identities
$C=r^3/G$, $\phi(z)=\sigma r/G$, $z=o/r$, and $3(V-R)=\sigma o$ give
\[
 \frac Cr=D+\frac{z\phi(z)}3,\qquad
 L=\frac14+\frac{\sigma o(o^2+2V)}{12V^2}
  =\frac14+\frac{rz\phi(z)(z^2+2)}{12C}
  =\frac14+\frac{z\phi(z)(z^2+2)}{12(D+z\phi(z)/3)}.
\]
The own-variance curvature inequality is
\[
 \gamma_i\le\frac{r\{15D+(3z-z^3)\phi(z)\}}{12C}=\frac32-L.
\]
For $z\ge0$ one has $L\ge1/4$. For $z=-y\in[-5/4,0]$, the function
$y(y^2+2)\phi(y)$ increases: its derivative is
$(2-y^2)(1+y^2)\phi(y)>0$. Since
$\phi(0)<2/5$ and $\sum_{j=0}^5(25/32)^j/j!>24/11$,
$\phi(5/4)<11/60$. Using $r<2$ and $C\ge49/120$ gives
\[
 y(y^2+2)\phi(y)<\frac{285}{64}\frac{11}{60}=\frac{209}{256},
 \qquad L>\frac14-\frac{20}{49}\frac{209}{256}
 =-\frac{261}{3136}>-\frac1{12}.
\]
The height identities, with $q$ allowed either sign, are
\[
 p=\frac32-\gamma_1=\frac{U}{A_1h}>0,\qquad
 q=\frac32-\gamma_2=\frac{G(h-A_2r_1)}{A_2hv},\qquad p,q\ge L.
\]
Thus $\gamma_2<4965/3136$ and $A_2<4965/1568$ by
\eqref{eq:contact-dd-diameter}. The exact mass and event identities read
\[
 m_2=\frac{pA_1h}{G},\quad
 r_1=\frac h{A_2}-\frac{qhv}{G},\quad
 q_1=\frac v{A_2}+\frac{qhv}{G},\quad
 F=\frac e{A_2}+\frac{m_2qhv}{G}.
\]
Put $x=p-L$, $y=q-L$, and $b_0=GV/(3A_1h^2v)$. The original
probability Hessian is negative semidefinite only if
\[
 x\ge0,\quad y\ge0,\quad x+y\ge b_0,\qquad
 4xy-(2b_0-x-y)^2\ge0.
\]
These are the signed inequalities used in the supplied positive-chart
calculation. They impose no positivity assumption on $q$ or $F-e/A_2$.

\subsection{A parameter bound for a positive middle contact}\label{contact:positive-parameter}
Keep the contact scale of Section~\ref{contact:two-triple-collision-221-reduction}.

In the positive chart for the collided two-triple $221$ event, $k<8/9$.
Indeed, with $u=a(1+k)$ and $A_1=u+h$, the centered marginal and selected
contact equations give
\[
 \beta_1=1-\frac{2a^2}{A_1u},\quad
 \sigma=\frac{aH_b}{A_1}+\left(J_c-\frac{H_b}{A_1}\right)Z,\quad
 J_c-\frac{H_b}{A_1}=3h\beta_1,\quad Z=a-A_1r_1.
\]
Thus $0<Z<a$. If $k\ge8/9$, then $A_1u>u^2>2a^2$, so
$\beta_1>0$ and $\sigma\le aJ_c\le J_c$. But
\[
 J_c=\frac{3(1-k)(k^2+4k+1)}{2(k^3+3k+2)}
\]
decreases on $[1/2,1]$: its derivative has the sign of
$k^4-4k^3-6k^2-4k+1<0$, as follows from $k^4\le k^3$ and evaluation
of the resulting decreasing polynomial at $1/2$. Consequently
$\sigma\le J_c(8/9)=1299/7828$.
On the other hand $\sigma=V\phi(z)/C$, and the bounds
$V>2401/5476$, $\phi(z)>9/50$, $C<227/500$ imply
\[
 \sigma>\frac{108045}{621526}
 =\frac{1299}{7828}+\frac{19206993}{2432652764}.
\]
This contradiction proves the parameter bound used in the positive-chart
certificate.

The fixed certificates \texttt{tt221\_negative} and
\texttt{tt221\_positive} prove that neither compact support chart
contains a point satisfying the necessary conditions in
Sections~\ref{contact:two-triple-collision-221-reduction}--\ref{contact:positive-parameter}.
Thus the collided $221$ event and its transpose $320$ cannot occur
at a maximizing pair. Section~\ref{sec:verification} gives the enclosure
and partition rules.
